\chapter{Elliptic curves over local fields: reduction, Kodaira types, Tate curves}\label{ch:reduction}
\chaptermark{Elliptic curves over local fields}

The paper works one bad prime at a time. At each bad prime $q$ of the optimal curve $E$ it reads off a handful of local
invariants: the Kodaira symbol, the geometric component group $\Phig_q$ and the Tamagawa number $c_q$, the valuations
$v_q(\Delta)$ and $v_q(j)$, the conductor exponent, the Tate parameter $q_E$ and its unit part $u$ at a potentially
multiplicative prime, the order of the image of inertia at a potentially good prime, and the effect on all of these of the
quadratic twist by $\Q(\sqrt{q^*})$. This chapter defines each of them, proves what can be proved with the tools of the
book, states the rest precisely with references, and computes everything on the paper's curves. The local field theory
(inertia, the tame character, Kummer theory, Artin conductors) is in \Cref{ch:local}; Weierstrass equations, the group law,
torsion, Tate modules and quadratic twists of models are in \Cref{ch:ellcurves}; the Galois representations built from the
present results are in \Cref{ch:galrep}.

\section*{Core results at a glance}

Throughout, $K$ is complete for a discrete valuation $v$ with $v(\varpi)=1$, valuation ring $\mathcal O$, perfect residue field
$k$ of characteristic $p$, and $\ell\ne p$ unless stated otherwise.

\begin{enumerate}[label=\textbf{\thechapter.\arabic*}]
\item \textbf{Minimal models.} A model with all $a_i\in\mathcal O$ is integral; it is minimal if $v(\Delta)$ is least among integral
models. Minimal models exist; two of them differ by $(u,r,s,t)$ with $u\in\mathcal O^\times$, $r,s,t\in\mathcal O$; the N\'eron differential
is defined up to $\mathcal O^\times$, and $v(\Delta_{\min})$, $v(c_4)$, $v(c_6)$ of a minimal model and $v(j)=3v(c_4)-v(\Delta)$ are invariants
of $E/K$. \emph{(proved here)} \Cref{def:reduction:minimal}, \Cref{prop:reduction:minimal-basic}.
\item \textbf{Minimality criteria.} An integral model with $v(\Delta)<12$, or $v(c_4)<4$, or $v(c_6)<6$ is minimal. For $p\ge5$ it is
minimal if and only if $v(\Delta)<12$ or $v(c_4)<4$, and otherwise $(\varpi,0,0,0)$ applied to $y^2=x^3-27c_4x-54c_6$ lowers $v(\Delta)$ by
$12$. \emph{(proved here)} Over $\Z$, $c_4,c_6$ with $\Delta=(c_4^3-c_6^2)/1728\in\Z\setminus\{0\}$ come from an integral model if and only if
$v_3(c_6)\ne2$ and either $c_6\equiv-1\pmod4$ or $v_2(c_4)\ge4$, $c_6\equiv0,8\pmod{32}$ (Kraus). \emph{(stated; reference)}
\Cref{prop:reduction:min-criterion}.
\item \textbf{Global minimal models.} Every $E/\Q$ has a model over $\Z$ minimal at every $p$, unique up to $(\pm1,r,s,t)$,
$r,s,t\in\Z$, and unique with $a_1,a_3\in\{0,1\}$, $a_2\in\{-1,0,1\}$; its discriminant is $\Delta_E$. \emph{(proof sketched here)}
\Cref{thm:reduction:global-minimal}.
\item \textbf{Reduction types.} For a minimal model: good iff $v(\Delta)=0$; multiplicative iff $v(\Delta)>0=v(c_4)$ (node; split iff
the tangent slopes are in $k$); additive iff $v(\Delta)>0$, $v(c_4)>0$ (cusp). \emph{(proved here)} \Cref{def:reduction:types}.
\item \textbf{Nonsingular points.} $\tilde E_{\mathrm{ns}}(k)\simeq k^+$ (cusp), $k^\times$ (split node), $\ker(\Nm\colon l^\times\to k^\times)$
(nonsplit node, $l/k$ quadratic); over $\F_p$ of orders $p$, $p-1$, $p+1$. Hence for $E/\Q$ with bad reduction at $p$,
$a_p=p+1-\#\tilde E(\F_p)$ is $1$, $-1$, $0$ for split, nonsplit, additive. \emph{(proved here)} \Cref{prop:reduction:Ens},
\Cref{cor:reduction:ap-bad}.
\item \textbf{Split versus nonsplit.} After moving the node to $(0,0)$, split iff $T^2+a_1T-a_2$ has a root in $k$; for $p$ odd, split iff
$-c_6$ (equivalently $-c_4/c_6$) is a square in $k^\times$; for $p$ odd and $k$ finite the unramified quadratic twist exchanges split and
nonsplit, and over the unramified quadratic extension the reduction is split. \emph{(proved here)} \Cref{prop:reduction:split}.
\item \textbf{The reduction filtration.} With $E_0(K)=\{P:\tilde P\text{ nonsingular}\}$ and $E_1(K)=\{P:\tilde P=\tilde O\}$, there is an exact
sequence $0\to E_1(K)\to E_0(K)\to\tilde E_{\mathrm{ns}}(k)\to0$. \emph{(proved here)} \Cref{def:reduction:filtration},
\Cref{prop:reduction:filtration}.
\item \textbf{Formal group.} $E_1(K)\simeq\hat E(\mathfrak m)$ \emph{(stated; reference)}; for $m$ prime to $p$, $[m]$ is bijective on $E_1(K)$,
which has no prime-to-$p$ torsion \emph{(proved here, given the formal group)}. \Cref{thm:reduction:formal-group}.
\item \textbf{Torsion and reduction.} For $m$ prime to $p$, $E_0(K)[m]\into\tilde E_{\mathrm{ns}}(k)$. If $E$ has good reduction, $I_K$ acts
trivially on $E[m]$ and on $T_\ell E$, and reduction $E[m]\to\tilde E[m](\bar k)$ is a Galois-equivariant isomorphism. \emph{(proved here)}
\Cref{cor:reduction:torsion-injects}.
\item \textbf{Index of $E_0$.} For $k$ finite, $E(K)/E_0(K)$ is finite, trivial for good reduction. \emph{(proved here)}
\Cref{prop:reduction:index-finite}.
\item \textbf{Potential reduction.} $E$ has potentially good reduction iff $v(j)\ge0$, and potentially multiplicative reduction iff $v(j)<0$;
then it is multiplicative over every extension where it is semistable. Multiplicative reduction has $v(j)=-v(\Delta_{\min})$.
\emph{(proof sketched here)} \Cref{def:reduction:potential}, \Cref{thm:reduction:potential-criterion}.
\item \textbf{Semistable reduction.} $E$ becomes semistable over a finite extension, for instance $K(E[m])$ with $m\ge3$ prime to $p$.
\emph{(proof sketched here)} \Cref{thm:reduction:semistable}.
\item \textbf{The four classes.} The Kodaira symbols fall into: good reduction, $\kod{I}_0$; multiplicative, $\kod{I}_n$ ($n\ge1$,
$n=v(\Delta_{\min})=-v(j)$); additive potentially multiplicative, $\kod{I}_n^*$ ($n\ge1$); additive potentially good, $\kod{II}$, $\kod{III}$,
$\kod{IV}$, $\kod{I}_0^*$, $\kod{IV}^*$, $\kod{III}^*$, $\kod{II}^*$, and at $p=2$ only also some $\kod{I}_m^*$ ($m\ge1$). \emph{(proved here, given the
classification)} \Cref{rem:reduction:types-of-the-four-cases}.
\item \textbf{N\'eron models.} The N\'eron model $\mathcal E$ (smooth separated group scheme of finite type over $\mathcal O$, generic fibre $E$,
N\'eron mapping property) exists and is unique; it is the smooth locus of the minimal regular model. \emph{(stated; reference)}
\Cref{def:reduction:neron}, \Cref{thm:reduction:neron-existence}.
\item \textbf{Component group.} $\mathcal E(\mathcal O)=E(K)$, $\mathcal E_k^0\simeq\tilde E_{\mathrm{ns}}$, $\mathcal E^0(\mathcal O)=E_0(K)$, and
the quotient $\mathcal E_k/\mathcal E_k^0$ is a finite \'etale group scheme over $k$, the component group scheme, whose groups of points
are written $\Phi_E(k)$ and $\Phi_E(\bar k)=\Phig$ \emph{(stated; reference)}; for $k$ finite $E(K)/E_0(K)\simeq\Phi_E(k)$
\emph{(proof sketched here)}. \Cref{prop:reduction:component-group}.
\item \textbf{$\Phig$ and $c$.} For $E/\Q$: $\Phig_p=\Phi_p(\bar\F_p)$ and $c_p=\#\Phi_p(\F_p)=[E(\Q_p):E_0(\Q_p)]$, the Frobenius invariants;
$c_p\mid\#\Phig_p$, with equality iff Frobenius acts trivially. The paper's statements are about $\#\Phig_q$. \emph{(proved here)}
\Cref{def:reduction:tamagawa}, \Cref{conv:reduction:phig}.
\item \textbf{Unramified base change.} N\'eron models commute with \'etale base change; type, Kodaira symbol, $\Phig$, $v(\Delta_{\min})$ and $f$
are invariant under unramified extension; $c$ is not. \emph{(stated; reference)} \Cref{prop:reduction:etale-base-change}.
\item \textbf{Tamagawa numbers of $\kod{I}_n$.} $\Phig\simeq\Z/n$; $c=n$ (split), $c=\gcd(2,n)$ (nonsplit); for $\ell$ odd the $\ell$-part of $c$
is $\ell^{v_\ell(n)}$ (split) or $1$ (nonsplit). \emph{(proved here, given Tate's theorem)} \Cref{rem:reduction:tamagawa-mult},
\Cref{prop:reduction:tate-tamagawa}.
\item \textbf{Inertia invariants and component groups.} Let $K/\Q_q$ be finite, $\ell\ne q$. If the reduction is additive,
$E[\ell^t]^{I_K}=E(K^{\ur})[\ell^t]\simeq\Phig[\ell^t]$, Galois-equivariantly, and $E(K)[\ell^\infty]\simeq\Phi_E(k)[\ell^\infty]$. If it is
multiplicative, $0\to\mu_{\ell^t}\otimes\varepsilon\to E[\ell^t]^{I_K}\to\Phig[\ell^t]\to0$ with $\Phig\simeq\Z/n$, $\varepsilon$ trivial iff split.
For $\ell$ odd, $\dim E[\ell]^{I_K}$ is $1+[\ell\mid n]$ for $\kod{I}_n$, $1$ for $\kod{IV},\kod{IV}^*$ if $\ell=3$, and $0$ for every other
additive type. \emph{(proved here)} \Cref{prop:reduction:component-invariants}.
\item \textbf{Kodaira-N\'eron classification.} The special fibres and $(m,\Phig)$ are: $\kod{I}_0$, $\kod{I}_1$, $\kod{II}$: $m=1$, $0$;
$\kod{I}_n$: $n$, $\Z/n$; $\kod{III}$: $2$, $\Z/2$; $\kod{IV}$: $3$, $\Z/3$; $\kod{I}_0^*$: $5$, $(\Z/2)^2$; $\kod{I}_n^*$: $n+5$, $\Z/4$ ($n$ odd) or
$(\Z/2)^2$ ($n$ even); $\kod{IV}^*$: $7$, $\Z/3$; $\kod{III}^*$: $8$, $\Z/2$; $\kod{II}^*$: $9$, $0$. Possible $c$: $\kod{IV},\kod{IV}^*$: $1,3$;
$\kod{I}_0^*$: $1,2,4$; $\kod{I}_n^*$: $2,4$; $\kod{III},\kod{III}^*$: $2$; $\kod{II},\kod{II}^*$: $1$. \emph{(stated; reference)}
\Cref{thm:reduction:kodaira-neron}.
\item \textbf{Tate's algorithm.} The eleven steps of \Cref{thm:reduction:tate-algorithm} compute the type, $m$, $f=v(\Delta)+1-m$, $c$ and a
minimal model; explicit translations for Step 2 at $p\ge5$, $p=3$, $p=2$, and completing the square for Step 6 at odd $p$. \emph{(proof
sketched here)} \Cref{thm:reduction:tate-algorithm}.
\item \textbf{The tame table.} For $p\ge5$ and a minimal model, $(v(c_4),v(c_6),v(\Delta))$ is $(0,0,n)$ for $\kod{I}_n$; $(\ge1,1,2)$ $\kod{II}$;
$(1,\ge2,3)$ $\kod{III}$; $(\ge2,2,4)$ $\kod{IV}$; $(\ge2,\ge3,6)$ $\kod{I}_0^*$; $(2,3,n+6)$ $\kod{I}_n^*$; $(\ge3,4,8)$ $\kod{IV}^*$; $(3,\ge5,9)$ $\kod{III}^*$;
$(\ge4,5,10)$ $\kod{II}^*$. So for potentially good reduction the type is determined by $v(\Delta)\in\{0,2,3,4,6,8,9,10\}$, and $f=2$ for every
additive type. \emph{(proved here, given Tate's algorithm, and Ogg's formula for the index of $\kod{I}_n^*$)} \Cref{cor:reduction:tame-table}.
\item \textbf{The cubic criterion for $\kod{I}_0^*$.} For $p$ odd and $y^2=x^3+a_2x^2+a_4x+a_6$ with $\varpi\mid a_2$, $\varpi^2\mid a_4$, $\varpi^3\mid a_6$:
type $\kod{I}_0^*$ iff $T^3+a_{2,1}T^2+a_{4,2}T+a_{6,3}$ has distinct roots mod $\varpi$, and then $c=1+\#$roots in $k$; every $\kod{I}_0^*$ at odd
$p$ has such a model. The twist by $q^*$ of a model with good reduction at odd $q$ has cubic $T^3+\epsilon a_2T^2+a_4T+\epsilon a_6$,
$\epsilon=q^*/q$, hence type $\kod{I}_0^*$. \emph{(proved here)} \Cref{prop:reduction:I0star-cubic}.
\item \textbf{PARI's Kodaira codes.} $1$: $\kod{I}_0$; $n+4$: $\kod{I}_n$ and $-(n+4)$: $\kod{I}_n^*$ \emph{for $n\ge1$ only}; $-1$: $\kod{I}_0^*$;
$\pm2,\pm3,\pm4$: $\kod{II}/\kod{II}^*$, $\kod{III}/\kod{III}^*$, $\kod{IV}/\kod{IV}^*$. \texttt{elllocalred} returns $[f,k,[u,r,s,t],c]$.
A session on \curve{726a1}, \curve{20a1}, \curve{539c1}, \curve{11a1} shows the output. \emph{(proved here, as a dictionary)}
\Cref{conv:reduction:pari-codes}, \Cref{ex:reduction:pari-session}.
\item \textbf{Conductor exponent.} $f=a(V_\ell E)=\dim V-\dim V^{I}+\operatorname{sw}(V)$, independent of $\ell$ \emph{(stated; reference)}; it is $0$
(good), $1$ (multiplicative, $V^I$ a line), $2+\operatorname{sw}\ge2$ (additive, $V^I=0$), and $\operatorname{sw}=0$ for $p\ge5$, and for $p$ odd at an
additive potentially multiplicative prime \emph{(proved here, given the criterion of N\'eron, Ogg and Shafarevich)}.
\Cref{def:reduction:conductor-exponent}, \Cref{prop:reduction:conductor-values}.
\item \textbf{Ogg's formula.} $f=v(\Delta_{\min})+1-m$. \emph{(stated; reference)} \Cref{thm:reduction:ogg}.
\item \textbf{Conductor bounds.} $f_p\le2$ for $p\ge5$ \emph{(proved here)}, $f_3\le5$, $f_2\le8$ \emph{(stated; reference)}.
\Cref{thm:reduction:conductor-bounds}.
\item \textbf{The conductor.} $N=\prod_pp^{f_p}$; $p\,\|\,N$ iff multiplicative, $p^2\mid N$ iff additive, $p^3\nmid N$ for $p\ge5$; isogenous curves
have the same conductor. \emph{(proved here)} \Cref{def:reduction:conductor}, \Cref{prop:reduction:conductor-properties}.
\item \textbf{Serre-Tate.} For potentially good reduction there is a smallest $L\supseteq K^{\ur}$ with good reduction, $L=K^{\ur}(E[m])$ ($m\ge3$,
$p\nmid m$); $G_{\mathrm{ss}}=\Gal(L/K^{\ur})$ is the image of $I_K$ in $\Aut(T_\ell E)$ for every $\ell\ne p$; it embeds in $\SL_2(\Z/m)$, its elements
have orders $1,2,3,4,6$, and for $p\ge5$ it is cyclic. \emph{(proof sketched here)} \Cref{thm:reduction:serre-tate}.
\item \textbf{Kraus.} For $p\ge5$ and potentially good reduction, $G_{\mathrm{ss}}$ is cyclic of order $12/\gcd(v(\Delta_{\min}),12)$: $1,6,4,3,2,3,4,6$
for $\kod{I}_0,\kod{II},\kod{III},\kod{IV},\kod{I}_0^*,\kod{IV}^*,\kod{III}^*,\kod{II}^*$ \emph{(proved here, given Serre-Tate)}. At $p=3$:
cyclic of order $1,2,3,4,6$ or $\Z/3\rtimes\Z/4$; at $p=2$: cyclic of order $1,2,3,4,6$, $Q_8$ or $\SL_2(\F_3)$ \emph{(stated; reference)}.
\Cref{thm:reduction:kraus}.
\item \textbf{Tame additive reduction.} Additive potentially good reduction over a finite $K/\Q_p$ is tame ($f=2$) iff $p\ge5$, or $p=3$ and type
$\kod{III},\kod{III}^*,\kod{I}_0^*$, or $p=2$ and type $\kod{IV},\kod{IV}^*$. \emph{(stated; reference)} \Cref{thm:reduction:tame-additive}.
\item \textbf{Types $\kod{IV}$, $\kod{IV}^*$ at $2$.} Such a curve over $\Q_2$ has potentially good reduction, $G_{\mathrm{ss}}\simeq\Z/3$, the image of
$I_2$ on $T_\ell E$ ($\ell$ odd) is generated by $\sigma$ of order $3$ with eigenvalues $\zeta_3^{\pm1}$, unipotent on $E[3]$, $f_2=2$, and
$\Frob\,\sigma\,\Frob^{-1}=\sigma^{-1}$ \emph{(proved here)} \Cref{cor:reduction:IV-at-2}. Consequently (the paper link after it, with Lemma
lem:euler and the normalisation of $P_q$ in \Cref{ch:lfunctions}) $P_2(X)=1+2X$ and $e_\ell(E,2)=v_\ell(3/2)$ \emph{(consequence; see
\Cref{ch:lfunctions})}.
\item \textbf{The Tate curve.} For $v(\mathsf q)>0$, $E_{\mathsf q}\colon y^2+xy=x^3-5s_3x-(5s_3+7s_5)/12$ has $\Delta=\mathsf q\prod(1-\mathsf q^n)^{24}$,
$j=\mathsf q^{-1}+744+196884\mathsf q+\cdots$, split multiplicative reduction of type $\kod{I}_{v(\mathsf q)}$ \emph{(proved here)}, and
$E_{\mathsf q}(\bar K)\simeq\bar K^\times/\mathsf q^{\Z}$ Galois-equivariantly \emph{(stated; reference)}. \Cref{thm:reduction:tate-curve}.
\item \textbf{Inverting $j$.} $\mathsf q(J)=J+744J^2+750420J^3+872769632J^4+\cdots\in J\Z[[J]]$ inverts $1/j(\mathsf q)$; for $v(j_0)<0$, $\mathsf q(1/j_0)$
is the unique $\mathsf q_0$ with $v(\mathsf q_0)>0$, $j(E_{\mathsf q_0})=j_0$, and $v(\mathsf q_0)=-v(j_0)$. \emph{(proved here)}
\Cref{prop:reduction:j-inverse}.
\item \textbf{Curves with $v(j)<0$.} $E$ is the quadratic twist of $E_{q_E}$ by $\gamma=-c_4/c_6$: split iff $\gamma$ is a square, nonsplit iff
$K(\sqrt\gamma)/K$ is unramified quadratic, additive iff it is ramified; for multiplicative reduction $n=v(\Delta_{\min})=-v(j)=v(q_E)$.
\emph{(proof sketched here)} \Cref{thm:reduction:tate-uniformisation}.
\item \textbf{Tate parameter and unit part.} For $E/\Q_q$ with $v_q(j)<0$: $q_E=\mathsf q(1/j_E)$, $\nu=v_q(q_E)=-v_q(j_E)$ ($=n_q$ at a multiplicative
prime), $u=q_Eq^{-\nu}\in\Z_q^\times$, $c(u)$ the class of $\bar u$ in $\F_q^\times/(\F_q^\times)^{\ell^{v_\ell(q-1)}}$. These depend only on $j_E$, so
$t_\ell(E,q)=v_\ell(\nu)$ is twist invariant \emph{(proved here)}; under an isogeny of prime degree $d$, $\nu\mapsto d\nu$ or $\nu/d$, so $t_\ell$ is
invariant under isogenies of degree prime to $\ell$ \emph{(proof sketched here)}. \Cref{def:reduction:tate-parameter},
\Cref{prop:reduction:tate-parameter-invariance}.
\item \textbf{Torsion of the Tate curve.} $E_{\mathsf q}[\ell^t]=\{\zeta^a(\mathsf q^{1/\ell^t})^b\}$, an extension $0\to\mu_{\ell^t}\to E_{\mathsf q}[\ell^t]\to\Z/\ell^t\to0$;
on $T_\ell E_{\mathsf q}$, $\sigma\mapsto\sm{\chi_\ell(\sigma)}{\kappa(\sigma)}{0}{1}$ with $\kappa$ the Kummer cocycle of $\mathsf q$; for the twist by $\varepsilon$,
$\sm{\chi_\ell\varepsilon}{\varepsilon\kappa}{0}{\varepsilon}$, with corner $\varepsilon\kappa$. \emph{(proved here, given Tate's theorem)}
\Cref{prop:reduction:tate-torsion}.
\item \textbf{Types at additive potentially multiplicative primes.} For $p$ odd the type is $\kod{I}_\nu^*$, $\nu=-v(j)$, with $f=2$,
$v(\Delta_{\min})=\nu+6$, and conversely $\kod{I}_n^*$ ($n\ge1$) at odd $p$ forces $v(j)=-n$ \emph{(proof sketched here, given Ogg's formula)}.
For $K=\Q_2$ the type is $\kod{I}_m^*$ with $m=\nu+2f_2-4\in\{\nu+4,\nu+8\}$ \emph{(stated; reference)}. \Cref{thm:reduction:potmult-additive}.
\item \textbf{Unramified twists (the paper's Lemma lem:twist (i)).} A twist by a quadratic character unramified at $p$ preserves at $p$ the type,
$\Phig_p$, $v_p(\Delta_{\min})$, $f_p$ and $q_E$; $c_p$ may change. \emph{(proved here)} \Cref{lem:reduction:twist-unramified}.
\item \textbf{Ramified twists (the paper's Lemma lem:twist (ii)).} For odd $q$ the twist by $q^*$ exchanges $\kod{I}_n\leftrightarrow\kod{I}_n^*$ (same
$n\ge1$) and $\kod{I}_0\leftrightarrow\kod{I}_0^*$; for $q\ge5$ also $\kod{II}\leftrightarrow\kod{IV}^*$, $\kod{III}\leftrightarrow\kod{III}^*$,
$\kod{IV}\leftrightarrow\kod{II}^*$ with $v_q(\Delta_{\min})$ shifted by $6$ mod $12$; the conductor exponent of the twist at $q$ is $0$ (good), $1$
($\kod{I}_n$) or $2$ (additive). \emph{(proved here; the equality of the indices in $\kod{I}_n\leftrightarrow\kod{I}_n^*$ uses Ogg's formula)}
\Cref{lem:reduction:twist-ramified}.
\item \textbf{The exchanges at $3$.} The twist by $-3$ exchanges $\kod{II}\leftrightarrow\kod{IV}^*$, $\kod{III}\leftrightarrow\kod{III}^*$,
$\kod{IV}\leftrightarrow\kod{II}^*$ at $3$ and preserves $f_3$. \emph{(proved here, given Ogg's formula)} \Cref{prop:reduction:twist-at-3}.
\item \textbf{Twist-symmetrised weight.} For odd $q$ and odd $\ell$,
\[
w_\ell(E,q)=\max\bigl(v_\ell\#\Phig_q(E),\,v_\ell\#\Phig_q(E^{(q)})\bigr),\qquad w_\ell(E,2)=v_\ell\#\Phig_2(E).
\] \emph{(proved here)} \Cref{prop:reduction:twist-weight}.
\item \textbf{What changes at $2$ and $3$.} For $p\in\{2,3\}$: the type determines neither $v(\Delta_{\min})$ nor $f$ (only
$f=v(\Delta_{\min})+1-m$); $f\le8$ at $2$, $f\le5$ at $3$; $G_{\mathrm{ss}}$ can be non-abelian, of order divisible by $p$; the type determines
$G_{\mathrm{ss}}$ only for $\kod{IV},\kod{IV}^*$ at $2$ ($\Z/3$) and $\kod{III},\kod{III}^*,\kod{I}_0^*$ at $3$ ($\Z/4,\Z/4,\Z/2$); at $2$ a symbol
$\kod{I}_m^*$ may be potentially good, and if not, $m\ne-v_2(j)$. So the adjoint Euler factor is not a function of the type at $2$, $3$.
\emph{(summary of results proved or stated elsewhere in the chapter)} \Cref{rem:reduction:wild-overview}.
\item \textbf{Inertia at $2$.} For additive potentially good reduction over $\Q_2$, the ramification pattern of the roots of $\psi_3$ is $(1,1,1,1)$,
$(1,3,3,3)$, $(2,2,2,2)$, $(4,4,4,4)$ according as $G_{\mathrm{ss}}$ is $\Z/2$ (a ramified twist of good reduction), cyclic of order $3$ or $6$,
$\Z/4$, or $Q_8$ or $\SL_2(\F_3)$; order $3$ occurs exactly for $\kod{IV}$, $\kod{IV}^*$. \emph{(proved here, given Kraus's list of groups
and \Cref{thm:reduction:tame-additive})} \Cref{prop:reduction:inertia-at-2}.
\item \textbf{The prime $3$.} Types $\kod{III},\kod{III}^*,\kod{I}_0^*$ at $3$ are tame with $G_{\mathrm{ss}}$ cyclic of order $4,4,2$; types
$\kod{II},\kod{II}^*,\kod{IV},\kod{IV}^*$ are wild with $3\mid e$ and $f_3\in\{3,4,5\}$; $\kod{I}_n^*$ is potentially multiplicative with $f_3=2$.
\emph{(proved here, given \Cref{thm:reduction:tame-additive})} \Cref{rem:reduction:wild-at-3}.
\item \textbf{Symbols $\kod{I}_m^*$ at $2$.} They are potentially good in $61\%$ of the pairs below conductor $10{,}000$; when potentially
multiplicative, $t_\ell(E,2)=v_\ell(\nu)$ with $\nu=m-4$ or $m-8$, not $v_\ell(m)$. \emph{(proved here, with data)}
\Cref{rem:reduction:Instar-at-2}.
\item \textbf{Worked data.} Tate's algorithm by hand on \curve{121a1}, \curve{121c2}, \curve{539c1}, \curve{539d1}, \curve{726a1}, \curve{20a1},
\curve{99a1}, \curve{980i1}; minimal models of twists; semistability defects; Tate parameters; twist pairs; local data of the running examples;
what Cremona's tables contain. \Cref{ex:reduction:minimal-examples}, \Cref{ex:reduction:tate-running},
\Cref{ex:reduction:defects}, \Cref{ex:reduction:tate-parameters}, \Cref{ex:reduction:twist-pairs}, \Cref{ex:reduction:running-table},
\Cref{rem:reduction:cremona-tables}.
\end{enumerate}

\section*{Where the paper uses this chapter}

\begin{itemize}
\item \textbf{Definition def:terms} ($t_\ell(E,q)=v_\ell(-v_q(j_E))$, potentially multiplicative primes, the symbols $\kod{I}_n^*$ at $2$): \Cref{prop:reduction:minimal-basic},
\Cref{thm:reduction:potential-criterion}, \Cref{prop:reduction:j-inverse},
\Cref{def:reduction:tate-parameter}, \Cref{prop:reduction:tate-parameter-invariance}, \Cref{thm:reduction:potmult-additive},
\Cref{rem:reduction:Instar-at-2}.
\item \textbf{The table of Section 1 and Lemma lem:euler} (types at $q\ge5$, inertia images): \Cref{cor:reduction:tame-table},
\Cref{thm:reduction:serre-tate}, \Cref{thm:reduction:kraus}, \Cref{lem:reduction:twist-ramified}.
\item \textbf{Section 1.4} ($\Phig_q$ versus $c_q$, the exchange of types, the weight $w_\ell$, the inequalities (eq:main), (eq:twist)):
\Cref{prop:reduction:component-group}, \Cref{def:reduction:tamagawa}, \Cref{thm:reduction:kodaira-neron}, \Cref{lem:reduction:twist-ramified},
\Cref{prop:reduction:twist-at-3}, \Cref{prop:reduction:twist-weight}.
\item \textbf{Lemma lem:twist (i), (ii)}: proved as \Cref{lem:reduction:twist-unramified} and \Cref{lem:reduction:twist-ramified}, using
\Cref{prop:reduction:min-criterion}, \Cref{prop:reduction:I0star-cubic}, \Cref{thm:reduction:tate-uniformisation} and
\Cref{thm:reduction:potmult-additive}.
\item \textbf{Lemmas lem:tamagawa-exp and lem:tamfactor} (the Tate curve over the unramified quadratic extension, $q_E$, $u$, $c(u)$, the matrices
with corner $\varepsilon\kappa$, the check with $c_q$): \Cref{prop:reduction:split}, \Cref{prop:reduction:component-invariants},
\Cref{thm:reduction:tate-uniformisation} to \Cref{prop:reduction:tate-tamagawa}; the lemmas themselves are proved in \Cref{ch:galrep} and
\Cref{ch:blochkato}.
\item \textbf{Lemma lem:IV and Remark rem:types} ($E[3]^{I_q}$ a line or $0$, Kraus's defect, the case $q=2$):
\Cref{prop:reduction:component-invariants}, \Cref{thm:reduction:serre-tate}, \Cref{thm:reduction:kraus}, \Cref{cor:reduction:IV-at-2}.
\item \textbf{Proposition prop:shapes and Lemma lem:tamadd} of the revised paper (the image of inertia at a potentially good prime is finite,
its elements have orders $1,2,3,4,6$, it is cyclic for $q\ge5$; a potentially multiplicative curve is a quadratic twist of a Tate curve):
\Cref{thm:reduction:serre-tate}, \Cref{thm:reduction:kraus}, \Cref{thm:reduction:tate-uniformisation}, \Cref{prop:reduction:inertia-at-2}.
\item \textbf{Remark rem:wild, Section sec:wild, Proposition prop:forms}: \Cref{thm:reduction:tame-additive}, \Cref{cor:reduction:IV-at-2},
\Cref{rem:reduction:wild-overview}, \Cref{prop:reduction:inertia-at-2}, \Cref{rem:reduction:wild-at-3}, \Cref{rem:reduction:Instar-at-2}.
\item \textbf{Theorems thm:qual, thm:quant and the proof of thm:main} (twists at $\kod{I}_0^*$, $\kod{I}_n^*$, $\kod{II}$, $\kod{II}^*$ primes; conductor
exponents; the prime $2$ of type $\kod{I}_n^*$): \Cref{prop:reduction:conductor-properties}, \Cref{lem:reduction:twist-ramified},
\Cref{thm:reduction:tate-uniformisation}.
\item \textbf{The appendix} (\texttt{ellglobalred}, \texttt{R[4]}, \texttt{R[5]}, \texttt{K[i][2]}, the Kodaira codes, \texttt{valuation(j, q)}):
\Cref{conv:reduction:pari-codes}, \Cref{ex:reduction:pari-session}.
\item \textbf{The Example of Section 3, the Remark after it, Section sec:anomaly}: \Cref{ex:reduction:tate-running},
\Cref{ex:reduction:tate-parameters}, \Cref{ex:reduction:twist-pairs}, \Cref{ex:reduction:running-table}.
\end{itemize}

\section{Integral and minimal Weierstrass models}\label{sec:reduction:minimal}

This section fixes the local setting and the one choice that every later statement depends on: a
Weierstrass model of $E$ over a complete discretely valued field that is integral and has the smallest
possible discriminant. Every local invariant the paper reads off at a bad prime $q$ (the Kodaira type, the
valuations $v_q(\Delta)$ and $v_q(c_4)$, the Tamagawa number $c_q$) is defined through such a model. The
paper uses minimality explicitly in the proof of Lemma lem:twist (ii), where the twisted model ``is minimal
exactly when the result is less than $12$, otherwise $12$ is subtracted'': that sentence is
\Cref{prop:reduction:min-criterion}(ii) below.

\begin{convention}[the local setting]\label{conv:reduction:local}
Throughout the chapter $K$ is a field of characteristic $0$, complete for a discrete valuation $v$, with
valuation ring $\mathcal O$, maximal ideal $\mathfrak m=\varpi\mathcal O$, perfect residue field $k=\mathcal O/\mathfrak m$ of
characteristic $p$, and $v(\varpi)=1$. In every application $K$ is $\Q_p$, a finite extension of $\Q_p$, or a
finite unramified extension of one of these; $K^{\ur}$ is the maximal unramified extension and $I_K$ the
inertia group (\Cref{ch:local}). When the paper's bad prime is meant we write $K=\Q_q$ and $p=q$. The letter
$q$ is also used for the Tate parameter in the literature; in this chapter the variable of the Tate curve is
written $\mathsf q$ and the Tate curve $E_{\mathsf q}$, and the Tate parameter of a given curve $E$ is always
$q_E$ (\Cref{sec:reduction:tate}). The arithmetic Frobenius is $\Frob$; the geometric Frobenius at a named
prime is written $\Frob_q^{-1}$, because $\Phi_q(\cdot)$ denotes a component group in this chapter (Appendix B,
\Cref{ch:notation}, conventions B-C2 and B-C11).
\end{convention}

\begin{convention}[models]\label{conv:reduction:models}
A \emph{model} of an elliptic curve $E/K$ is a Weierstrass equation
\[
y^2+a_1xy+a_3y=x^3+a_2x^2+a_4x+a_6,\qquad a_i\in K,
\]
with the general coefficients $a_1,\dots,a_6$ and the quantities $b_2,b_4,b_6,b_8,c_4,c_6,\Delta,j$ of
\Cref{ch:ellcurves}. Two models of $E$ are related by a change of coordinates
$x=u^2x'+r$, $y=u^3y'+su^2x'+t$ with $u\in K^\times$, $r,s,t\in K$, written $(u,r,s,t)$, and then
\[
u^4c_4'=c_4,\qquad u^6c_6'=c_6,\qquad u^{12}\Delta'=\Delta,\qquad j'=j,
\]
together with the formulas for the $a_i'$ and $b_i'$ recalled in \Cref{ch:ellcurves}. At $p=2$ and $p=3$ a
short model $y^2=x^3+Ax+B$ is in general not minimal, so short models are used only when $p\ge5$.
\end{convention}

\begin{definition}[integral and minimal models]\label{def:reduction:minimal}
A model of $E/K$ is \emph{integral} if all $a_i\in\mathcal O$; then $b_2,b_4,b_6,b_8,c_4,c_6,\Delta\in\mathcal O$, as they
are polynomials with integer coefficients in the $a_i$. An integral model is \emph{minimal} if $v(\Delta)$ is
smallest among all integral models of $E$; its discriminant is the \emph{minimal discriminant}
$\Delta_{\min}$, well defined up to $\mathcal O^\times$. For a minimal model the invariant differential
$\omega=dx/(2y+a_1x+a_3)$ is the \emph{N\'eron differential}.
\end{definition}

The reason for minimising is that the reduction of a model modulo $\varpi$ depends on the model, and a careless choice
of coordinates makes a perfectly good curve look singular. Let $p\ge5$ and take $y^2=x^3+x+1$ over $\Q_p$; its
discriminant is $-16\cdot31=-496$, so for $p\ne31$ it reduces to a smooth cubic. The change $(p^{-1},0,0,0)$ multiplies
$a_i$ by $p^i$ and gives the integral model $y^2=x^3+p^4x+p^6$ of the same curve, with discriminant $p^{12}\cdot(-496)$;
it reduces to the cuspidal cubic $y^2=x^3$. Since $u^{12}\Delta'=\Delta$, the valuations $v(\Delta)$ of the integral
models of $E$ all lie in one class modulo $12$, and a change of coordinates can only add multiples of $12$. Minimising
removes these spurious factors $\varpi^{12}$, and \Cref{prop:reduction:minimal-basic} below shows that what is left (the
reduction of a minimal model, up to a change of coordinates over $k$) depends only on $E$.

\begin{proposition}[existence and uniqueness of minimal models]\label{prop:reduction:minimal-basic}
Let $E/K$ be an elliptic curve.
\begin{enumerate}[label=(\roman*)]
\item $E$ has a minimal model.
\item Two minimal models of $E$ are related by a change $(u,r,s,t)$ with $u\in\mathcal O^\times$ and
$r,s,t\in\mathcal O$.
\item The N\'eron differential is well defined up to multiplication by $\mathcal O^\times$. The integers
$v(\Delta_{\min})$, $v(c_4)$ and $v(c_6)$ of a minimal model depend only on $E/K$, and
$v(j)=3v(c_4)-v(\Delta)$ is the same for every model.
\end{enumerate}
\end{proposition}

\begin{proof}
(i) Start from any model. The change $(u,0,0,0)$ multiplies $a_i$ by $u^{-i}$; with $u=\varpi^{-n}$ the new
coefficients are $\varpi^{ni}a_i$, which lie in $\mathcal O$ for $n$ large. So integral models exist, the values
$v(\Delta)$ of integral models are nonnegative integers, and a smallest one is attained.

(ii) Any two models of $E$ are related by some $(u,r,s,t)$ (\Cref{ch:ellcurves}). If both are minimal, then
$v(\Delta)=v(\Delta')$ and $u^{12}\Delta'=\Delta$ give $v(u)=0$. The transformation formulas for $b_6$ and $b_8$
read
\begin{align*}
4r^3+b_2r^2+2b_4r+b_6&=u^6b_6',\\
3r^4+b_2r^3+3b_4r^2+3b_6r+b_8&=u^8b_8',
\end{align*}
with right-hand sides in $\mathcal O$. Suppose $v(r)<0$. If $p\ne3$, in the second equation the term $3r^4$ has
valuation $4v(r)$, strictly less than the valuation of every other term ($\ge3v(r)$, $2v(r)$, $v(r)$ or $0$),
which is impossible. If $p=3$ then $4$ is a unit and the same argument applies to $4r^3$ in the first
equation. So $r\in\mathcal O$. Next, $ua_1'=a_1+2s$ and $u^2a_2'=a_2-sa_1+3r-s^2$ show that $s$ is a root of the
monic polynomial $X^2+a_1X+(u^2a_2'-a_2-3r)$ with coefficients in $\mathcal O$; as $\mathcal O$ is integrally closed,
$s\in\mathcal O$. Likewise $u^6a_6'=a_6+ra_4+r^2a_2+r^3-ta_3-t^2-rta_1$ shows that $t$ is a root of a monic
quadratic over $\mathcal O$, so $t\in\mathcal O$.

(iii) Under $(u,r,s,t)$ the invariant differential is multiplied by $u$ (\Cref{ch:ellcurves}), and $u$ is a unit
by (ii). The relations $u^4c_4'=c_4$, $u^6c_6'=c_6$, $u^{12}\Delta'=\Delta$ with $u$ a unit give the invariance of
the three valuations. Finally $j=c_4^3/\Delta$ does not depend on the model at all.
\end{proof}

\begin{proposition}[criteria for minimality]\label{prop:reduction:min-criterion}
Let a model of $E/K$ be integral.
\begin{enumerate}[label=(\roman*)]
\item (any $p$) If $v(\Delta)<12$, or $v(c_4)<4$, or $v(c_6)<6$, the model is minimal.
\item ($p\ge5$) The model is minimal if and only if $v(\Delta)<12$ or $v(c_4)<4$. If it is not minimal, the
model $y^2=x^3-27\varpi^{-4}c_4x-54\varpi^{-6}c_6$ is integral with discriminant of valuation $v(\Delta)-12$;
repeating this step finitely often gives a minimal model.
\item (Kraus's conditions; $K=\Q$) Let $c_4,c_6\in\Z$ with $\Delta=(c_4^3-c_6^2)/1728\in\Z\setminus\{0\}$. There is a
Weierstrass model with coefficients in $\Z$ whose invariants are $c_4,c_6$ if and only if $v_3(c_6)\ne2$, and
either $c_6\equiv-1\pmod4$, or $v_2(c_4)\ge4$ and $c_6\equiv0$ or $8\pmod{32}$. Consequently, for
$p\in\{2,3\}$, an integral model over $\Z$ with $p^4\mid c_4$, $p^6\mid c_6$ and $p^{12}\mid\Delta$ fails to be
minimal at $p$ exactly when $(c_4/p^4,c_6/p^6)$ satisfies the condition at $p$ ($v_3(c_6/3^6)\ne2$ for $p=3$;
the congruence condition for $p=2$).
\end{enumerate}
\end{proposition}

\begin{proof}
(i) Suppose an integral model with invariants $c_4',c_6',\Delta'$ and $v(\Delta')<v(\Delta)$ exists, related by
$(u,r,s,t)$. Then $u^{12}\Delta'=\Delta$ forces $v(u)\ge1$, hence $v(\Delta)=v(\Delta')+12v(u)\ge12$,
$v(c_4)=v(c_4')+4v(u)\ge4$ and $v(c_6)=v(c_6')+6v(u)\ge6$, as $c_4',c_6',\Delta'$ are integral. Each of the three
inequalities in the hypothesis contradicts one of these.

(ii) By (i) the condition is sufficient. Conversely let $v(c_4)\ge4$ and $v(\Delta)\ge12$. From
$c_6^2=c_4^3-1728\Delta$ we get $v(c_6^2)\ge12$, so $v(c_6)\ge6$. The curve has the model $y^2=x^3-27c_4x-54c_6$
(\Cref{ch:ellcurves}), whose discriminant is $6^{12}\Delta$. (For a short model $y^2=x^3+Ax+B$ one has $b_2=0$, $b_4=2A$,
$b_6=4B$, hence $c_4=-48A$ and $c_6=-864B$; with $A=-27c_4$, $B=-54c_6$ this gives the invariants $1296c_4=6^4c_4$ and
$46656c_6=6^6c_6$, so the short model is obtained by a change $(u,r,s,t)$ with $u=1/6$, and its discriminant is
$u^{-12}\Delta=6^{12}\Delta$.) Applying $(\varpi,0,0,0)$ to it gives
$y^2=x^3-27\varpi^{-4}c_4x-54\varpi^{-6}c_6$, which is integral, with discriminant $6^{12}\varpi^{-12}\Delta$. As
$p\ge5$, $6$ is a unit and the valuation is $v(\Delta)-12$. So the model was not minimal. Each step lowers
$v(\Delta)$ by $12$, so the process stops.

(iii) We do not prove this; see Kraus \cite{Kraus89}, and \cite[Section 3.2]{Cremona97} for its use in
algorithms. The consequence follows by applying the criterion to the invariants $c_4/p^4$, $c_6/p^6$ of the
model obtained by $(p,0,0,0)$: the conditions at the other prime are unchanged, since $p$ is a unit there
(for $p=3$: multiplication by $3^{-6}\equiv9\pmod{32}$ preserves the classes $0$ and $8$ modulo $32$ and the class $-1$ modulo $4$).
\end{proof}

We checked (iii) on the data: $20{,}000$ random integral models with coefficients in $[-50,50]$ all satisfy
Kraus's conditions, and in all $8{,}248$ pairs (curve, $p\in\{2,3\}$) of Cremona's tables below conductor
$10{,}000$ whose minimal model has $p^4\mid c_4$, $p^6\mid c_6$, $p^{12}\mid\Delta$, the scaled invariants fail
the condition at $p$, as they must.

Part (i) is only a sufficient condition when $p\in\{2,3\}$, and (iii) is what replaces the converse of (ii) there. For
example \curve{32a1} is $y^2=x^3+4x$, with $c_4=-48\cdot4=-2^6\cdot3$, $c_6=0$ and $\Delta=-16\cdot4\cdot4^3=-2^{12}$: all three
valuations at $2$ reach the thresholds of (i), and yet the model is minimal at $2$. Indeed a model with smaller
discriminant would come from $(2,r,s,t)$ and have invariants $c_4/2^4=-12$, $c_6/2^6=0$ and $\Delta/2^{12}=-1$; these fail
Kraus's condition at $2$ ($0\not\equiv-1\pmod4$, and $v_2(-12)=2<4$), so no integral model has them.

\begin{theorem}[global minimal models]\label{thm:reduction:global-minimal}
Every elliptic curve $E/\Q$ has a model with coefficients in $\Z$ that is minimal at every prime $p$. It is
unique up to a change $(u,r,s,t)$ with $u=\pm1$ and $r,s,t\in\Z$, and its discriminant $\Delta_E$ is the
\emph{minimal discriminant} of $E$. Imposing $a_1,a_3\in\{0,1\}$ and $a_2\in\{-1,0,1\}$ makes it unique.
\end{theorem}

\begin{proof}[Sketch of proof]
Start from a model over $\Z$. At each of the finitely many primes $p$ dividing its discriminant choose a
change $(u_p,r_p,s_p,t_p)$ over $\Q_p$ to a minimal model. Since $\Z$ is a principal ideal domain the ideal
$\prod_pp^{v_p(u_p)}$ is generated by $u=\prod_pp^{v_p(u_p)}$, and the Chinese remainder theorem lets one choose
$r,s,t\in\Z$ approximating the $r_p,s_p,t_p$ closely enough that the change $(u,r,s,t)$ is integral and
minimal at every $p$ simultaneously. Uniqueness follows from
\Cref{prop:reduction:minimal-basic}(ii) at every prime: $u$ is a unit of $\Z$ and $r,s,t\in\Z_p$ for all $p$.
The normalisation of $a_1,a_3,a_2$ uses $s$, $t$ and $r$ modulo $2$, $2$ and $3$. See
\cite[VIII.8]{SilvermanAEC} and \cite[Section 3.2]{Cremona97}.
\end{proof}

Cremona's tables \cite{Cremona97,ecdata} and PARI's \texttt{ellminimalmodel} give exactly this normalised
global minimal model. All models printed in this book are of this form.

\begin{example}[minimal models]\label{ex:reduction:minimal-examples}
(a) The curve \curve{121a1} is $[1,1,1,-30,-76]$ with $\Delta=-11^2$. Its twist by $-11$, written with the
$b$-invariants as in \Cref{ch:ellcurves}, is $y^2=x^3+db_2x^2+8d^2b_4x+16d^3b_6$ with $d=-11$; its discriminant
is $2^{12}d^6\Delta$, so $v_{11}(\Delta)=2+6=8<12$ and the model is minimal at $11$ by
\Cref{prop:reduction:min-criterion}(i). It is the curve \curve{121c2}, with $\Delta=-11^8$. Twisting
\curve{121c2} back by $-11$ gives $v_{11}(\Delta)=8+6=14$ and $v_{11}(c_4)=3+2=5\ge4$, so that model is not
minimal at $11$; one step of \Cref{prop:reduction:min-criterion}(ii) gives $v_{11}(\Delta)=2$, and the curve is
\curve{121a1} again.

(b) The curve \curve{140b1} is $[0,0,0,32,212]$, with $v_2(\Delta)=8$, $v_2(c_4)=9$ and $v_2(j)=27-8=19$. It is
minimal at $2$ by \Cref{prop:reduction:min-criterion}(i), although $c_4$ is very divisible by $2$.

(c) The short model $y^2=x^3-27c_4x-54c_6$ of \curve{11a1} has discriminant $6^{12}\Delta$, so
$v_2(\Delta)=v_3(\Delta)=12$; it is not minimal at $2$ or $3$, since the model $[0,-1,1,-10,-20]$ has
$\Delta=-11^5$.
\end{example}

\begin{paperlink}
Lemma lem:twist (ii), proof: ``the twisted model has $v_q(\Delta)$ increased by $6$, and is minimal exactly when
the result is less than $12$, otherwise $12$ is subtracted''. For $q\ge5$ and potentially good reduction this is
\Cref{prop:reduction:min-criterion}: if the result is at least $12$, then $v_q(j)\ge0$ gives
$3v_q(c_4)\ge v_q(\Delta)\ge12$, so $v_q(c_4)\ge4$ and part (ii) applies once. The full argument is
\Cref{lem:reduction:twist-ramified}.
\end{paperlink}

\section{Reduction modulo the uniformiser}\label{sec:reduction:reduction}

This section defines good, multiplicative and additive reduction, splits multiplicative reduction into split
and nonsplit, and studies the reduction map on points. The paper uses the trichotomy everywhere (the numbers
$\delta_p=\dim V^{I_p}=2,1,0$ of its Section 1.1; the PARI test \texttt{k >= 5} of its appendix), the values
$a_p\in\{1,-1,0\}$ at bad primes (Lemma lem:old, Definition def:terms), and the split/nonsplit distinction in
Lemmas lem:tamagawa-exp and lem:tamfactor. The injectivity of reduction on prime-to-$p$ torsion
(\Cref{cor:reduction:torsion-injects}) is the easy half of the criterion of N\'eron, Ogg and Shafarevich
(\Cref{ch:galrep}) and the input of \Cref{prop:reduction:component-invariants}.

\begin{definition}[reduction types]\label{def:reduction:types}
Let $y^2+a_1xy+a_3y=x^3+a_2x^2+a_4x+a_6$ be a minimal model of $E/K$ and let $\tilde E$ be the plane cubic over
$k$ defined by the reduced coefficients $\tilde a_i$. The curve $E$ has
\begin{itemize}
\item \emph{good reduction} if $\tilde E$ is smooth, that is, $v(\Delta)=0$;
\item \emph{multiplicative reduction} if $\tilde E$ has a node, that is, $v(\Delta)>0$ and $v(c_4)=0$;
\item \emph{additive reduction} if $\tilde E$ has a cusp, that is, $v(\Delta)>0$ and $v(c_4)>0$.
\end{itemize}
Multiplicative reduction is \emph{split} if the two tangent directions at the node are defined over $k$, and
\emph{nonsplit} otherwise. \emph{Semistable} means good or multiplicative. For $E/\Q$ and a prime $p$ we say
that $E$ has good (multiplicative, additive) reduction at $p$ if $E/\Q_p$ has.
\end{definition}

The definition does not depend on the minimal model: two minimal models differ by $(u,r,s,t)$ with
$u\in\mathcal O^\times$ and $r,s,t\in\mathcal O$ (\Cref{prop:reduction:minimal-basic}), which reduces to a change of
coordinates over $k$. A Weierstrass cubic is smooth if and only if its discriminant is nonzero, and if
$\tilde\Delta=0$ it has exactly one singular point, a node if $\tilde c_4\ne0$ and a cusp if $\tilde c_4=0$
(\Cref{ch:ellcurves}). The singular point is $k$-rational: it is unique, hence fixed by $\Gal(\bar k/k)$, and
$k$ is perfect.

\begin{proposition}[the group of nonsingular points]\label{prop:reduction:Ens}
Let $\tilde E/k$ be a Weierstrass cubic with a singular point $S$, and let $\tilde E_{\mathrm{ns}}(k)$ be the set of
$k$-points other than $S$. The chord-tangent construction makes $\tilde E_{\mathrm{ns}}(k)$ an abelian group with
neutral element $\tilde O$, and
\begin{enumerate}[label=(\roman*)]
\item if $S$ is a cusp, $\tilde E_{\mathrm{ns}}(k)\simeq k^+$;
\item if $S$ is a split node, $\tilde E_{\mathrm{ns}}(k)\simeq k^\times$;
\item if $S$ is a nonsplit node and $l/k$ is the quadratic extension containing the tangent slopes,
$\tilde E_{\mathrm{ns}}(k)\simeq\ker(\Nm_{l/k}\colon l^\times\to k^\times)$.
\end{enumerate}
For $k=\F_p$ the orders are $p$, $p-1$ and $p+1$.
\end{proposition}

\begin{proof}
A line through two nonsingular points does not pass through $S$: a line through $S$ meets the cubic there with
multiplicity at least $2$, hence in at most one further point. So the chord-tangent law, defined exactly as for
a smooth cubic (\Cref{ch:ellcurves}), stays inside $\tilde E_{\mathrm{ns}}$; associativity can be checked as in the
smooth case, and also follows from the isomorphisms below, which we now construct.

In both cases we first translate $S$ to $(0,0)$ (a change $(1,r,0,t)$ over $k$; $S$ is $k$-rational). Then $\tilde a_6=0$
because $(0,0)$ lies on the curve, and $\tilde a_4=\tilde a_3=0$ because the partial derivatives
$\partial F/\partial x=a_1y-3x^2-2a_2x-a_4$ and $\partial F/\partial y=2y+a_1x+a_3$ of
$F=y^2+a_1xy+a_3y-x^3-a_2x^2-a_4x-a_6$ vanish at $(0,0)$. So the equation is $y^2+a_1xy-a_2x^2=x^3$, and the quadratic
form $y^2+a_1xy-a_2x^2$ is the tangent cone at $S$: a product of two distinct linear forms at a node, a square at a cusp.

(i) At a cusp the quadratic part is $(y-\alpha x)^2$, with $\alpha=-a_1/2$ if $\operatorname{char}k\ne2$; if $\operatorname{char}k=2$
then $a_1=0$ (the discriminant $a_1^2+4a_2=a_1^2$ vanishes) and $\alpha=\sqrt{a_2}$, which lies in $k$ because $k$ is perfect.
The change $y\mapsto y+\alpha x$ makes the single tangent direction the line $y=0$, and the curve becomes $y^2=x^3$, that
is, $Y^2Z=X^3$ in projective coordinates. Its points other than
$S=[0:0:1]$ have $Y\ne0$ and are $[t:1:t^3]$, $t\in k$, with $\tilde O=[0:1:0]$ at $t=0$. A line
$\alpha X+\beta Y+\gamma Z=0$ not through $S$ has $\gamma\ne0$ and meets the curve at the points with
$\gamma t^3+\alpha t+\beta=0$, counted with multiplicity. The coefficient of $t^2$ is $0$, so the three values
of $t$ add up to $0$. Three points add up to $\tilde O$ in the group law exactly when they are the intersection
of a line with the curve, so $P\mapsto t(P)$ is a homomorphism $\tilde E_{\mathrm{ns}}(k)\to k^+$, and it is
bijective.

(ii) With $S$ at $(0,0)$ as above, the quadratic part factors as $(y-\alpha x)(y-\beta x)$ with $\alpha\ne\beta$ the roots of
$T^2+a_1T-a_2$; split means $\alpha,\beta\in k$. A nonsingular point $P\ne\tilde O$ has $x\ne0$ (the only point of the curve
with $x=0$ is $S$, as $y^2=0$ there), so it lies on the line $y=\lambda x$ through $S$ with $\lambda=y/x$. Substituting
$y=\lambda x$ in the equation gives $x^2(\lambda^2+a_1\lambda-a_2)=x^3$, so $x=(\lambda-\alpha)(\lambda-\beta)$ and $y=\lambda x$;
conversely every $\lambda\in k\setminus\{\alpha,\beta\}$ gives such a point. So $P\mapsto\lambda$ is a bijection from
$\tilde E_{\mathrm{ns}}(k)\setminus\{\tilde O\}$ to $k\setminus\{\alpha,\beta\}$, and $\tilde O$ corresponds to
$\lambda=\infty$. Put $\phi(P)=(\lambda-\alpha)/(\lambda-\beta)$ and $\phi(\tilde O)=1$; this is a bijection onto
$k^\times$. A line $y=mx+c_0$ with $c_0\ne0$ meets the curve at the points with $\lambda x=mx+c_0$, that is, at the
roots of $(\lambda-\alpha)(\lambda-\beta)(\lambda-m)-c_0=\prod_{i=1}^3(\lambda-\lambda_i)$ (the line does not pass through
$S$, which is why $c_0\ne0$). Evaluating at $\lambda=\alpha$ and at $\lambda=\beta$ gives $\prod(\alpha-\lambda_i)=-c_0=\prod(\beta-\lambda_i)$,
hence $\prod_i(\lambda_i-\alpha)=c_0=\prod_i(\lambda_i-\beta)$, and dividing, $\prod_i\phi(P_i)=1$. A vertical line
$x=x_0$ ($x_0\ne0$) meets the curve at $\tilde O$ and at the two roots of $(\lambda-\alpha)(\lambda-\beta)-x_0$, and the same
evaluation gives $\phi(P_1)\phi(P_2)=1=\phi(\tilde O)^{-1}$. So in all cases three collinear points have
$\phi(P_1)\phi(P_2)\phi(P_3)=1$. Since $P$, $-P$, $\tilde O$ are collinear, $\phi(-P)=\phi(P)^{-1}$; and if $P,Q,R$ are collinear then
$P+Q=-R$, so $\phi(P+Q)=\phi(R)^{-1}=\phi(P)\phi(Q)$. So $\phi$ is a homomorphism to $k^\times$. (The same two lines, with sums in
place of products, turn ``the three values of $t$ add up to $0$'' into the homomorphism property in (i).)

(iii) Apply (ii) over $l$, where $\alpha,\beta$ are conjugate under the generator $\sigma$ of $\Gal(l/k)$. For a point
$P$ with slope $\lambda\in l$, the conjugate point has slope $\sigma(\lambda)$, and
$\phi(\sigma P)=(\sigma\lambda-\alpha)/(\sigma\lambda-\beta)=\sigma\bigl((\lambda-\beta)/(\lambda-\alpha)\bigr)=\sigma(\phi(P))^{-1}$. So $P$
is $k$-rational if and only if $\phi(P)\sigma(\phi(P))=1$, that is, $\Nm_{l/k}\phi(P)=1$.

For $k=\F_p$: $\#\F_p=p$, $\#\F_p^\times=p-1$, and the norm $\F_{p^2}^\times\to\F_p^\times$ is surjective, so its
kernel has $(p^2-1)/(p-1)=p+1$ elements.
\end{proof}

\begin{corollary}[$a_p$ at a bad prime]\label{cor:reduction:ap-bad}
Let $E/\Q$ have bad reduction at $p$ and put $a_p=p+1-\#\tilde E(\F_p)$, the singular point included in the
count. Then $a_p=1$, $-1$, $0$ according as the reduction is split multiplicative, nonsplit multiplicative or
additive.
\end{corollary}

\begin{proof}
$\#\tilde E(\F_p)=\#\tilde E_{\mathrm{ns}}(\F_p)+1$, and \Cref{prop:reduction:Ens} gives $p$, $p-1$, $p+1$ for the three
cases.
\end{proof}

\begin{proposition}[split and nonsplit]\label{prop:reduction:split}
Let $E/K$ have multiplicative reduction.
\begin{enumerate}[label=(\roman*)]
\item After a change $(1,r,0,t)$ with $r,t\in\mathcal O$ moving the node to $(0,0)$, the reduction is split if and
only if $T^2+a_1T-a_2$ has a root in $k$.
\item If $p$ is odd, the reduction is split if and only if $-c_6$ is a square in $k^\times$ for one (equivalently
every) minimal model; equivalently, if and only if $-c_4/c_6$ is a square in $k^\times$.
\item If $p$ is odd and $k$ is finite, the twist of $E$ by the unramified quadratic extension of $K$ exchanges split and
nonsplit reduction, and over the unramified quadratic extension of $K$ the reduction is split.
\end{enumerate}
\end{proposition}

\begin{proof}
(i) is the definition, by the factorisation of the quadratic part in the proof of \Cref{prop:reduction:Ens}(ii).
(ii) After moving the node to $(0,0)$, $a_3,a_4,a_6\in\mathfrak m$, so $b_4\equiv0$ and $b_6\equiv0$, and
$c_4=b_2^2-24b_4\equiv b_2^2$, $c_6=-b_2^3+36b_2b_4-216b_6\equiv-b_2^3$. The discriminant of $T^2+a_1T-a_2$ is
$a_1^2+4a_2=b_2$, which is nonzero mod $\mathfrak m$ because $c_4$ is a unit. So split means $\tilde b_2$ is a
square, and $-\tilde c_6=\tilde b_2\cdot\tilde b_2^2$ is a square if and only if $\tilde b_2$ is, as is
$-\tilde c_4/\tilde c_6=\tilde b_2^{-1}$. Another minimal model changes $c_6$ by $u^6$ and $c_4/c_6$ by $u^{-2}$, with
$u$ a unit, which does not change the square classes.
(iii) Since $k$ is finite, $k^\times/(k^\times)^2$ has order $2$, so $K$ has a unique unramified quadratic extension
$K'=K(\sqrt d)$, with $d$ a unit whose reduction is not a square. For $p$ odd we may take a model $y^2=x^3+a_2x^2+a_4x+a_6$; its twist by $d$ is
$y^2=x^3+da_2x^2+d^2a_4x+d^3a_6$ (\Cref{ch:ellcurves}), which multiplies $c_6$ by $d^3$. As $\tilde d$ is not a square, this changes the square class of $-\tilde c_6$, and the twisted model is still
minimal ($v(\Delta)$ is unchanged). Over the unramified quadratic extension $K'$ every element of $k^\times$
becomes a square in the residue field $k'$, so the reduction is split there.
\end{proof}

We checked (ii) on all $127{,}592$ pairs (curve, odd prime of multiplicative reduction) of Cremona's tables below
conductor $10{,}000$: in every case $a_p=\bigl(\frac{-c_6}{p}\bigr)$.

\begin{paperlink}
Lemma lem:tamagawa-exp ends with ``This holds for split and for nonsplit reduction'', and its proof passes to
``the unramified quadratic extension of $\Q_q$ (which does not change the inertia group)'', over which $E$ is
split by \Cref{prop:reduction:split}(iii) and a Tate curve by \Cref{thm:reduction:tate-uniformisation}. In Lemma
lem:tamfactor the character $\varepsilon$ is unramified of order at most $2$ with $\varepsilon(\Frob_q)=a_q$, and in Lemma
lem:old $\varepsilon(\Frob_p)=a_p(f)=\pm1$ by \Cref{cor:reduction:ap-bad}.
\end{paperlink}

\begin{definition}[the reduction filtration]\label{def:reduction:filtration}
For a minimal model of $E/K$ and $P\in E(K)$, write $P=[X:Y:Z]$ with $X,Y,Z\in\mathcal O$ not all in $\mathfrak m$, and
let $\tilde P=[\tilde X:\tilde Y:\tilde Z]\in\tilde E(k)$ be its reduction. Put
\[
E_0(K)=\{P\in E(K):\tilde P\in\tilde E_{\mathrm{ns}}(k)\},\qquad E_1(K)=\{P\in E(K):\tilde P=\tilde O\}.
\]
\end{definition}

The point of the definition is that reduction is a group homomorphism only where it lands in a group: $\tilde E(k)$ is not a
group when $\tilde E$ is singular, but $\tilde E_{\mathrm{ns}}(k)$ is (\Cref{prop:reduction:Ens}). So $E_0(K)$, the set of points whose
reduction avoids the singular point, is the natural domain of a reduction homomorphism, and $E_1(K)$, the points $p$-adically close to $O$ (in affine coordinates, those with
$v(x)<0$), is its kernel. The three pieces $E_1(K)$, $E_0(K)/E_1(K)$ and $E(K)/E_0(K)$ are then studied separately: the first by the
formal group, the second by \Cref{prop:reduction:Ens}, the third by the component group of \Cref{sec:reduction:neron}.

\begin{proposition}[the reduction map]\label{prop:reduction:filtration}
$E_0(K)$ is a subgroup of $E(K)$, $E_1(K)$ is a subgroup of $E_0(K)$, and reduction is a surjective homomorphism
$E_0(K)\to\tilde E_{\mathrm{ns}}(k)$ with kernel $E_1(K)$:
\[
0\longrightarrow E_1(K)\longrightarrow E_0(K)\longrightarrow\tilde E_{\mathrm{ns}}(k)\longrightarrow0 .
\]
If $E$ has good reduction, $E_0(K)=E(K)$.
\end{proposition}

\begin{proof}
Reduction is well defined, since the primitive coordinates are unique up to $\mathcal O^\times$. Clearly
$\tilde O=[0:1:0]$ is nonsingular, and $-P=(x,-y-a_1x-a_3)$ reduces to $-\tilde P$.
Let $P_1,P_2\in E_0(K)$ and let $P_3$ be the third intersection of the line $L$ through them (the tangent if
$P_1=P_2$) with $E$. Write $L$ as $\alpha X+\beta Y+\gamma Z=0$ with $\alpha,\beta,\gamma\in\mathcal O$ not all in $\mathfrak m$; its
reduction $\tilde L$ is a line. Parametrising $L$ linearly and substituting into the cubic gives a cubic form in two
variables with coefficients in $\mathcal O$ whose three roots are $P_1,P_2,P_3$; reducing it shows that $\tilde L$
meets $\tilde E$ exactly in $\tilde P_1,\tilde P_2,\tilde P_3$, counted with multiplicity (the reduced form is not
identically zero, as $\tilde E$ is irreducible and does not contain the line $\tilde L$). If $\tilde P_3$ were the
singular point, $\tilde L$ would meet $\tilde E$ with multiplicity at least $2$ there and also at the nonsingular
points $\tilde P_1,\tilde P_2$, four intersections in all, which is impossible. So $\tilde P_3$ is nonsingular and
$\tilde P_1+\tilde P_2+\tilde P_3=\tilde O$ in $\tilde E_{\mathrm{ns}}(k)$. Hence $E_0(K)$ is closed under addition
and negation, and reduction is a homomorphism on it; its kernel is $E_1(K)$ by definition.

Surjectivity: let $(\tilde x_0,\tilde y_0)$ be a nonsingular affine point of $\tilde E(k)$ (the point $\tilde O$ is
the image of $O$). With $F(x,y)=y^2+a_1xy+a_3y-x^3-a_2x^2-a_4x-a_6$, one of $\partial F/\partial x$,
$\partial F/\partial y$ is nonzero at the point. If $\partial F/\partial y(\tilde x_0,\tilde y_0)\ne0$, choose
any lift $x_0\in\mathcal O$ of $\tilde x_0$; the polynomial $F(x_0,y)$ has the simple root $\tilde y_0$ modulo $\mathfrak m$,
and Hensel's lemma (\Cref{ch:local}) lifts it to $y_0\in\mathcal O$ with $F(x_0,y_0)=0$. The other case is the same
with the roles of $x$ and $y$ exchanged. The point $(x_0,y_0)\in E(K)$ reduces to $(\tilde x_0,\tilde y_0)$.
\end{proof}

\begin{theorem}[the formal group]\label{thm:reduction:formal-group}
Let a minimal model of $E/K$ be given.
\begin{enumerate}[label=(\roman*)]
\item In the parameter $z=-x/y$ near $O$, the curve is described by a power series $w(z)=-1/y=z^3+a_1zw+\cdots$
with coefficients in $\Z[a_1,\dots,a_6]$, and the group law by a power series $F(z_1,z_2)\in\Z[a_1,\dots,a_6][[z_1,z_2]]$;
$F$ is a one-parameter commutative formal group $\hat E$, and $P\mapsto z(P)$ is an isomorphism
$E_1(K)\simeq\hat E(\mathfrak m)$, the set $\mathfrak m$ with the group law $F$.
\item For $r>v(p)/(p-1)$ the formal logarithm gives an isomorphism $\hat E(\mathfrak m^r)\simeq\mathfrak m^r$ with
the additive group.
\item For every integer $m$ prime to $p$, multiplication by $m$ is a bijection of $E_1(K)$. In particular $E_1(K)$
has no nonzero torsion of order prime to $p$, and it is uniquely $m$-divisible.
\end{enumerate}
\end{theorem}

\begin{proof}[Sketch of proof]
(i) and (ii) are \cite[IV.1 to IV.6 and VII.2]{SilvermanAEC}: one solves the equation of the curve for $w$ as a
power series in $z$ by iteration, computes the third intersection of a line with the curve in these
coordinates, and checks that the resulting series converge on $\mathfrak m$ and match the reduction filtration.
(iii) is a complete argument given (i): a formal group law has the form $F(z_1,z_2)=z_1+z_2+(\text{terms of degree}\ge2)$,
so by induction $[m](z)=F([m-1](z),z)=(m-1)z+z+\cdots$ is $mz+(\text{terms of degree}\ge2)$ for $m\ge1$, and
$[-1](z)=-z+\cdots$ handles negative $m$; here $m\in\mathcal O^\times$ because $p\nmid m$. A power series $f(z)=mz+\cdots$ with
coefficients in $\mathcal O$ and $m$ a unit has a compositional inverse $g(z)=m^{-1}z+\cdots$ with coefficients in $\mathcal O$:
writing $g=\sum_{n\ge1}g_nz^n$, the coefficient of $z^n$ in $g(f(z))$ is $g_nm^n$ plus a polynomial with integer coefficients in
$g_1,\dots,g_{n-1}$ and the coefficients of $f$, so the equations $g(f(z))=z$ determine the $g_n$ recursively, dividing only by
powers of $m$. As $f$ and $g$ have coefficients in $\mathcal O$ and no constant term, both converge on $\mathfrak m$ and map
$\mathfrak m$ into $\mathfrak m$, and the identities $g\circ f=f\circ g=z$ of power series become identities of maps. So $[m]$ is a
bijection of $\hat E(\mathfrak m)$, hence of $E_1(K)$; in particular its kernel $E_1(K)[m]$ is $0$.
\end{proof}

\begin{corollary}[torsion and reduction]\label{cor:reduction:torsion-injects}
Let $m$ be prime to $p$.
\begin{enumerate}[label=(\roman*)]
\item Reduction $E_0(K)[m]\to\tilde E_{\mathrm{ns}}(k)$ is injective.
\item If $E$ has good reduction, $E(K)[m]\into\tilde E(k)$; the inertia group $I_K$ acts trivially on $E[m]$, so
$E[m]\subset E(K^{\ur})$; and reduction is an isomorphism $E[m]\to\tilde E[m](\bar k)$ compatible with the actions
of $G_K$ and $\Gal(\bar k/k)$. Consequently $I_K$ acts trivially on $T_\ell E$ for every prime $\ell\ne p$.
\end{enumerate}
\end{corollary}

\begin{proof}
(i) The kernel is $E_1(K)[m]$, which is $0$ by \Cref{thm:reduction:formal-group}(iii).
(ii) The first statement is (i) with $E_0(K)=E(K)$. Let $L=K(E[m])$, a finite extension; the given model has unit
discriminant over $L$ too, so $E/L$ has good reduction and (i) over $L$ shows that $E(L)[m]=E[m]$ injects into
$\tilde E(k_L)$. Let $\sigma\in I_K$ and $P\in E[m]$. Since $\sigma$ acts trivially on $k_L$, the point $\sigma P-P\in E[m]$
reduces to $\widetilde{\sigma P}-\tilde P=\tilde\sigma\tilde P-\tilde P=\tilde O$, so $\sigma P=P$. Hence $E[m]\subset E(K^{\ur})$.
Reduction from $E[m]$ to $\tilde E[m](\bar k)$ is injective, and both groups have $m^2$ elements ($m$ is prime to the
characteristics; \Cref{ch:ellcurves}), so it is bijective. It commutes with the Galois actions because reduction
commutes with automorphisms of $K^{\ur}$. Passing to the limit over $m=\ell^n$ gives the statement for $T_\ell E$.
\end{proof}

\begin{proposition}[the index of $E_0$]\label{prop:reduction:index-finite}
$E(K)/E_0(K)$ is finite when $k$ is finite; it is trivial when the reduction is good, and its order (the Tamagawa
number) is computed in \Cref{prop:reduction:component-group}, \Cref{thm:reduction:kodaira-neron} and
\Cref{prop:reduction:tate-tamagawa}.
\end{proposition}

\begin{proof}
For good reduction $E_0(K)=E(K)$ by definition. In general $E_0(K)$ is open in the compact group $E(K)$ (reduction is locally constant on $E(K)$: two points whose primitive
coordinates are congruent modulo $\mathfrak m$ have the same reduction), so its cosets form an open cover of the compact group $E(K)$, and the
index is finite. Compactness: as $k$ is finite, $\mathcal O=\varprojlim\mathcal O/\mathfrak m^n$ is compact, so the set of primitive triples
(those in $\mathcal O^3$ not all in $\mathfrak m$) is compact; it maps continuously onto $\mathbb P^2(K)$, which is therefore compact, and
$E(K)$ is closed in it, being the zero set of a homogeneous polynomial. The finer
statement that $E(K)/E_0(K)$ is the group of $k$-points of the component group of the N\'eron model is
\Cref{prop:reduction:component-group}; see \cite[VII.6]{SilvermanAEC} and \cite[IV.9]{SilvermanATAEC}.
\end{proof}

\paragraph{An example: torsion outside $E_0$.} The curve \curve{11a1}, $y^2+y=x^3-x^2-10x-20$, has split multiplicative
reduction of type $\kod{I}_5$ at $11$ with $c_{11}=5$ (\Cref{ex:reduction:running-table}), and $E(\Q)_{\mathrm{tors}}\simeq\Z/5$, with
nonzero points $(5,5)$, $(16,-61)$, $(16,60)$, $(5,-6)$ (the multiples $P,2P,3P,4P$ of $P=(5,5)$, computed with PARI). Modulo $11$
all four reduce to $(5,5)$, since $16\equiv5$, $-61\equiv60\equiv-6\equiv5$. With $F=y^2+y-x^3+x^2+10x+20$ one has $F(5,5)=0$,
$\partial F/\partial y=2y+1=11$ and $\partial F/\partial x=-3x^2+2x+10=-55$ at $(5,5)$; both partials vanish modulo $11$, so $(5,5)$ is
the node of $\tilde E$. Hence no nonzero torsion point lies in $E_0(\Q_{11})$, the subgroup $E(\Q)_{\mathrm{tors}}$ meets $E_0(\Q_{11})$
trivially, and it maps isomorphically onto $E(\Q_{11})/E_0(\Q_{11})$, a group of order $c_{11}=5$. This does not contradict
\Cref{cor:reduction:torsion-injects}(i), which constrains only torsion inside $E_0$.

\begin{paperlink}
Section 1.1 of the paper: $\delta_p=\dim V^{I_p}$ is $2$ at a good prime because $I_p$ acts trivially on $T_\ell E$
(\Cref{cor:reduction:torsion-injects}(ii)); the values $1$ and $0$ at multiplicative and additive primes come
from \Cref{prop:reduction:tate-torsion} and \Cref{prop:reduction:conductor-values}. Section 1.4: the Tamagawa
number is $c_q=[E(\Q_q):E_0(\Q_q)]$.
\end{paperlink}

\section{Potentially good and potentially multiplicative reduction}\label{sec:reduction:potential}

The reduction type can change under finite extensions of $K$, and the paper's local terms depend on which
reduction a curve acquires: the term $t_\ell(E,q)$ of Definition def:terms lives at the potentially
multiplicative primes, and the term $e_\ell(E,q)$ at a tame additive prime is read off the potentially good
types. This section proves the criterion by $v(j)$ and states the semistable reduction theorem.

\begin{definition}[potential reduction]\label{def:reduction:potential}
$E/K$ has \emph{potentially good} (resp.\ \emph{potentially multiplicative}) reduction if there is a finite
extension $L/K$ such that $E/L$ has good (resp.\ multiplicative) reduction.
\end{definition}

The definition makes sense because a minimal model over $K$ need not stay minimal over an extension: the change
$(u,0,0,0)$ with $u$ in the larger field can remove a factor of the discriminant that no $u\in K$ can remove. For
example, let $p$ be odd and $E\colon y^2=x^3+px$ over $\Q_p$. Then $c_4=-48p$ and $\Delta=-16\cdot4p^3=-64p^3$, so $v(\Delta)=3<12$,
the model is minimal (\Cref{prop:reduction:min-criterion}(i)) and the reduction is additive ($v(c_4)>0$); PARI gives type
$\kod{III}$ at $p=3$ and $p=5$. Over $L=\Q_p(p^{1/4})$ the change $(u,0,0,0)$ with $u=p^{1/4}$, that is $x=p^{1/2}x'$, $y=p^{3/4}y'$,
gives $p^{3/2}y'^2=p^{3/2}x'^3+p^{3/2}x'$, that is, $y'^2=x'^3+x'$, with discriminant $-64$, a unit: $E$ has good reduction over $L$.
Accordingly $j=1728$ is integral, $v(j)\ge0$, as the criterion below requires.

\begin{theorem}[the $j$-invariant criterion]\label{thm:reduction:potential-criterion}
Let $E/K$ be an elliptic curve.
\begin{enumerate}[label=(\roman*)]
\item $E$ has potentially good reduction if and only if $v(j)\ge0$.
\item $E$ has potentially multiplicative reduction if and only if $v(j)<0$; in that case $E$ has multiplicative
reduction over every finite extension over which it is semistable, and never good reduction.
\item If $E$ has multiplicative reduction, $v(j)=-v(\Delta_{\min})<0$.
\end{enumerate}
\end{theorem}

\begin{proof}[Sketch of proof]
(iii) For a minimal model $c_4$ is a unit, so $v(j)=3v(c_4)-v(\Delta)=-v(\Delta)<0$.

(i), the direction ``$\Rightarrow$'': if $E/L$ has good reduction, a minimal model over $L$ has $\Delta$ a unit and
$c_4$ integral, so $j=c_4^3/\Delta$ is integral in $L$, hence $v(j)\ge0$.

(i), the direction ``$\Leftarrow$'', for $p$ odd: completing the square (\Cref{ch:ellcurves}) gives a model $y^2=g(x)$ with $g$
a monic cubic with distinct roots $e_1,e_2,e_3$, and over the finite extension $L=K(E[2])=K(e_1,e_2,e_3)$ it reads
$y^2=(x-e_1)(x-e_2)(x-e_3)$. Over $L'=L(\sqrt{e_2-e_1})$ the change $x=e_1+(e_2-e_1)x'$, $y=(e_2-e_1)^{3/2}y'$ turns the three
factors into $(e_2-e_1)x'$, $(e_2-e_1)(x'-1)$ and $(e_2-e_1)(x'-\lambda)$ with $\lambda=(e_3-e_1)/(e_2-e_1)\ne0,1$; dividing by
$(e_2-e_1)^3$ gives the Legendre model $y'^2=x'(x'-1)(x'-\lambda)$, for which
\[
j=2^8\frac{(\lambda^2-\lambda+1)^3}{\lambda^2(\lambda-1)^2},\qquad\Delta=16\lambda^2(\lambda-1)^2 .
\]
(Here $v$ denotes the valuation of $L'$ extending that of $K$; the sign of $v(j)$ does not depend on the normalisation, and $v(2)=0$.)
If $v(\lambda)<0$ the numerator has valuation $6v(\lambda)$ and the denominator $4v(\lambda)$, so $v(j)=2v(\lambda)<0$. If
$v(\lambda)>0$ then $\lambda^2-\lambda+1$ and $\lambda-1$ are units, so $v(j)=-2v(\lambda)<0$, and likewise if $v(\lambda-1)>0$. So $v(j)\ge0$ forces $\lambda$ and $\lambda-1$ to be
units (if $v(\lambda)=0$ then $v(\lambda-1)\ge0$, and the remaining case is $v(\lambda-1)=0$), and then the Legendre model is
integral with unit discriminant ($p$ odd): good reduction over $L'$. For $p=2$
one uses a different cubic normal form; see \cite[VII.5.5]{SilvermanAEC}.

(ii) If $v(j)<0$, $E$ never has good reduction, by (i). By the semistable reduction theorem
(\Cref{thm:reduction:semistable}) it is semistable over some finite extension, hence multiplicative there, and
over every larger extension. Conversely multiplicative reduction over $L$ gives $v_L(j)<0$ by (iii), hence
$v(j)<0$.
\end{proof}

\begin{theorem}[semistable reduction]\label{thm:reduction:semistable}
For every elliptic curve $E/K$ there is a finite extension $L/K$ such that $E/L$ is semistable. One may take
$L=K(E[m])$ for any integer $m\ge3$ prime to $p$.
\end{theorem}

\begin{proof}[Sketch of proof]
If $v(j)\ge0$, \Cref{thm:reduction:potential-criterion}(i) gives good reduction over a finite extension. If
$v(j)<0$, $E$ becomes isomorphic to a Tate curve, which has split multiplicative reduction, over an extension
of degree at most $2$ (\Cref{thm:reduction:tate-uniformisation}). The statement about $K(E[m])$ is the theorem of
Serre and Tate \cite[Section 2]{SerreTate68}; see also \cite[VII.5]{SilvermanAEC} and
\Cref{thm:reduction:serre-tate} below.
\end{proof}

\begin{remark}[the four classes]\label{rem:reduction:types-of-the-four-cases}
Combining \Cref{def:reduction:types}, \Cref{thm:reduction:potential-criterion} and the classification of
\Cref{sec:reduction:kodaira}, the Kodaira symbols fall into four classes:
\begin{itemize}
\item good reduction: $\kod{I}_0$;
\item multiplicative reduction: $\kod{I}_n$, $n\ge1$, with $n=v(\Delta_{\min})=-v(j)$;
\item additive, potentially multiplicative reduction: $\kod{I}_n^*$ with $n\ge1$ (\Cref{thm:reduction:potmult-additive});
\item additive, potentially good reduction: $\kod{II}$, $\kod{III}$, $\kod{IV}$, $\kod{I}_0^*$, $\kod{IV}^*$, $\kod{III}^*$, $\kod{II}^*$,
and, at $p=2$ only, also some symbols $\kod{I}_m^*$ with $m\ge1$.
\end{itemize}
Until \Cref{sec:reduction:kodaira} the symbols are used only as names for these classes.
\end{remark}

\begin{warning}\label{warn:reduction:potmult}
``Potentially multiplicative'' does not mean ``multiplicative'': a curve of type $\kod{I}_n^*$ ($n\ge1$) at an
odd prime has additive reduction ($f_p\ge2$, $a_p=0$, $\delta_p=0$), but $v(j)<0$; at $p=2$ a symbol $\kod{I}_m^*$ may also be
potentially good (\Cref{rem:reduction:Instar-at-2}). The paper's $n_q=v_q(\Delta_E)$ is defined at
multiplicative primes; at an additive potentially multiplicative prime the relevant integer is $\nu=-v_q(j_E)$,
which for odd $q$ is the index $n$ of the symbol $\kod{I}_n^*$ (\Cref{thm:reduction:potmult-additive}) but for
$q=2$ is not.
\end{warning}

\begin{paperlink}
Definition def:terms: $t_\ell(E,q)=v_\ell(-v_q(j_E))$ if $v_q(j_E)<0$ and $0$ otherwise, ``the $\ell$-part of the
valuation of the Tate parameter at a potentially multiplicative prime''. Lemma lem:euler: ``the inertia group acts
on $V$ through a finite cyclic quotient when the reduction is potentially good, and through a quadratic character
times a unipotent when it is potentially multiplicative'' (\Cref{thm:reduction:serre-tate},
\Cref{thm:reduction:tate-uniformisation}). Remark rem:wild: a symbol $\kod{I}_n^*$ at $2$ can be potentially good
(\Cref{rem:reduction:Instar-at-2}).
\end{paperlink}

\section{N\'eron models and component groups}\label{sec:reduction:neron}

The paper's inequalities (eq:main) and (eq:twist) are stated in terms of the geometric component group
$\Phig_q$ of the N\'eron model of $E$ at $q$, not the Tamagawa number $c_q$. This section defines both, explains
why they differ, and proves the fact that links component groups to Galois representations:
the inertia invariants of $E[\ell^t]$ at an additive prime are the $\ell^t$-torsion of $\Phig_q$
(\Cref{prop:reduction:component-invariants}). That proposition is the decisive input of Lemma lem:IV and Remark
rem:types (via \Cref{ch:galrep}), and of the treatment of the types $\kod{IV}$, $\kod{IV}^*$ at $2$ in
\Cref{cor:reduction:IV-at-2}. We use scheme-theoretic language only as far as the statements need it.

A minimal Weierstrass model gives a scheme over $\mathcal O$ whose generic fibre is $E$ and whose special fibre is
the cubic $\tilde E$. When the reduction is bad this model has two defects: its special fibre is singular, and
its points over $\mathcal O$ do not see the group structure of $E(K)$ well (only $E_0(K)$ reduces into the group
$\tilde E_{\mathrm{ns}}(k)$). N\'eron's idea is to look for the best smooth model: one that is a group scheme and
whose $\mathcal O$-points are all of $E(K)$. Its special fibre is a smooth commutative algebraic group over $k$,
usually disconnected; the identity component is $\tilde E_{\mathrm{ns}}$, and the group of components is a finite
abelian group with an action of $\Gal(\bar k/k)$.

\begin{definition}[N\'eron model]\label{def:reduction:neron}
Let $R=\mathcal O$. A \emph{N\'eron model} of $E/K$ is a smooth separated group scheme $\mathcal E$ of finite type over
$R$ with generic fibre $\mathcal E_K=E$, such that for every smooth $R$-scheme $S$ every $K$-morphism $S_K\to E$
extends uniquely to an $R$-morphism $S\to\mathcal E$ (the \emph{N\'eron mapping property}).
\end{definition}

\begin{theorem}[N\'eron]\label{thm:reduction:neron-existence}
Every elliptic curve $E/K$ has a N\'eron model $\mathcal E$, unique up to unique isomorphism. It is the smooth
locus of the minimal regular proper model of $E$ over $R$; if $E$ has good reduction it is the minimal
Weierstrass model.
\end{theorem}

We do not prove this; see N\'eron \cite{NeronModels}, \cite[Sections 1.2, 1.4]{BLR} and
\cite[IV.5, IV.6]{SilvermanATAEC}.

\begin{proposition}[the component group]\label{prop:reduction:component-group}
Let $\mathcal E$ be the N\'eron model of $E/K$, $\mathcal E_k$ its special fibre and $\mathcal E_k^0$ the connected
component of the identity.
\begin{enumerate}[label=(\roman*)]
\item $\mathcal E(R)=E(K)$. For a minimal Weierstrass model, $\mathcal E_k^0\simeq\tilde E_{\mathrm{ns}}$ and
$\mathcal E^0(R)=E_0(K)$, where $\mathcal E^0$ is the open subgroup scheme with special fibre $\mathcal E_k^0$. The quotient
$\mathcal E_k/\mathcal E_k^0$ is a finite \'etale group scheme over $k$, the \emph{component group scheme} of $E$; we write its
groups of points with an argument, $\Phi_E(k)$ and $\Phi_E(\bar k)$, never as a bare letter. Concretely it is a finite abelian
group $\Phi_E(\bar k)$ with a continuous action of $\Gal(\bar k/k)$, and $\Phi_E(k)=\Phi_E(\bar k)^{\Gal(\bar k/k)}$.
\item If $k$ is finite, $E(K)/E_0(K)\simeq\Phi_E(k)$.
\end{enumerate}
\end{proposition}

\begin{proof}[Sketch of proof]
(i) is \cite[IV.5, IV.9]{SilvermanATAEC} and \cite[Section 7.1]{BLR}; the equality $\mathcal E(R)=E(K)$ is the mapping
property applied to $S=\Spec R$. (ii) Since $\mathcal E$ is smooth, Hensel's lemma shows that reduction
$\mathcal E(R)\to\mathcal E_k(k)$ is surjective. Since $k$ is finite and $\mathcal E_k^0$ is connected, Lang's theorem applies: a connected algebraic group $G$ over a
finite field $k$ has a $k$-point on every principal homogeneous space (a variety over $k$ on which $G$ acts simply
transitively over $\bar k$); in the language of Galois cohomology (\Cref{ch:cohomology}), $H^1(k,G)=0$. A component
defined over $k$, that is a point of $\Phi_E(k)$, is such a space for $G=\mathcal E_k^0$, so it has a $k$-point, and
$\mathcal E_k(k)\to\Phi_E(k)$ is surjective. The kernel of the composite
$E(K)=\mathcal E(R)\to\Phi_E(k)$ is the set of points reducing into the identity component, $\mathcal E^0(R)=E_0(K)$. See
\cite[IV.9]{SilvermanATAEC}.
\end{proof}

\begin{definition}[geometric component group and Tamagawa number]\label{def:reduction:tamagawa}
Let $E/\Q$ and let $p$ be a prime. Write $\Phi_p(\F_p)$ and $\Phi_p(\bar\F_p)$ for the groups of points of the component group
scheme over $\F_p$ of the N\'eron model of $E$ over $\Z_p$ (\Cref{prop:reduction:component-group}). The
\emph{geometric component group} is $\Phig_p=\Phi_p(\bar\F_p)$, and the \emph{Tamagawa number} is
\[
c_p=\#\Phi_p(\F_p)=[E(\Q_p):E_0(\Q_p)].
\]
At a good prime $\Phig_p=0$ and $c_p=1$. The product $\prod_pc_p$ is the Tamagawa factor in the conjecture of
Birch and Swinnerton-Dyer; Cremona's file \texttt{allbsd} lists it (\Cref{ch:computation}).
\end{definition}

\begin{convention}[$\Phig$ versus $c$]\label{conv:reduction:phig}
The component group is always written with an argument: $\Phi_p(\F_p)$ or $\Phig_p=\Phi_p(\bar\F_p)$ (over a general local
field, $\Phi_E(k)$ or $\Phig=\Phi_E(\bar k)$); a bare $\Phi$ is the geometric Frobenius (Appendix B). Since $\Phi_p(\F_p)$
is the subgroup of $\Gal(\bar\F_p/\F_p)$-invariants of $\Phig_p$, $c_p$ divides $\#\Phig_p$, with equality if and only
if Frobenius acts trivially on $\Phig_p$. The paper's statements are about $\#\Phig_q$, never about $c_q$: for
instance \curve{121c1} and \curve{121c2} at $11$ and \curve{980i1} at $2$ have $\#\Phig=3$ but $c=1$
(\Cref{ex:reduction:running-table}). We never write $c_p$ for $\#\Phig_p$, and we never use a bare $c$ for the class of
the unit part of the Tate parameter, which is $c(u)$ (\Cref{def:reduction:tate-parameter}).
\end{convention}

\begin{proposition}[unramified base change]\label{prop:reduction:etale-base-change}
The formation of the N\'eron model commutes with \'etale base change; in particular with base change from $K$
to a finite unramified extension $K'$, or to the completion of $K^{\ur}$. Consequently the reduction type, the
Kodaira symbol, the group $\Phig$ with its structure, $v(\Delta_{\min})$ and the conductor exponent are
unchanged under unramified extension of $K$, and a minimal model over $K$ stays minimal over $K'$. The Tamagawa
number can change, since only the invariants of $\Gal(\bar k/k')$ on $\Phig$ are counted.
\end{proposition}

We do not prove this; see \cite[Section 7.2]{BLR} and \cite[IV.6]{SilvermanATAEC}. (For the conductor
exponent the invariance is also immediate from \Cref{def:reduction:conductor-exponent}, since $I_{K'}=I_K$.)

\begin{remark}[Tamagawa numbers at multiplicative primes]\label{rem:reduction:tamagawa-mult}
At a prime of type $\kod{I}_n$ one has $\Phig\simeq\Z/n$, and $c=n$ for split, $c\in\{1,2\}$ (according to the parity
of $n$) for nonsplit reduction (\Cref{prop:reduction:tate-tamagawa}). So for $\ell$ odd the $\ell$-part of $c$ is
$\ell^{v_\ell(n)}$ in the split case and $1$ in the nonsplit case, while the $\ell$-part of $\#\Phig$ is
$\ell^{v_\ell(n)}$ in both. This is why the paper's inequality (eq:main) uses $\Phig_q$, and why the history paragraph
of its Section 1.5 (Ribet and Takahashi) speaks of the $n_q$.

The facts $\Phig\simeq\Z/n$ and the values of $c$ at $\kod{I}_n$ are proved in \Cref{sec:reduction:tate}
(\Cref{prop:reduction:tate-tamagawa}), from the Tate curve. That proof uses only the uniformisation of
\Cref{thm:reduction:tate-uniformisation} and \Cref{prop:reduction:etale-base-change}, not the proposition below, so the
forward references to it in this remark and in the proof of \Cref{prop:reduction:component-invariants}(ii) are not
circular.
\end{remark}

\begin{proposition}[inertia invariants and component groups]\label{prop:reduction:component-invariants}
Let $K/\Q_q$ be finite, with residue field $k$, let $\ell\ne q$ be a prime, $t\ge1$ an integer, and let $E/K$ have bad
reduction, with geometric component group $\Phig=\Phi_E(\bar k)$.
\begin{enumerate}[label=(\roman*)]
\item If the reduction is additive, there is a $\Gal(\bar k/k)$-equivariant isomorphism
\[
E[\ell^t]^{I_K}=E(K^{\ur})[\ell^t]\ \simeq\ \Phig[\ell^t].
\]
In particular $E(K)[\ell^\infty]\simeq\Phi_E(k)[\ell^\infty]$.
\item If the reduction is multiplicative, there is an exact sequence of $\Gal(\bar k/k)$-modules
\[
0\longrightarrow\mu_{\ell^t}\otimes\varepsilon\longrightarrow E[\ell^t]^{I_K}\longrightarrow\Phig[\ell^t]\longrightarrow0,
\]
where $\Phig\simeq\Z/n$ ($n=v(\Delta_{\min})$) and $\varepsilon$ is the unramified character of order at most $2$
which is trivial if and only if the reduction is split.
\item Hence, for $\ell$ odd, $\dim_{\F_\ell}E[\ell]^{I_K}$ is $1+[\ell\mid n]$ for type $\kod{I}_n$ ($n\ge1$), $1$ for
types $\kod{IV}$ and $\kod{IV}^*$ when $\ell=3$, and $0$ for every other additive type.
\end{enumerate}
\end{proposition}

\begin{proof}
An element of $E[\ell^t]\subset E(\bar K)$ is fixed by $I_K$ if and only if its coordinates lie in the fixed field
$\bar K^{I_K}=K^{\ur}$; this is the first equality in (i). Let $K'$ run over the finite unramified extensions of $K$,
with residue fields $k'$. By \Cref{prop:reduction:etale-base-change} a minimal model over $K$ stays minimal over
each $K'$, the component group scheme of $E/K'$ is the base change of that of $E/K$ (so $\Phi_E(k')$ is the group of
$k'$-points of the same scheme and $\Phig$ is unchanged), and by \Cref{prop:reduction:filtration} and
\Cref{prop:reduction:component-group}(ii) there are exact sequences
\[
0\to E_1(K')\to E_0(K')\to\tilde E_{\mathrm{ns}}(k')\to0,\qquad 0\to E_0(K')\to E(K')\to\Phi_E(k')\to0,
\]
compatible with the inclusions $K'\subset K''$ and with the action of $\Gal(K^{\ur}/K)=\Gal(\bar k/k)$. Taking the union
over $K'$ (a direct limit of exact sequences is exact) gives
\[
0\to E_1^{\ur}\to E_0^{\ur}\to\tilde E_{\mathrm{ns}}(\bar k)\to0,\qquad 0\to E_0^{\ur}\to E(K^{\ur})\to\Phig\to0,
\]
with $E_i^{\ur}=\bigcup_{K'}E_i(K')$. Now:
\begin{itemize}
\item $E_1^{\ur}$ is uniquely $\ell$-divisible, by \Cref{thm:reduction:formal-group}(iii) applied over each $K'$.
\item In the additive case $\tilde E_{\mathrm{ns}}(\bar k)\simeq\bar k^+$ (\Cref{prop:reduction:Ens}), which is uniquely
$\ell$-divisible as $\ell\ne q$. In the multiplicative case $\tilde E_{\mathrm{ns}}(\bar k)\simeq\bar k^\times$, which is
$\ell$-divisible ($\bar k$ is algebraically closed) with $\ell^t$-torsion $\mu_{\ell^t}(\bar k)\simeq\Z/\ell^t$. The
Galois action is the natural one if the reduction is split, and is twisted by the nontrivial unramified quadratic
character $\varepsilon$ if it is nonsplit, by the formula $\phi(\sigma P)=\sigma(\phi(P))^{-1}$ in the proof of
\Cref{prop:reduction:Ens}(iii); reduction identifies $\mu_{\ell^t}(\bar k)$ with $\mu_{\ell^t}\subset K^{\ur}$.
\end{itemize}
Apply the snake lemma (\Cref{ch:modules}) to multiplication by $\ell^t$ on the first sequence. Since $E_1^{\ur}$ has
no $\ell^t$-torsion and $E_1^{\ur}/\ell^t=0$, we get $E_0^{\ur}[\ell^t]\simeq\tilde E_{\mathrm{ns}}(\bar k)[\ell^t]$ and
$E_0^{\ur}/\ell^tE_0^{\ur}\simeq\tilde E_{\mathrm{ns}}(\bar k)/\ell^t=0$. Then the snake lemma for $\ell^t$ on the second
sequence gives
\[
0\to E_0^{\ur}[\ell^t]\to E(K^{\ur})[\ell^t]\to\Phig[\ell^t]\to E_0^{\ur}/\ell^tE_0^{\ur}=0 .
\]
In the additive case $E_0^{\ur}[\ell^t]=0$, which is (i); in the multiplicative case it is $\mu_{\ell^t}\otimes\varepsilon$,
which is (ii), with $\Phig\simeq\Z/n$ by \Cref{prop:reduction:tate-tamagawa}. Every map is
$\Gal(\bar k/k)$-equivariant. Taking $\Gal(\bar k/k)$-invariants in (i) gives
$E(K)[\ell^t]=(E[\ell^t]^{I_K})^{G_K/I_K}\simeq\Phig[\ell^t]^{\Gal(\bar k/k)}=\Phi_E(k)[\ell^t]$.

(iii) In case (ii), $\#E[\ell]^{I_K}=\ell\cdot\#(\Z/n)[\ell]$, which is $\ell^2$ if $\ell\mid n$ and $\ell$ otherwise. In case
(i), $\Phig$ is one of $0$, $\Z/2$, $\Z/3$, $(\Z/2)^2$, $\Z/4$ (\Cref{thm:reduction:kodaira-neron}), and for $\ell$ odd its
$\ell$-torsion is nonzero only when $\ell=3$ and $\Phig\simeq\Z/3$, that is, for the types $\kod{IV}$ and $\kod{IV}^*$, where
it is a line.
\end{proof}

\begin{paperlink}
The paragraph before Lemma lem:IV in the revised paper states (i) with the proof above (``the snake lemma for
multiplication by $\ell^k$ gives the isomorphism'', citing \cite[Section 1]{SerreTate68}). Lemma lem:IV: ``$E[3]^{I_q}$
is a line'' is (iii) for $\kod{IV}$, $\kod{IV}^*$. Remark rem:types: ``$E[3]^{I_q}=0$'' for the other additive types is
(iii). Lemma lem:tamfactor, the final check with $M=T$: the paper's module $(\Z/\ell^{v_\ell(n)})\bar e_1$ (the torsion of
$H^1(I,T)$, generated by the class of $e_1$), with Frobenius acting by $\varepsilon$, is isomorphic as a Galois module to
$\Phig[\ell^\infty]$, the quotient in (ii) generated by the image of $e_2$ (Frobenius acts on $\Phig\simeq\Z/n$ by $\varepsilon$,
\Cref{prop:reduction:tate-tamagawa}); so its Frobenius invariants have the order of $\Phi_E(k)[\ell^\infty]$, the $\ell$-part of
$c_q(E)$ (\Cref{rem:reduction:tamagawa-mult}). They are not the invariants of the sub $\mu_{\ell^t}\otimes\varepsilon$, on which
Frobenius acts by $q\varepsilon$. Chapter 11 (\Cref{ch:galrep}, its item 2.2) cites this proposition instead of
proving it again.
\end{paperlink}

\section{Kodaira symbols and Tate's algorithm}\label{sec:reduction:kodaira}

The paper names the reduction of $E$ at a bad prime by its Kodaira symbol, reads $\#\Phig_q$ off the symbol
(Section 1.4), and for a tame prime $q\ge5$ reads the omitted Euler factor off the symbol (the table of Section 1,
Lemma lem:euler). It computes the symbols with Tate's algorithm, through PARI (Section 2, Section sec:data, the
appendix). This section states the classification of Kodaira and N\'eron, states Tate's algorithm in full, derives
from it the table for $p\ge5$ and the criterion for $\kod{I}_0^*$ that Lemma lem:twist (ii) uses, fixes PARI's codes,
and runs the algorithm by hand on the paper's curves.

\paragraph{Regular models, blowing up, multiplicities.} The classification below is a statement about a surface over
$\mathcal O$, and we explain the words it uses for a reader without algebraic geometry. A \emph{model} of $E$ over $\mathcal O$ is a
two-dimensional scheme, proper and flat over $\mathcal O$, with generic fibre $E$; the minimal Weierstrass model is one, given by
one equation $F(x,y)=0$ in the coordinates $x,y$ over $\mathcal O$ (plus the point at infinity). Such a surface is
\emph{regular} at a point $P$ of its special fibre if the local ring there has a maximal ideal generated by two elements. For
a surface given by one equation, and $P=(x_0,y_0)$ in the special fibre, this holds if and only if $F$ does not lie in
$\mathfrak n_P^2$, where $\mathfrak n_P=(\varpi,x-\tilde x_0,y-\tilde y_0)$ with lifts $\tilde x_0,\tilde y_0$; informally, the three
partial ``derivatives'' of $F$ in the directions $\varpi$, $x$, $y$ do not all vanish at $P$. A Weierstrass model is regular at
every smooth point of $\tilde E$, and may fail to be regular at the singular point. \emph{Blowing up} a point $P$ replaces it by
the projective line of directions through it: in coordinates, one substitutes $x-\tilde x_0=\varpi x_1$, $y-\tilde y_0=\varpi y_1$
(and similarly in the other charts, where $x-\tilde x_0$ or $y-\tilde y_0$ is the divisor) and removes the factor of $\varpi$ that
appears. Repeating this at the non-regular points, and then contracting the components that can be contracted without
losing regularity, gives the \emph{minimal regular proper model}, unique up to isomorphism (Lichtenbaum, Shafarevich; see
\cite[IV.7]{SilvermanATAEC}). Its special fibre is the divisor of $\varpi$, a union of curves $C_i$ counted with
\emph{multiplicities} $m_i\ge1$: $m_i$ is the order of vanishing of $\varpi$ along $C_i$. The \emph{dual graph} has a vertex for each
$C_i$ and an edge for each intersection point. A point of $E(K)$ extends, by properness, to an $\mathcal O$-point of the regular model, which meets the special fibre with
intersection number $1$; so it reduces to a smooth point of a component of multiplicity $1$; this is why only those components
contribute to the N\'eron model and to $\Phig$.

\begin{example}[blowing up a node: type $\kod{I}_2$]\label{ex:reduction:blowup}
Let $p\ge5$ and $E\colon y^2=x^3+x^2+p^2$ over $\Q_p$. Here $b_2=4$, $b_4=0$, $b_6=4p^2$, $b_8=4p^2$, so $c_4=16$ is a unit and
$\Delta=-16p^2(4+27p^2)$ has $v_p(\Delta)=2$: multiplicative reduction, with $\tilde E\colon y^2=x^3+x^2$ nodal at $(0,0)$, with tangents
$y=\pm x$ (split). The equation $F=y^2-x^2-x^3-p^2$ lies in $(p,x,y)^2$, so the model is not regular at the node. Blow it up:
with $x=px_1$, $y=py_1$ one gets $F=p^2(y_1^2-x_1^2-px_1^3-1)$, and the new surface $y_1^2-x_1^2-1-px_1^3=0$ is regular along the
special fibre of this chart (the partial derivatives $2y_1$ and $-2x_1$ do not both vanish on $y_1^2-x_1^2=1$; the other charts
are checked in the same way). Its special
fibre in this chart is the conic $y_1^2-x_1^2=1$, a new component $C_1\simeq\mathbb P^1$ along which $\varpi=p$ vanishes to order $1$.
The strict transform $C_0$ of $\tilde E$ is the nodal cubic with its node pulled apart into the two points of $C_1$ at infinity in
the directions $y_1=\pm x_1$, which are the tangent directions of the node. So the special fibre is two rational curves of
multiplicity $1$ meeting in two points: the cycle $\tilde A_1$ of type $\kod{I}_2$. Removing the two intersection points leaves two
copies of $\mathbb G_m$, and $\Phig\simeq\Z/2$, as in \Cref{tab:reduction:kodaira}. In the additive types higher multiplicities
appear: in type $\kod{I}_0^*$ the central component has multiplicity $2$, so it contributes nothing to the N\'eron model, and
$\Phig$ comes from the four outer components.
\end{example}

\begin{theorem}[Kodaira, N\'eron]\label{thm:reduction:kodaira-neron}
Let $E/K$ be an elliptic curve and let $m$ be the number of irreducible components (over $\bar k$) of the special
fibre of the minimal regular proper model of $E$ over $\mathcal O$. The special fibre over $\bar k$ is one of the
configurations of \Cref{tab:reduction:kodaira}. The N\'eron model's special fibre is the union of the
multiplicity-one components minus their intersection points with the other components, and $\Phig$ is
the group of multiplicity-one components, with the structure given in the table. Over a finite residue field
$k$ the Tamagawa number takes only the values in the last column.
\end{theorem}

\begin{table}[ht]
\centering\footnotesize\setlength{\tabcolsep}{4pt}
\begin{tabular}{p{1.5cm}lp{4.6cm}p{1.6cm}lp{2.3cm}}
\toprule
symbol & $m$ & special fibre over $\bar k$ & $\Phig$ & $\#\Phig$ & possible $c$\\
\midrule
$\kod{I}_0$ & $1$ & smooth curve of genus $1$ & $0$ & $1$ & $1$\\
$\kod{I}_1$ & $1$ & rational curve with a node & $0$ & $1$ & $1$\\
$\kod{I}_n$, $n\ge2$ & $n$ & $n$ smooth rational curves of multiplicity $1$ forming a cycle ($\tilde A_{n-1}$) & $\Z/n$ & $n$ & $n$ (split); $1$ or $2$ for $n$ odd or even (nonsplit)\\
$\kod{II}$ & $1$ & rational curve with a cusp & $0$ & $1$ & $1$\\
$\kod{III}$ & $2$ & two rational curves meeting at one point with a common tangent & $\Z/2$ & $2$ & $2$\\
$\kod{IV}$ & $3$ & three rational curves through one point, with distinct tangents & $\Z/3$ & $3$ & $1$ or $3$\\
$\kod{I}_0^*$ & $5$ & one curve of multiplicity $2$ meeting four of multiplicity $1$ ($\tilde D_4$) & $(\Z/2)^2$ & $4$ & $1$, $2$ or $4$\\
$\kod{I}_n^*$, $n\ge1$ & $n+5$ & a chain of $n+1$ curves of multiplicity $2$, with two curves of multiplicity $1$ at each end ($\tilde D_{n+4}$) & $\Z/4$ ($n$ odd), $(\Z/2)^2$ ($n$ even) & $4$ & $2$ or $4$\\
$\kod{IV}^*$ & $7$ & $\tilde E_6$, multiplicities $1,2,3,2,1,2,1$ & $\Z/3$ & $3$ & $1$ or $3$\\
$\kod{III}^*$ & $8$ & $\tilde E_7$, multiplicities $1,2,3,4,3,2,1,2$ & $\Z/2$ & $2$ & $2$\\
$\kod{II}^*$ & $9$ & $\tilde E_8$, multiplicities $1,2,3,4,5,6,4,2,3$ & $0$ & $1$ & $1$\\
\bottomrule
\end{tabular}
\caption{The Kodaira-N\'eron classification. The components are smooth rational curves except in the first, second and
fourth rows; $\tilde A$, $\tilde D$, $\tilde E$ are the extended Dynkin diagrams formed by the dual graph (a vertex for each
component, an edge for each intersection point).}\label{tab:reduction:kodaira}
\end{table}

We do not prove \Cref{thm:reduction:kodaira-neron}; see Kodaira \cite{Kodaira}, N\'eron \cite{NeronModels}, Tate
\cite{Tate75} and \cite[IV.8, IV.9 and Table 4.1]{SilvermanATAEC}. The orders $\#\Phig$ are the ones quoted in Section
1.4 of the paper ($n$; $1$; $2$; $3$; $4$ for $\kod{I}_n$; $\kod{II},\kod{II}^*$; $\kod{III},\kod{III}^*$;
$\kod{IV},\kod{IV}^*$; $\kod{I}_0^*,\kod{I}_n^*$). In the language of Lie theory (a physics reader's aside): $\Phig$
is the centre of the simply connected compact Lie group of type $A_{n-1},D_{n+4},E_6,E_7,E_8$ for $\kod{I}_n,\kod{I}_n^*,\kod{IV}^*,\kod{III}^*,\kod{II}^*$
($\kod{III}$ and $\kod{IV}$ correspond to $A_1$ and $A_2$; $\kod{I}_0$, $\kod{I}_1$ and $\kod{II}$ have trivial $\Phig$), which is the dual
of the root lattice modulo the root lattice; the same diagrams classify the finite subgroups of $\mathrm{SU}(2)$, and
the binary tetrahedral group $\SL_2(\F_3)$ met in \Cref{thm:reduction:kraus} is one of them.

On Cremona's tables below conductor $10{,}000$ ($211{,}455$ pairs of a curve and a bad prime) every value of $c$ in
the last column occurs, and no other: for $\kod{IV},\kod{IV}^*$ one finds $c\in\{1,3\}$, for $\kod{I}_0^*$
$c\in\{1,2,4\}$, for $\kod{I}_n^*$ $c\in\{2,4\}$, and for $\kod{I}_n$ the values predicted by
\Cref{prop:reduction:tate-tamagawa} in all $153{,}343$ multiplicative pairs.

\begin{theorem}[Tate's algorithm]\label{thm:reduction:tate-algorithm}
Let an integral model of $E/K$ be given, with $\varpi$, $k$ as in \Cref{conv:reduction:local}, and write
$a_{i,j}=a_i/\varpi^j$. The following procedure terminates and returns the Kodaira symbol, the number $m$ of
components, the conductor exponent $f=v(\Delta)+1-m$ (\Cref{thm:reduction:ogg}), the Tamagawa number $c$ and a
minimal model (the model at the step where it stops).
\begin{enumerate}[label=\textbf{Step \arabic*.},leftmargin=*]
\item If $\varpi\nmid\Delta$: type $\kod{I}_0$, $m=1$, $f=0$, $c=1$.
\item Change coordinates by $(1,r,0,t)$, $r,t\in\mathcal O$, so that the singular point of $\tilde E$ is $(0,0)$; then
$\varpi\mid a_3,a_4,a_6$. If $\varpi\nmid b_2$: type $\kod{I}_n$ with $n=v(\Delta)$, $m=n$, $f=1$; the reduction is split
if $T^2+a_1T-a_2$ has a root in $k$, and then $c=n$; otherwise $c=1$ ($n$ odd) or $2$ ($n$ even).
\item If $\varpi^2\nmid a_6$: type $\kod{II}$, $m=1$, $f=v(\Delta)$, $c=1$.
\item If $\varpi^3\nmid b_8$: type $\kod{III}$, $m=2$, $f=v(\Delta)-1$, $c=2$.
\item If $\varpi^3\nmid b_6$: type $\kod{IV}$, $m=3$, $f=v(\Delta)-2$; $c=3$ if $Y^2+a_{3,1}Y-a_{6,2}$ has a root in $k$,
else $c=1$.
\item Change coordinates by $(1,0,s,t)$ so that $\varpi\mid a_1,a_2$, $\varpi^2\mid a_3,a_4$, $\varpi^3\mid a_6$. Let
$P(T)=T^3+a_{2,1}T^2+a_{4,2}T+a_{6,3}$. If $\tilde P$ has distinct roots: type $\kod{I}_0^*$, $m=5$,
$f=v(\Delta)-4$, $c=1+\#\{\text{roots of }\tilde P\text{ in }k\}$.
\item If $\tilde P$ has one simple and one double root: type $\kod{I}_n^*$ for some $n\ge1$, $m=n+5$,
$f=v(\Delta)-4-n$. To find $n$, translate the double root to $0$ (a change $(1,\varpi r_0,0,0)$); then
$\varpi\,\|\,a_2$, $\varpi^3\mid a_4$, $\varpi^4\mid a_6$. For $s=1,2,\dots$ alternately consider the quadratics
\[
Y^2+a_{3,(s+3)/2}Y-a_{6,s+3}\ (s\text{ odd}),\qquad a_{2,1}X^2+a_{4,(s+4)/2}X+a_{6,s+3}\ (s\text{ even});
\]
if the current one has distinct roots in $\bar k$, stop with $n=s$, and $c=4$ if its roots lie in $k$, else $c=2$;
if it has a double root, translate it to $0$ (by $t$ for $s$ odd, by $r$ for $s$ even, with the appropriate power
of $\varpi$) and continue with $s+1$.
\item If $\tilde P$ has a triple root, translate it to $0$; then $\varpi^2\mid a_2$, $\varpi^3\mid a_4$, $\varpi^4\mid a_6$. If
$Y^2+a_{3,2}Y-a_{6,4}$ has distinct roots: type $\kod{IV}^*$, $m=7$, $f=v(\Delta)-6$, $c=3$ if the roots are in $k$,
else $c=1$.
\item Otherwise translate its double root to $0$, so that $\varpi^3\mid a_3$, $\varpi^5\mid a_6$. If $\varpi^4\nmid a_4$: type
$\kod{III}^*$, $m=8$, $f=v(\Delta)-7$, $c=2$.
\item If $\varpi^6\nmid a_6$: type $\kod{II}^*$, $m=9$, $f=v(\Delta)-8$, $c=1$.
\item Otherwise the model is not minimal: replace $a_i$ by $a_i/\varpi^i$ (the change $(\varpi,0,0,0)$) and return to
Step 1.
\end{enumerate}
The translations can be made explicit. For $K=\Q_p$, Step 2 is achieved by the following $(r,t)$ modulo $p$:
\begin{itemize}
\item $p\ge5$: $r=-b_2/12$ if $p\mid c_4$, $r=-(c_6+b_2c_4)/(12c_4)$ otherwise, and $t=-(a_1r+a_3)/2$;
\item $p=3$: $r=-b_6$ if $3\mid b_2$, $r=-b_2b_4$ otherwise, and $t=a_1r+a_3$;
\item $p=2$: $r=a_4$, $t=r(1+a_2+a_4)+a_6$ if $2\mid b_2$, and $r=a_3$, $t=r+a_4$ otherwise.
\end{itemize}
For $p$ odd, Step 6 is completing the square: $s=-a_1/2$, $t=-a_3/2$, after which $a_1=a_3=0$; at $p=2$ one solves the
congruences for $s$ and $t$ directly.
\end{theorem}

\begin{proof}[Sketch of proof]
Each step either recognises a configuration of the special fibre of the minimal regular model, by blowing up the
singular point of the current model and examining the exceptional curves, or performs the change of coordinates
that puts the next singularity at the origin. The conditions $\varpi^i\mid a_j$ record which monomials vanish to which
order at the singular point. That the procedure stops and returns a minimal model is part of the theorem; the values
of $f$ come from Ogg's formula. See Tate \cite{Tate75}, \cite[IV.9]{SilvermanATAEC} and
\cite[Section 3.2]{Cremona97}. The formulas for Step 2 are elementary: for $p\ge5$ the singular point of
$\tilde E$ has $x$-coordinate the multiple root of $4x^3+b_2x^2+2b_4x+b_6$, which is $-b_2/12$ for a triple root and
$-(c_6+b_2c_4)/(12c_4)$ for a double root, and its $y$-coordinate is the zero of $\partial F/\partial y=2y+a_1x+a_3$; the
formulas at $3$ and $2$ come from the same computation in characteristic $3$ and $2$. We checked the three
formulas on all $211{,}455$ pairs (curve, bad prime) below conductor $10{,}000$: in every case the translated model has
$p\mid a_3,a_4,a_6$.
\end{proof}

\begin{corollary}[the tame table]\label{cor:reduction:tame-table}
Let $p\ge5$ and let a minimal model of $E/K$ be given. Then
\begin{enumerate}[label=(\roman*)]
\item $v(\Delta)=0$ if and only if the type is $\kod{I}_0$;
\item $v(c_4)=0<v(\Delta)=n$ if and only if the type is $\kod{I}_n$;
\item if $v(c_4)>0$ and $v(j)\ge0$, the type is determined by $v(\Delta)\in\{2,3,4,6,8,9,10\}$, which gives
$\kod{II},\kod{III},\kod{IV},\kod{I}_0^*,\kod{IV}^*,\kod{III}^*,\kod{II}^*$ respectively;
\item if $v(c_4)>0$ and $v(j)<0$, then $v(c_4)=2$, $v(c_6)=3$, $v(\Delta)=n+6$ with $n=-v(j)\ge1$, and the type is
$\kod{I}_n^*$.
\end{enumerate}
More precisely, the triples $(v(c_4),v(c_6),v(\Delta))$ are as in \Cref{tab:reduction:tame}. In particular $f=2$
for every additive type.
\end{corollary}

\begin{table}[ht]
\centering\small
\begin{tabular}{lccccc}
\toprule
type & $v(c_4)$ & $v(c_6)$ & $v(\Delta)$ & $m$ & $f$\\
\midrule
$\kod{I}_0$ & $\ge0$ & $\ge0$ & $0$ & $1$ & $0$\\
$\kod{I}_n$ & $0$ & $0$ & $n$ & $n$ & $1$\\
$\kod{II}$ & $\ge1$ & $1$ & $2$ & $1$ & $2$\\
$\kod{III}$ & $1$ & $\ge2$ & $3$ & $2$ & $2$\\
$\kod{IV}$ & $\ge2$ & $2$ & $4$ & $3$ & $2$\\
$\kod{I}_0^*$ & $\ge2$ & $\ge3$ & $6$ & $5$ & $2$\\
$\kod{I}_n^*$ & $2$ & $3$ & $n+6$ & $n+5$ & $2$\\
$\kod{IV}^*$ & $\ge3$ & $4$ & $8$ & $7$ & $2$\\
$\kod{III}^*$ & $3$ & $\ge5$ & $9$ & $8$ & $2$\\
$\kod{II}^*$ & $\ge4$ & $5$ & $10$ & $9$ & $2$\\
\bottomrule
\end{tabular}
\caption{Valuations of a minimal model at $p\ge5$. For $\kod{I}_0^*$ at least one of $v(c_4)=2$, $v(c_6)=3$ holds;
$v(c_4)=\infty$ or $v(c_6)=\infty$ (when $c_4=0$ or $c_6=0$) is allowed wherever ``$\ge$'' appears.}\label{tab:reduction:tame}
\end{table}

\begin{proof}
As $p\ge5$, the change $(1/6,\ast,\ast,\ast)$ to the short model $y^2=x^3+Ax+B$ with $A=-27c_4$, $B=-54c_6$ has $u$ a unit,
so the short model is integral and minimal, with $v(A)=v(c_4)$, $v(B)=v(c_6)$ and the same $v(\Delta)$, where
$\Delta=-16(4A^3+27B^2)$ up to a unit. For it $a_1=a_2=a_3=0$, $a_4=A$, $a_6=B$, $b_2=0$, $b_6=4B$, $b_8=-A^2$. Run
\Cref{thm:reduction:tate-algorithm}.

If $v(\Delta)=0$ the type is $\kod{I}_0$. If $v(\Delta)>0$ and $v(A)=0$ the reduction is multiplicative, of type
$\kod{I}_n$ with $n=v(\Delta)$; then $v(B)=0$ too, since $4A^3\equiv-27B^2$. Otherwise $v(A)\ge1$, hence $v(B)\ge1$, the
singular point is already at $(0,0)$ and $b_2=0$, so we pass to Step 3.
\begin{itemize}
\item Step 3: $\varpi^2\nmid a_6$ means $v(B)=1$: type $\kod{II}$; then $v(27B^2)=2<3\le v(4A^3)$, so $v(\Delta)=2$.
\item Step 4: $\varpi^3\nmid b_8=-A^2$ means $v(A)=1$ (with $v(B)\ge2$): type $\kod{III}$, and
$v(\Delta)=\min(3v(A),2v(B))=3$, as $2v(B)\ge4$. (Throughout, $v(\Delta)=v(4A^3+27B^2)$, which equals $\min(3v(A),2v(B))$ whenever $3v(A)\ne2v(B)$.)
\item Step 5: $\varpi^3\nmid b_6=4B$ means $v(B)=2$ (with $v(A)\ge2$): type $\kod{IV}$, $v(\Delta)=4$, as $3v(A)\ge6>4$.
\item Step 6: now $v(A)\ge2$, $v(B)\ge3$, and the conditions of Step 6 hold without change. $P(T)=T^3+aT+b$ with
$a=A/\varpi^2$, $b=B/\varpi^3$, whose discriminant is $-(4a^3+27b^2)=\Delta/(16\varpi^6)$ up to a unit. So $\tilde P$ has distinct
roots if and only if $v(\Delta)=6$: type $\kod{I}_0^*$.
\item $\tilde P$ has a triple root $\rho$ if and only if $3\rho=0$ (no $T^2$ term), so $\rho=0$, that is, $v(A)\ge3$ and
$v(B)\ge4$. If $v(\Delta)>6$ and the root is not triple, then $v(A)=2$ or $v(B)=3$; but $v(A)=2$, $v(B)\ge4$ or
$v(A)\ge3$, $v(B)=3$ would give $v(\Delta)=6$. So Step 7 occurs exactly when $v(A)=2$, $v(B)=3$, $v(\Delta)>6$; then
$v(j)=3v(c_4)-v(\Delta)=6-v(\Delta)<0$, and by \Cref{thm:reduction:potmult-additive} the type is $\kod{I}_n^*$ with
$n=-v(j)=v(\Delta)-6$.
\item Step 8: $v(A)\ge3$, $v(B)\ge4$, root already at $0$; $Y^2-a_{6,4}$ has distinct roots if and only if $a_{6,4}=B/\varpi^4$ is a
unit ($p$ is odd), that is $v(B)=4$: type $\kod{IV}^*$, and $v(\Delta)=\min(3v(A),2v(B))=8$ since $3v(A)\ge9$.
\item Step 9: $v(B)\ge5$; the double root of $Y^2-a_{6,4}$ is already $0$, so no translation is needed; $\varpi^4\nmid a_4$ means
$v(A)=3$: type $\kod{III}^*$, and $v(\Delta)=9$ since $3v(A)=9<10\le2v(B)$.
\item Step 10: $v(A)\ge4$; $\varpi^6\nmid a_6$ means $v(B)=5$: type $\kod{II}^*$, and $v(\Delta)=10$ since $2v(B)=10<12\le3v(A)$.
\item Step 11 would need $v(A)\ge4$, $v(B)\ge6$, impossible for a minimal model (\Cref{prop:reduction:min-criterion}(ii)).
\end{itemize}
This gives \Cref{tab:reduction:tame}. In (iii), $v(j)\ge0$ excludes $\kod{I}_n$ and $\kod{I}_n^*$ ($v(j)=-n$ there), and
the seven remaining additive types have the distinct values $v(\Delta)=2,3,4,6,8,9,10$. The values of $f$ follow from
$f=v(\Delta)+1-m$.
\end{proof}

All $9{,}699$ potentially good additive pairs (curve, $p\ge5$) of Cremona's tables below conductor $10{,}000$ have
$v_p(\Delta)$ as in (iii), and all pairs with $p\ge5$ fit \Cref{tab:reduction:tame}. At $p=2,3$ no such table exists:
Papadopoulos \cite{Papadopoulos93} gives the corresponding tables, which need congruence conditions on the unit
parts of $c_4$, $c_6$, $\Delta$ besides the valuations.

\begin{proposition}[the cubic criterion for $\kod{I}_0^*$]\label{prop:reduction:I0star-cubic}
Let $p$ be odd and let $E/K$ have a model $y^2=x^3+a_2x^2+a_4x+a_6$ with $\varpi\mid a_2$, $\varpi^2\mid a_4$, $\varpi^3\mid a_6$.
Then $E$ has type $\kod{I}_0^*$ if and only if
\[
P(T)=T^3+(a_2/\varpi)T^2+(a_4/\varpi^2)T+a_6/\varpi^3
\]
has distinct roots modulo $\varpi$, and then $c=1+\#\{\text{roots of }\tilde P\text{ in }k\}$. Conversely every curve of
type $\kod{I}_0^*$ at an odd $p$ has such a model. In particular, if $y^2=x^3+a_2x^2+a_4x+a_6$ is a model over $\Z$ with
good reduction at an odd prime $q$, its twist by $q^*$,
\[
y^2=x^3+q^*a_2x^2+(q^*)^2a_4x+(q^*)^3a_6,
\]
has cubic $T^3+\epsilon a_2T^2+a_4T+\epsilon a_6$ with $\epsilon=q^*/q=\pm1$, whose roots are $\epsilon$ times those of
$x^3+a_2x^2+a_4x+a_6$; so it has type $\kod{I}_0^*$ at $q$.
\end{proposition}

\begin{proof}
For the given model $a_1=a_3=0$, so $b_2=4a_2$, $b_6=4a_6$, $b_8=4a_2a_6-a_4^2$, and $(0,0)$ is the singular point of
$\tilde E$ (the reduced equation is $y^2=x^3$). Run \Cref{thm:reduction:tate-algorithm}: Step 2 passes as
$\varpi\mid b_2$; Step 3 passes as $v(a_6)\ge3$; Step 4 passes as $v(b_8)\ge\min(v(a_2)+v(a_6),2v(a_4))\ge4$; Step 5
passes as $v(b_6)\ge3$; and the conditions of Step 6 already hold. So the type is $\kod{I}_0^*$ exactly when $\tilde P$
has distinct roots, with the stated $c$.

Conversely let $E$ have type $\kod{I}_0^*$ at odd $p$. Completing the square (a change $(1,0,s,t)$ with
$s=-a_1/2$, $t=-a_3/2$, integral since $p$ is odd) gives a minimal model $y^2=f(x)=x^3+a_2x^2+a_4x+a_6$. The singular point
of $\tilde E$ is $(x_0,0)$ with $x_0$ a multiple root of $\tilde f$; translating $x$ by a lift of $x_0$ keeps the shape
of the equation and gives $\varpi\mid a_4,a_6$. As the reduction is additive, $\tilde f$ has a triple root at $0$, so
$\varpi\mid a_2$. Since Steps 3, 4 and 5 are passed, $\varpi^2\mid a_6$, $\varpi^3\mid b_8=4a_2a_6-a_4^2$ and $\varpi^3\mid b_6=4a_6$.
Then $v(4a_2a_6)\ge3$ gives $v(a_4^2)\ge3$, so $\varpi^2\mid a_4$, and $\varpi^3\mid a_6$: the model has the stated shape.

For the twist: $P(T)=T^3+(q^*/q)a_2T^2+(q^*/q)^2a_4T+(q^*/q)^3a_6=T^3+\epsilon a_2T^2+a_4T+\epsilon a_6$, as $\epsilon^2=1$.
If $\theta$ is a root of $g(x)=x^3+a_2x^2+a_4x+a_6$ then
$(\epsilon\theta)^3+\epsilon a_2(\epsilon\theta)^2+a_4\epsilon\theta+\epsilon a_6=\epsilon g(\theta)=0$. The roots of $\tilde g$ are distinct because the
model has good reduction at $q$ ($\Delta=16\operatorname{disc}(g)$ and $q$ is odd), so those of $\tilde P$ are too.
\end{proof}

\begin{warning}\label{warn:reduction:sign}
The sign $\epsilon=q^*/q$ matters when $q\equiv3\pmod4$: the cubic of the twist is $T^3+\epsilon a_2T^2+a_4T+\epsilon a_6$, not
the cubic of $E$. Its roots are $\epsilon$ times the old ones, so the conclusion ``distinct roots'' is unaffected, and so is the
number of roots in $k$, which gives $c$. The twist must be by $q^*$, not by $q$; see
\Cref{warn:reduction:q-versus-qstar}.
\end{warning}

\begin{convention}[PARI's Kodaira codes]\label{conv:reduction:pari-codes}
For a curve \texttt{E} and a prime \texttt{p}, PARI's function \texttt{elllocalred(E,p)} returns the vector $[f_p,k_p,[u,r,s,t],c_p]$: the
conductor exponent, the Kodaira code, the change of coordinates to a minimal model at $p$, and the Tamagawa number. The function
\texttt{ellglobalred(E)} returns $[N,[u,r,s,t],\prod_pc_p,F,L]$ with $F$ the factorisation matrix of $N$ and $L$ the list of local data at the bad
primes, in the order of $F$; in PARI 2.17.2 the entries of $L$ are $[f_p,k_p,0,c_p]$. So the paper's \texttt{R[4]} is the
factorised conductor and \texttt{R[4][,1]} its column of primes; with \texttt{K = R[5]}, the entry \texttt{K[i][2]} is the
Kodaira code at the $i$-th bad prime. The codes are given in \Cref{tab:reduction:codes}.
\end{convention}

\begin{table}[ht]
\centering\small
\begin{tabular}{rlcccc}
\toprule
code $k_p$ & symbol & $m$ & $\#\Phig$ & $f$ at $p\ge5$ & $v(\Delta_{\min})$ at $p\ge5$\\
\midrule
$1$ & $\kod{I}_0$ & $1$ & $1$ & $0$ & $0$\\
$n+4$ ($n\ge1$) & $\kod{I}_n$ & $n$ & $n$ & $1$ & $n$\\
$2$ & $\kod{II}$ & $1$ & $1$ & $2$ & $2$\\
$3$ & $\kod{III}$ & $2$ & $2$ & $2$ & $3$\\
$4$ & $\kod{IV}$ & $3$ & $3$ & $2$ & $4$\\
$-1$ & $\kod{I}_0^*$ & $5$ & $4$ & $2$ & $6$\\
$-(n+4)$ ($n\ge1$) & $\kod{I}_n^*$ & $n+5$ & $4$ & $2$ & $n+6$\\
$-4$ & $\kod{IV}^*$ & $7$ & $3$ & $2$ & $8$\\
$-3$ & $\kod{III}^*$ & $8$ & $2$ & $2$ & $9$\\
$-2$ & $\kod{II}^*$ & $9$ & $1$ & $2$ & $10$\\
\bottomrule
\end{tabular}
\caption{PARI's Kodaira codes. The formulas $n+4$ and $-(n+4)$ hold for $n\ge1$ only: $-(0+4)=-4$ is $\kod{IV}^*$, and
$\kod{I}_0^*$ is $-1$.}\label{tab:reduction:codes}
\end{table}

So \texttt{k >= 5} detects multiplicative reduction, \texttt{k <= -5} the symbols $\kod{I}_n^*$ with $n\ge1$, and
\texttt{abs(k) == 3} the types $\kod{III}$, $\kod{III}^*$. Since \texttt{ellglobalred} lists only the bad primes, the code
$1$ never occurs in its output.

\begin{example}[Tate's algorithm on the paper's curves]\label{ex:reduction:tate-running}
All models are the global minimal models of Cremona's tables; every number below was recomputed with PARI.
\begin{enumerate}[label=(\alph*)]
\item \curve{121a1} $=[1,1,1,-30,-76]$ at $11$: $v(\Delta)=2$, $v(c_4)=1$. Here $b_2=a_1^2+4a_2=5$ and $11\mid c_4$, so the
formula $r=-b_2/12$ of the theorem applies: $12r\equiv-5\pmod{11}$ gives $r\equiv6$, and $t=-(a_1r+a_3)/2=-7/2\equiv2\pmod{11}$. Step 2 with
$(r,t)=(6,2)$ gives $[1,19,11,88,-22]$, with $11\mid a_3,a_4,a_6$ and $11\mid b_2=1+4\cdot19=77$; Step 3: $v(a_6)=1<2$. Type $\kod{II}$, $c=1$,
$f=v(\Delta)=2$.
\item \curve{121c2} $=[1,1,0,-3632,82757]$ at $11$: $v(\Delta)=8$. Step 2 with $(r,t)=(6,8)$ gives
$[1,19,22,-3520,61105]$, with $v(a_6)=2$, $v(b_8)\ge3$, $v(b_6)\ge3$, so Steps 3 to 5 are passed. Step 6 with
$(s,t)=(5,110)$ gives $[11,-11,242,-4840,46585]$ and $P(T)=T^3-T^2-40T+35\equiv(T-4)^3\pmod{11}$: a triple root, so
Step 8. Type $\kod{IV}^*$, $c=1$, $f=2$.
\item \curve{539d1} $=[0,1,1,-16,-66]$ at $7$: Step 2 with $(r,t)=(2,3)$ gives $[0,7,7,0,-98]$; Step 6 with $(s,t)=(0,21)$
gives $[0,7,49,0,-686]$ and $P(T)=T^3+T^2-2=(T-1)(T^2+2T+2)$, where $T^2+2T+2$ has discriminant $-4\equiv3$, a non-square
modulo $7$. Distinct roots, one in $\F_7$: type $\kod{I}_0^*$, $c=2$, $f=2$.
\item \curve{539c1} $=[1,0,1,170,-3237]$ at $7$: Step 2 with $(r,t)=(4,1)$, then Step 6 with $(s,t)=(3,21)$ give
$[7,0,49,49,-3087]$ and $P(T)=T^3+T-9\equiv(T-1)(T-3)^2\pmod 7$. Step 7: translating the double root by $r=21$ gives
$[7,63,196,1372,7203]$; the quadratic $Y^2+4Y-3\equiv(Y-5)^2$ has a double root, and $t=245$ gives
$[7,63,686,-343,-100842]$; the quadratic $2X^2+6X+1\equiv2(X-2)^2$ has a double root, and $r=98$ gives
$[7,357,1372,40817,1411788]$; the quadratic $Y^2+4Y+2\equiv(Y-1)(Y-2)$ has distinct roots in $\F_7$. So $n=3$: type
$\kod{I}_3^*$, $c=4$, and $f=v(\Delta)-4-n=9-7=2$.
\item \curve{726a1} $=[1,1,0,-35,-51]$: at $11$, Step 2 with $(r,t)=(6,8)$ gives $[1,19,22,77,-121]$, with $v(a_6)=2$ and
$v(b_8)=2<3$: type $\kod{III}$, $c=2$. At $2$: $v(\Delta)=6$, $v(c_4)=0$, and the tangent cone gives nonsplit reduction:
$\kod{I}_6$, $c=2$. At $3$: $\kod{I}_3$ nonsplit, $c=1$.
\item \curve{20a1} $=[0,1,0,4,4]$ at $2$: $v(\Delta)=8$, $v(c_4)=4$. The singular point is already at $(0,0)$, and Steps 3 to
5 are passed ($v(a_6)=2$, $b_8=0$, $v(b_6)=4$). Step 6 with $(s,t)=(1,2)$: for a change $(1,0,s,t)$ of a model with
$a_1=a_3=0$ the formulas of \Cref{ch:ellcurves} reduce to $a_1'=2s$, $a_2'=a_2-s^2$, $a_3'=2t$, $a_4'=a_4-2st$, $a_6'=a_6-t^2$, which
here give $[2,0,4,0,0]$, that is $y^2+2xy+4y=x^3$. Now $2\mid a_1,a_2$, $4\mid a_3,a_4$, $8\mid a_6$, and $P(T)=T^3$, a triple
root; Step 8: $Y^2+a_{3,2}Y-a_{6,4}=Y^2+Y$ has the distinct roots $0,1\in\F_2$. Type $\kod{IV}^*$, $c=3$, $f=8+1-7=2$.
\item \curve{99a1} $=[1,-1,1,-2,0]$ at $3$: Step 2 with $(r,t)=(2,0)$ gives $[1,5,3,6,0]$, $v(b_2)=1$, $a_6=0$, and
$v(b_8)=2<3$: type $\kod{III}$, $c=2$, $f=v(\Delta)-1=2$.
\item \curve{980i1} $=[0,0,0,1568,-72716]$: at $2$, type $\kod{IV}^*$ with $c=1$; at $5$, $\kod{I}_1$; at $7$, the model
already satisfies the conditions of Step 6 ($v(a_4)=2$, $v(a_6)=3$), with $P(T)=T^3+32T-212\equiv(T-2)(T+1)^2\pmod7$; the
subprocedure of Step 7 passes four double roots and stops at $s=5$ with the quadratic $Y^2+1$, irreducible over $\F_7$:
type $\kod{I}_5^*$, $c=2$, $v_7(\Delta)=11=5+6$.
\end{enumerate}
\end{example}

\begin{paperlink}
The Example of Section 3 of the paper (types $\kod{II}$ of \curve{121a1}, $\kod{I}_3^*$ of \curve{539c1}, $\kod{I}_5^*$ of
\curve{980i1}), the Remark after it (\curve{99a1}, \curve{539d1}, \curve{726a1}), Section sec:anomaly (\curve{726a1}, type
$\kod{III}$ at $11$) and the appendix (\texttt{rule} on \curve{121a1} and \curve{726a1}) all use the types computed here. In the
appendix, \texttt{k >= 5} selects the multiplicative primes (where the naive factor is removed), \texttt{abs(k) != 3}
excludes $\kod{III}$, $\kod{III}^*$ unless \texttt{iii = 1}, and \texttt{k == 4 || k == 2}, \texttt{k == -4 || k == -2},
\texttt{k == -3} select $\kod{IV},\kod{II}$, then $\kod{IV}^*,\kod{II}^*$, then $\kod{III}^*$ for the weights $w_3$ at $q=3$.
\end{paperlink}

\section{The conductor and the semistability defect}\label{sec:reduction:conductor}

This section defines the conductor exponent $f_p$ and the conductor $N$, and studies how inertia acts on the Tate
module at a potentially good prime: through a finite group $G_{\mathrm{ss}}$, whose order is given by Kraus's formula
$12/\gcd(v(\Delta),12)$ when $p\ge5$. The paper uses the conductor exponents in Lemma lem:twist and Theorem thm:qual,
the defect in Lemma lem:IV, Remark rem:types and Lemma lem:euler, and the wild primes in Remark rem:wild and
Section sec:wild. The section ends with the theorem that types $\kod{IV}$ and $\kod{IV}^*$ at $2$ are tame with inertia
image cyclic of order $3$ (\Cref{cor:reduction:IV-at-2}), which removes the exception that the original version of
the paper made in Proposition prop:forms.

\begin{convention}[wild part and defect]\label{conv:reduction:sw}
The Swan conductor (wild part of the conductor) is written $\operatorname{sw}$, as in \Cref{ch:local}; the letter
$\delta_p$ is reserved for the paper's $\dim V^{I_p}$. For potentially good reduction the group measuring the defect of
semistability is written $G_{\mathrm{ss}}$ (it is called $\Phi_K$ in some sources; we avoid that letter, which denotes
component groups here), and $e=\#G_{\mathrm{ss}}$ is the \emph{semistability defect}.
\end{convention}

\begin{definition}[conductor exponent]\label{def:reduction:conductor-exponent}
Let $E/K$ and $\ell\ne p$ a prime, and let $V=V_\ell E$ as a representation of $G_K$. The \emph{conductor exponent} of
$E$ is the Artin conductor of $V$ (\Cref{ch:local}):
\[
f(E)=a(V)=\dim V-\dim V^{I_K}+\operatorname{sw}(V),
\]
where $\operatorname{sw}(V)\ge0$ is the Swan conductor, which vanishes if and only if the wild inertia group $P_K$ acts
trivially on $V$. For $E/\Q$ and a prime $p$ we write $f_p=f_p(E)$ for the exponent of $E/\Q_p$.
\end{definition}

The conductor exponent measures how far $V_\ell E$ is from being unramified. The tame part $\dim V-\dim V^{I_K}$ counts
the dimensions that inertia moves, which is $0$, $1$ or $2$; the Swan term is a weighted count of how deep in the
ramification filtration the action of wild inertia reaches, and is nonzero exactly when the pro-$p$ group $P_K$ acts
nontrivially, which forces $p$ to divide the order of the image of inertia. For \curve{11a1} at $11$ the action of inertia fixes a line (\Cref{prop:reduction:conductor-values}(ii)
below), so $f_{11}=1$; for \curve{20a1} at $2$, of type $\kod{IV}^*$, inertia fixes no vector and acts through a group of order $3$, prime to $2$, so
$f_2=2$ (\Cref{cor:reduction:IV-at-2}).

That $f(E)$ does not depend on $\ell$ is a theorem of Ogg and of Serre and Tate \cite{Ogg,SerreTate68}; we do not prove it.
The following values are proved here, given the criterion of N\'eron, Ogg and Shafarevich (\Cref{ch:galrep}: $E$ has
good reduction if and only if $I_K$ acts trivially on $T_\ell E$) and two results of this chapter proved below.

\begin{proposition}[values of the conductor exponent]\label{prop:reduction:conductor-values}
Let $\ell\ne p$.
\begin{enumerate}[label=(\roman*)]
\item Good reduction: $V^{I_K}=V$ and $f=0$.
\item Multiplicative reduction: $V^{I_K}$ is a line, $I_K$ acts tamely through a nontrivial unipotent, and $f=1$.
\item Additive reduction: $V^{I_K}=0$ and $f=2+\operatorname{sw}(V)\ge2$.
\item If $p\ge5$, $\operatorname{sw}(V)=0$, so $f=2$ for additive reduction. If $p$ is odd and the reduction is additive
and potentially multiplicative, also $\operatorname{sw}(V)=0$ and $f=2$.
\end{enumerate}
\end{proposition}

\begin{proof}
(i) is \Cref{cor:reduction:torsion-injects}(ii). (ii) By \Cref{prop:reduction:tate-torsion}(iii), in a suitable basis
$G_K$ acts on $T_\ell E$ by $\sm{\chi_\ell\varepsilon}{\varepsilon\kappa}{0}{\varepsilon}$ with $\varepsilon$ unramified. On $I_K$ the
cyclotomic character $\chi_\ell$ is trivial ($\mu_{\ell^\infty}\subset K^{\ur}$), $\varepsilon$ is trivial, and $\kappa$
restricts to $\nu$ times the tame character $t_\ell\colon I_K\to\Z_\ell(1)$, with $\nu=v(q_E)\ge1$ (\Cref{ch:local}); as in
Appendix B, the tame character always carries its domain or an argument, and $t_\ell(E,q)$ always its two arguments. So
$\sigma\in I_K$ acts by $\sm{1}{\nu t_\ell(\sigma)}{0}{1}$: unipotent, nontrivial, trivial on $P_K$, with invariants the line spanned by $e_1$; so
$f=2-1+0=1$.

(iii) Suppose first that the reduction is potentially good. By \Cref{thm:reduction:serre-tate} the image $\rho(I_K)$ is
finite, and it lies in $\SL_2$ since $\det\rho=\chi_\ell$ is trivial on $I_K$. If $V^{I_K}\ne0$, then in a basis $(v,w)$
with $v$ fixed every element of $\rho(I_K)$ has the form $\sm1b0d$ with $d=1$; a unipotent matrix of finite order in
characteristic $0$ is the identity, so $\rho(I_K)=1$ and the reduction would be good by the criterion of N\'eron, Ogg
and Shafarevich. If the reduction is potentially multiplicative, then by \Cref{thm:reduction:tate-uniformisation} $E$ is
the twist of a Tate curve by a ramified quadratic character $\chi$, and by \Cref{prop:reduction:tate-torsion}
$\sigma\in I_K$ acts by $\chi(\sigma)\sm{1}{\nu t_\ell(\sigma)}{0}{1}$. An element with $\chi(\sigma)=-1$ has the single eigenvalue
$-1$ and fixes no nonzero vector. In both cases $V^{I_K}=0$ and $f=2+\operatorname{sw}(V)$.

(iv) If the reduction is potentially good, $\rho(I_K)\simeq G_{\mathrm{ss}}$ has order dividing $24$
(\Cref{thm:reduction:serre-tate}(iii)); for $p\ge5$ this is prime to $p$, so the pro-$p$ group $P_K$ acts trivially and
$\operatorname{sw}(V)=0$. If it is potentially multiplicative and $p$ is odd, the quadratic character $\chi$ is trivial on the
pro-$p$ group $P_K$, and so is $t_\ell\colon I_K\to\Z_\ell(1)$; so $P_K$ acts trivially again.
\end{proof}

\begin{theorem}[Ogg's formula]\label{thm:reduction:ogg}
Let $E/K$, with $K/\Q_p$ finite, and let $m$ be the number of irreducible components over $\bar k$ of the special
fibre of its minimal regular model (\Cref{tab:reduction:kodaira}). Then
\[
f=v(\Delta_{\min})+1-m .
\]
\end{theorem}

We do not prove this; it is due to Ogg \cite{Ogg}, whose argument has a gap in some cases at $p=2$, and in general to
Saito \cite{Saito88}; see \cite[IV.11]{SilvermanATAEC}. For $p\ge5$ it can be read off \Cref{tab:reduction:tame}. PARI computes
$f$ through this formula, so a comparison with PARI's output is a consistency check and not an independent test.

\begin{theorem}[bounds for the conductor exponent]\label{thm:reduction:conductor-bounds}
Let $E/\Q$. Then $f_p\le2$ for $p\ge5$, $f_3\le5$ and $f_2\le8$.
\end{theorem}

The first bound is \Cref{prop:reduction:conductor-values}; we do not prove the other two. They are special cases of the
bounds of Brumer and Kramer \cite{BrumerKramer94} and of Lockhart, Rosen and Silverman \cite{LockhartRosenSilverman93} for
abelian varieties. Below conductor $10{,}000$ every value $f_2\in\{2,\dots,8\}$ and $f_3\in\{2,\dots,5\}$ occurs at an additive
prime ($4{,}750$, $5{,}390$, $7{,}067$, $2{,}278$, $3{,}946$, $892$, $388$ pairs at $2$; $13{,}194$, $3{,}014$, $914$, $356$ at
$3$).

\begin{definition}[the conductor]\label{def:reduction:conductor}
The \emph{conductor} of $E/\Q$ is $N=\prod_pp^{f_p}$.
\end{definition}

\begin{proposition}[properties of the conductor]\label{prop:reduction:conductor-properties}
Let $E/\Q$ have conductor $N$.
\begin{enumerate}[label=(\roman*)]
\item $p\mid N$ if and only if $E$ has bad reduction at $p$; $p\,\|\,N$ if and only if the reduction is multiplicative;
$p^2\mid N$ if and only if it is additive; and $p^3\nmid N$ for $p\ge5$.
\item Isogenous curves over $\Q$ have the same conductor.
\end{enumerate}
\end{proposition}

\begin{proof}
(i) is \Cref{prop:reduction:conductor-values}. (ii) An isogeny $E\to E'$ over $\Q$ induces a $G_\Q$-equivariant
isomorphism $V_\ell E\simeq V_\ell E'$ (\Cref{ch:ellcurves}), and $f_p$ depends only on $V_\ell E$ restricted to $G_{\Q_p}$, for
any $\ell\ne p$.
\end{proof}

The conductors of the paper's curves, all checked with PARI, are in the column $N$ of
\Cref{ex:reduction:running-table}.

\begin{theorem}[Serre and Tate]\label{thm:reduction:serre-tate}
Let $E/K$ have potentially good reduction, $K/\Q_p$ finite.
\begin{enumerate}[label=(\roman*)]
\item There is a smallest extension $L$ of $K^{\ur}$ over which $E$ has good reduction. It is Galois over $K^{\ur}$, and
$L=K^{\ur}(E[m])$ for every integer $m\ge3$ prime to $p$.
\item Put $G_{\mathrm{ss}}=\Gal(L/K^{\ur})$ and $e=\#G_{\mathrm{ss}}$. For every prime $\ell\ne p$ the kernel of
$I_K\to\Aut(T_\ell E)$ is $\Gal(\bar K/L)$, so the image of $I_K$ in $\Aut(T_\ell E)$ is isomorphic to $G_{\mathrm{ss}}$,
independently of $\ell$.
\item $G_{\mathrm{ss}}$ embeds in $\SL_2(\Z/m)$ for every $m\ge3$ prime to $p$; in particular $e$ divides $24$. Every
element of $G_{\mathrm{ss}}$ has order $1$, $2$, $3$, $4$ or $6$. If $p\ge5$, $G_{\mathrm{ss}}$ is cyclic of order $1$, $2$, $3$, $4$
or $6$.
\item $G_{\mathrm{ss}}$ acts faithfully on the reduction $\tilde E_L$ of $E$ over $L$; if $e>2$ then $j(\tilde E_L)\in\{0,1728\}$.
\end{enumerate}
\end{theorem}

\begin{proof}[Sketch of proof]
(ii) given (i): for an extension $L'$ of $K^{\ur}$ inside $\bar K$, $E$ has good reduction over $L'$ if and only if
$\Gal(\bar K/L')$ (the inertia group of $L'$) acts trivially on $T_\ell E$, by the criterion of N\'eron, Ogg and
Shafarevich over $L'$ (\Cref{ch:galrep}). So the fixed field of the kernel of $I_K\to\Aut(T_\ell E)$ is the smallest such
$L'$; this proves (i) except for the description by $E[m]$, and (ii). For (iii): the image of $I_K$ in $\Aut(E[m])$ is
$\Gal(K^{\ur}(E[m])/K^{\ur})=G_{\mathrm{ss}}$ by (i), and its determinant is the cyclotomic character modulo $m$, trivial on
$I_K$. So $e$ divides $\#\SL_2(\Z/m)$ for every $m\ge3$ prime to $p$. If $p\ne3$, take $m=3$: $\#\SL_2(\F_3)=24$. If $p=3$, take
$m=4$ and $m=5$: $e$ divides $\#\SL_2(\Z/4)=48$ and $\#\SL_2(\F_5)=120$, hence $\gcd(48,120)=24$. The characteristic polynomial of an element of $\rho(I_K)$ has integer coefficients
independent of $\ell$ \cite[Section 2]{SerreTate68}, and its roots are roots of unity of degree at most $2$
over $\Q$, of order $1,2,3,4$ or $6$. If $p\ge5$, $p\nmid24$, so $L/K^{\ur}$ is tamely ramified and $G_{\mathrm{ss}}$ is a
quotient of the procyclic tame inertia group (\Cref{ch:local}), hence cyclic, of an order that is the order of one of its
elements. The description $L=K^{\ur}(E[m])$ in (i) and statement (iv) are in \cite[Section 2]{SerreTate68};
see also \cite[Section 5.6]{Serre72}. We do not prove them.
\end{proof}

\begin{theorem}[Kraus]\label{thm:reduction:kraus}
Let $K/\Q_p$ be finite and let $E/K$ have potentially good reduction.
\begin{enumerate}[label=(\roman*)]
\item If $p\ge5$, $G_{\mathrm{ss}}$ is cyclic of order $e=12/\gcd(v(\Delta_{\min}),12)$. By type: $\kod{I}_0$, $\kod{II}$, $\kod{III}$,
$\kod{IV}$, $\kod{I}_0^*$, $\kod{IV}^*$, $\kod{III}^*$, $\kod{II}^*$ give $e=1,6,4,3,2,3,4,6$.
\item If $p=3$, $G_{\mathrm{ss}}$ is cyclic of order $1,2,3,4$ or $6$, or the semidirect product $\Z/3\rtimes\Z/4$ of order $12$.
\item If $p=2$, $G_{\mathrm{ss}}$ is cyclic of order $1,2,3,4$ or $6$, or the quaternion group $Q_8$, or $\SL_2(\F_3)$ of order
$24$. For $K=\Q_2$ and $K=\Q_3$, Kraus determines $G_{\mathrm{ss}}$ from $v(c_4)$, $v(c_6)$, $v(\Delta)$ and congruences on
their unit parts.
\end{enumerate}
\end{theorem}

\begin{proof}
(i), given \Cref{thm:reduction:serre-tate}. Put $e=12/\gcd(v(\Delta),12)$ for a minimal model, so that $ev(\Delta)/12$ is an
integer, and let $K'=K^{\ur}(\varpi^{1/e})$, a cyclic extension of $K^{\ur}$ of degree $e$, totally ramified with uniformiser
$\varpi'=\varpi^{1/e}$ ($K^{\ur}$ contains $\mu_e$ as $p\nmid e$). Let $v'$ be the valuation of $K'$, so $v'=ev$ on $K$. Apply
$(u,0,0,0)$ with $u=\varpi'^{\,ev(\Delta)/12}$. The new invariants satisfy
\begin{align*}
v'(\Delta')&=ev(\Delta)-12\cdot ev(\Delta)/12=0,\\
v'(c_4')&=ev(c_4)-4ev(\Delta)/12=e\,v(j)/3\ge0,\\
v'(c_6')&=ev(c_6)-6ev(\Delta)/12\ge0,
\end{align*}
using $v(j)\ge0$ (\Cref{thm:reduction:potential-criterion}) for $c_4'$, and $2v(c_6)\ge\min(3v(c_4),v(\Delta))=v(\Delta)$ (from
$c_6^2=c_4^3-1728\Delta$) for $c_6'$. As $p\ge5$, the model $y^2=x^3-27c_4'x-54c_6'$ is integral with unit discriminant, so $E$
has good reduction over $K'$. Conversely, if $E$ has good reduction over an extension $L'$ of $K^{\ur}$ of ramification
index $e'$, a minimal model over $L'$ has unit discriminant $u^{-12}\Delta$, so $e'v(\Delta)\equiv0\pmod{12}$ and $e\mid e'$. Apply
this to the smallest extension $L$ of \Cref{thm:reduction:serre-tate}: $e$ divides $[L:K^{\ur}]$. But $L\subseteq K'$, as $E$
has good reduction over $K'$, so $[L:K^{\ur}]\le e$. Hence $L=K'$ and $G_{\mathrm{ss}}=\Gal(K'/K^{\ur})\simeq\mu_e$ is cyclic of
order $e$. The list by type is \Cref{cor:reduction:tame-table}(iii): $v(\Delta)=0,2,3,4,6,8,9,10$ gives
$e=1,6,4,3,2,3,4,6$.

(ii) and (iii) we do not prove; see Kraus \cite{Kraus90} and \cite[Section 5.6]{Serre72}.
\end{proof}

The theorem numbers inside \cite{Kraus90} are quoted inconsistently in the secondary literature, so we cite the paper as a
whole.

\begin{theorem}[tame additive reduction]\label{thm:reduction:tame-additive}
Let $K/\Q_p$ be finite and let $E/K$ have additive, potentially good reduction. Then the reduction is tame
(equivalently $\operatorname{sw}=0$, equivalently $f=2$) if and only if
\begin{enumerate}[label=(\alph*)]
\item $p\ge5$; or
\item $p=3$ and the type is $\kod{III}$, $\kod{III}^*$ or $\kod{I}_0^*$; or
\item $p=2$ and the type is $\kod{IV}$ or $\kod{IV}^*$.
\end{enumerate}
\end{theorem}

We do not prove this; see Dokchitser and Dokchitser \cite[Theorem 3.1]{DokDok15}, who deduce it from Kramer and Tunnell
\cite[Proposition 8.20]{KramerTunnell82} at $p=2$ and from Kraus \cite{Kraus90} at $p=3$. The direction ``type $\kod{IV}$ or
$\kod{IV}^*$ at $2$ implies tame'' is proved independently in \Cref{cor:reduction:IV-at-2}. The data agree: below conductor
$10{,}000$, $f_2=2$ occurs exactly for the pairs of type $\kod{IV}$ or $\kod{IV}^*$ at $2$ and for no other additive type
(\Cref{tab:reduction:inertia2}), and at $3$ the types $\kod{III}$, $\kod{III}^*$, $\kod{I}_0^*$ always have $f_3=2$ while $\kod{II}$, $\kod{II}^*$,
$\kod{IV}$, $\kod{IV}^*$ always have $f_3\in\{3,4,5\}$.

\begin{corollary}[types $\kod{IV}$ and $\kod{IV}^*$ at $2$]\label{cor:reduction:IV-at-2}
Let $E/\Q_2$ have Kodaira type $\kod{IV}$ or $\kod{IV}^*$. Then:
\begin{enumerate}[label=(\roman*)]
\item $E$ has potentially good reduction;
\item $G_{\mathrm{ss}}$ is cyclic of order $3$, so for every odd prime $\ell$ the image of $I_2$ in $\Aut(T_\ell E)$ is cyclic of
order $3$, generated by an element $\sigma$ with $\det\sigma=1$ and eigenvalues the two primitive cube roots of unity; on
$E[3]$ the image is a nontrivial unipotent group and $E[3]^{I_2}$ is a line;
\item the reduction is tame and $f_2=2$;
\item every Frobenius lift $\Frob\in G_{\Q_2}$ satisfies $\rho(\Frob)\,\sigma\,\rho(\Frob)^{-1}=\sigma^{-1}$.
\end{enumerate}
\end{corollary}

\begin{proof}
By \Cref{prop:reduction:component-invariants}(i) and \Cref{thm:reduction:kodaira-neron}, $E[3]^{I_2}\simeq\Phig[3]\simeq\Z/3$
is a line.

(i) Suppose $v_2(j)<0$. By \Cref{thm:reduction:tate-uniformisation}, $E$ is the twist of the Tate curve $E_{q_E}$ by a
quadratic character $\chi$, ramified since the reduction is additive, and by \Cref{prop:reduction:tate-torsion} an element
$\sigma\in I_2$ acts on $T_3E$ by $\chi(\sigma)\sm{1}{\kappa(\sigma)}{0}{1}$, where $\kappa(\sigma)=\nu\,t_3(\sigma)$ is $\nu$ times the $3$-adic tame
character $t_3\colon I_2\to\Z_3(1)$ (\Cref{ch:local}; not the term $t_3(E,2)$). On $I_2$ the cyclotomic character $\chi_3$ is
trivial, so the cocycle relation $\kappa(gh)=\chi_3(g)\kappa(h)+\kappa(g)$ makes $\kappa|_{I_2}$ a homomorphism to $\Z_3$, and
$H=\{\sigma\in I_2:\kappa(\sigma)\equiv0\pmod3\}$ is the kernel of its reduction $I_2\to\Z/3$, a subgroup of index $1$ or $3$. The
subgroup $\ker\chi|_{I_2}$ has index $2$, because $\chi$ is ramified. If $H$ were contained in $\ker\chi$, the index $2$ would divide
$[I_2:H]\in\{1,3\}$; so $H\not\subseteq\ker\chi$, and some $\sigma\in H$ has $\chi(\sigma)=-1$. It acts on $E[3]$ by
$-\sm1001=-1$, which fixes no nonzero vector. This contradicts $E[3]^{I_2}\ne0$. So $v_2(j)\ge0$.

(ii) By \Cref{thm:reduction:serre-tate} the image $\rho_3(I_2)$ is a finite subgroup of $\SL_2(\Z_3)$. Reduction modulo $3$ is
injective on it: if $g=1+3^kA$ with $k\ge1$ and $A\not\equiv0\pmod3$ had finite order, a power of it would have prime order $r$
and the same shape (with possibly larger $k$), and the binomial theorem (valid as $1$ and $A$ commute) gives
\begin{gather*}
g^r=1+r3^kA+\tbinom r2 3^{2k}A^2+\cdots\equiv1+r3^kA\pmod{3^{k+1}}\quad(r\ne3),\\
g^3=1+3^{k+1}A+3^{2k+1}A^2+3^{3k}A^3\equiv1+3^{k+1}A\pmod{3^{k+2}},
\end{gather*}
using $2k\ge k+1$, $2k+1\ge k+2$ and $3k\ge k+2$ for $k\ge1$. As $r$ is then a unit and $A\not\equiv0\pmod 3$, neither power is
$1$, a contradiction. The image modulo $3$ is nontrivial (otherwise $\rho_3(I_2)=1$ and the reduction would be good) and fixes
the line $E[3]^{I_2}$ pointwise. In a basis $(v,w)$ of $E[3]$ with $v$ spanning that line, its elements are
$\sm{1}{b}{0}{d}$ with determinant $d=1$: the group $\{\sm1b01\}$ of order $3$. So $\rho_3(I_2)$ is cyclic of order $3$,
hence $G_{\mathrm{ss}}\simeq\Z/3$ and $\rho_\ell(I_2)\simeq\Z/3$ for every $\ell\ne2$ (\Cref{thm:reduction:serre-tate}(ii)). A generator
$\sigma$ has determinant $1$ and $\sigma^3=1$, $\sigma\ne1$; its eigenvalues $\lambda,\lambda^{-1}$ are cube roots of unity, and not both $1$
(a unipotent element of finite order is trivial), so they are $\zeta_3^{\pm1}$.

(iii) The order $3$ is prime to $2$, so the wild inertia group, a pro-$2$ group, acts trivially: $\operatorname{sw}=0$. With
$V^{I_2}=0$ (\Cref{prop:reduction:conductor-values}(iii)), $f_2=2$.

(iv) As $\rho(P_2)=1$, $\rho|_{I_2}$ factors through the tame quotient, and $\sigma=\rho(t)$ for a topological generator $t$ of
it. The tame relation $\Frob\,t\,\Frob^{-1}=t^2$ (\Cref{ch:local}) gives $\rho(\Frob)\sigma\rho(\Frob)^{-1}=\sigma^2=\sigma^{-1}$.
\end{proof}

We checked the conclusions on all pairs of type $\kod{IV}$ or $\kod{IV}^*$ at $2$ below conductor $10{,}000$ (the first two rows of
\Cref{tab:reduction:inertia2}): all have
$v_2(j)\ge0$ and $f_2=2$, and the four roots of the $3$-division polynomial generate over $\Q_2$ fields of ramification
indices $1,3,3,3$, as (ii) predicts (\Cref{prop:reduction:inertia-at-2}).

\begin{paperlink}
Lemma lem:IV of the revised paper is stated for every $q\ne3$; its proof is the argument above (component group, then
injectivity of reduction on finite subgroups of $\GL_2(\Z_3)$). For $q\ge5$, Kraus's formula gives the order $3$ directly: the revised
paper notes after the lemma that ``for $q\ge5$ the order $3$ is also the semistability defect $12/\gcd(v_q(\Delta_E),12)$ of
Kraus'' (the original version wrote ``for $v_q(\Delta_E)\in\{4,8\}$ this is $3$''). The Galois-representation consequences, including the Serre conductor
exponent $1$ of $\bar\rho_{E,3}$ at $q$, are proved in \Cref{ch:galrep}. With (ii) and (iv), the computation of Lemma lem:euler
for $\kod{II},\kod{II}^*,\kod{IV},\kod{IV}^*$ applies verbatim at $q=2$: the trace-zero endomorphisms commuting with $\sigma$ form
the line spanned by $2\sigma+1$, Frobenius inverts $\sigma$ and so acts on that line by $-1=\chi_{-3}(2)$, so
$L_2(A_f,s)^{-1}=1+2^{-s}$, $P_2(X)=1+2X$ and $e_\ell(E,2)=v_\ell(3/2)$, which is $1$ for $\ell=3$ and $0$ for $\ell\ge5$
(\Cref{ch:lfunctions}). This is why Proposition prop:forms holds with no exception (the original version excepted these
types and checked $e_3(E,2)\ge1$ on $92{,}812$ curves), why Remark rem:wild can say that a type $\kod{IV}$ or $\kod{IV}^*$ at $2$
``always has $P_2(X)=1+2X$'', and why case (a) of Theorem thm:qual holds at $q=2$.
\end{paperlink}

\begin{example}[semistability defects of the paper's curves]\label{ex:reduction:defects}
At the additive primes of the running examples: \curve{121a1} at $11$ (type $\kod{II}$, $v(\Delta)=2$): $e=6$; \curve{121c1}
($\kod{IV}$, $v=4$) and \curve{121c2} ($\kod{IV}^*$, $v=8$) at $11$: $e=3$; \curve{539d1} at $7$ ($\kod{I}_0^*$): $e=2$; \curve{726a1}
at $11$ ($\kod{III}$, $v=3$): $e=4$. At the wild primes: \curve{99a1} at $3$ (type $\kod{III}$, tame by
\Cref{thm:reduction:tame-additive}) has $G_{\mathrm{ss}}$ cyclic of order $4$ (\Cref{rem:reduction:wild-at-3});
\curve{20a1}, \curve{140b1} and \curve{980i1} at $2$ (type $\kod{IV}^*$) have $e=3$ by \Cref{cor:reduction:IV-at-2};
\curve{80b1}, the twist of \curve{20a1} by $-1$, has type $\kod{I}_0^*$ and $f_2=4$ at $2$, and $G_{\mathrm{ss}}$ cyclic of order $6$
(the order-$3$ group of \curve{20a1} times the image of the wildly ramified character of $\Q_2(i)$;
\Cref{prop:reduction:inertia-at-2}). The potentially multiplicative primes (\curve{539c1} and \curve{980i1} at $7$, \curve{121d1}
at $11$) have no defect: the image of inertia is infinite (\Cref{prop:reduction:tate-torsion}).
\end{example}

\section{The Tate curve}\label{sec:reduction:tate}

At a prime of multiplicative reduction, and at an additive prime with $v(j)<0$, the curve is (a quadratic twist of)
a Tate curve $\bar K^\times/q_E^{\Z}$. Everything the paper says at such a prime comes from this description: the term
$t_\ell(E,q)=v_\ell(-v_q(j_E))$ of Definition def:terms is the $\ell$-part of the valuation of the Tate parameter; Lemma
lem:tamagawa-exp reads the ramification of $E[\ell^t]$ off an $\ell^t$-th root of $q_E$; Lemma lem:tamfactor uses the
unit part $u$ of $q_E$ and the matrices of the Galois action. This section constructs the Tate curve, defines $q_E$
and $u$, describes $E_{\mathsf q}[\ell^t]$, computes the Tamagawa number, determines the type at an additive
potentially multiplicative prime, and tabulates the Tate parameters of the paper's curves.

For the reader with a physics background: over $\C$, the exponential $z\mapsto e^{2\pi iz}$ identifies the torus
$\C/(\Z+\Z\tau)$ with $\C^\times/\mathsf q^{\Z}$, $\mathsf q=e^{2\pi i\tau}$, the ``multiplicative'' picture of the same lattice. Tate's
observation is that the second description makes sense over a $p$-adic field, where $\C/\Lambda$ does not.

\begin{theorem}[Tate]\label{thm:reduction:tate-curve}
Let $\mathsf q\in K^\times$ with $v(\mathsf q)>0$. Put
\[
s_k(\mathsf q)=\sum_{n\ge1}\frac{n^k\mathsf q^n}{1-\mathsf q^n},\qquad a_4(\mathsf q)=-5s_3(\mathsf q),\qquad
a_6(\mathsf q)=-\frac{5s_3(\mathsf q)+7s_5(\mathsf q)}{12}.
\]
\begin{enumerate}[label=(\roman*)]
\item The series $a_4,a_6$ have coefficients in $\Z$ and converge. The equation
\[
E_{\mathsf q}\colon\ y^2+xy=x^3+a_4(\mathsf q)x+a_6(\mathsf q)
\]
defines an elliptic curve over $K$ with
\begin{gather*}
c_4=1+240s_3(\mathsf q),\qquad c_6=-1+504s_5(\mathsf q),\qquad\Delta=\mathsf q\prod_{n\ge1}(1-\mathsf q^n)^{24},\\
j(E_{\mathsf q})=\frac1{\mathsf q}+744+196884\,\mathsf q+21493760\,\mathsf q^2+\cdots .
\end{gather*}
\item This model is minimal, and $E_{\mathsf q}$ has split multiplicative reduction of type $\kod{I}_n$ with
$n=v(\mathsf q)=-v(j(E_{\mathsf q}))$.
\item With
\[
X(w,\mathsf q)=\sum_{n\in\Z}\frac{\mathsf q^nw}{(1-\mathsf q^nw)^2}-2s_1(\mathsf q),\qquad
Y(w,\mathsf q)=\sum_{n\in\Z}\frac{(\mathsf q^nw)^2}{(1-\mathsf q^nw)^3}+s_1(\mathsf q),
\]
the map $\phi\colon\bar K^\times/\mathsf q^{\Z}\to E_{\mathsf q}(\bar K)$, $w\mapsto(X(w,\mathsf q),Y(w,\mathsf q))$ for $w\notin\mathsf q^{\Z}$ and
$w\mapsto O$ for $w\in\mathsf q^{\Z}$, is an isomorphism of groups commuting with the action of $G_K$. In particular
$E_{\mathsf q}(L)\simeq L^\times/\mathsf q^{\Z}$ for every algebraic extension $L/K$.
\end{enumerate}
\end{theorem}

\begin{proof}[Sketch of proof]
(i) $s_k(\mathsf q)=\sum_{n\ge1}\sigma_k(n)\mathsf q^n$ with $\sigma_k(n)=\sum_{d\mid n}d^k$, so $s_k\in\mathsf q\Z[[\mathsf q]]$. For $a_6$ it suffices that
$5d^3+7d^5\equiv0\pmod{12}$ for every integer $d$: modulo $3$, $d^3(5+7d^2)\equiv d^3\cdot12\equiv0$ if $3\nmid d$, and trivially $d^3\equiv0$ if $3\mid d$; modulo $4$, if $d$ is
odd then $d^2\equiv1$ and $5+7=12$, and if $d$ is even then $8\mid d^3$. Power series with integer coefficients converge
at $\mathsf q$ with $v(\mathsf q)>0$. For the model, $a_1=1$ and $a_2=a_3=0$ give $b_2=a_1^2+4a_2=1$,
$b_4=a_1a_3+2a_4=2a_4$, $b_6=a_3^2+4a_6=4a_6$; then $c_4=b_2^2-24b_4=1-48a_4=1+240s_3$, and
\[
c_6=-b_2^3+36b_2b_4-216b_6=-1+72a_4-864a_6=-1-360s_3+72(5s_3+7s_5)=-1+504s_5 ;
\]
$a_4=-5s_3$ is chosen to make $c_4=1+240s_3$, and then $a_6=-(5s_3+7s_5)/12$ is chosen so that the $s_3$ terms cancel in
$c_6$, leaving $-1+504s_5$. So $c_4$ and $c_6$
are $E_4$ and $-E_6$, the normalised Eisenstein series of weights $4$ and $6$ in the variable $\mathsf q$ (\Cref{ch:modforms}),
and the classical identity $E_4^3-E_6^2=1728\,\mathsf q\prod(1-\mathsf q^n)^{24}$ of power series with
integer coefficients gives $\Delta$; $j=c_4^3/\Delta$ follows, and its coefficients are those of the $q$-expansion of the
modular function $j$ (computed with PARI). As $v(\Delta)=v(\mathsf q)$ is finite, $\Delta\ne0$.

(ii) $c_4\in1+\mathfrak m$ is a unit and $v(\Delta)=v(\mathsf q)$, so the model is minimal (\Cref{prop:reduction:min-criterion}(i)) with
multiplicative reduction, of type $\kod{I}_n$, $n=v(\Delta)$. The reduced curve is $y^2+xy=x^3$ with the node at $(0,0)$, and the
tangent cone $y^2+xy=y(y+x)$ has the slopes $0$ and $-1$ in $k$: split. Finally $v(j)=3v(c_4)-v(\Delta)=-n$.

(iii) Over $\C$, with $\mathsf q=e^{2\pi i\tau}$ and $w=e^{2\pi iz}$, one has $X=(2\pi i)^{-2}\wp(z;\Z+\Z\tau)-\frac1{12}$ and
$Y=(2\pi i)^{-3}\tfrac12\wp'(z)-\tfrac12X$, so the identities $Y^2+XY=X^3+a_4X+a_6$ and the addition law hold as identities of
formal series with rational coefficients (we also checked the first numerically to $40$ digits). Tate's theorem is that
they hold $p$-adically and that $\phi$ is bijective; see \cite[V.3]{SilvermanATAEC}. We do not prove it.
\end{proof}

\begin{proposition}[inverting $j$]\label{prop:reduction:j-inverse}
There is a unique power series
\[
\mathsf q(J)=J+744J^2+750420J^3+872769632J^4+1102652742882J^5+\cdots\in J\,\Z[[J]]
\]
such that $j(\mathsf q(J))=1/J$ as formal series. For $j_0\in K$ with $v(j_0)<0$, the element $\mathsf q_0=\mathsf q(1/j_0)$ is the unique
$\mathsf q_0\in K^\times$ with $v(\mathsf q_0)>0$ and $j(E_{\mathsf q_0})=j_0$, and $v(\mathsf q_0)=-v(j_0)$.
\end{proposition}

\begin{proof}
Write $j(\mathsf q)=\mathsf q^{-1}U(\mathsf q)$ with $U\in1+\mathsf q\Z[[\mathsf q]]$. Then $J(\mathsf q):=1/j(\mathsf q)=\mathsf q\,U(\mathsf q)^{-1}\in\mathsf q+\mathsf q^2\Z[[\mathsf q]]$. A power
series $J(\mathsf q)=\mathsf q+\sum_{n\ge2}b_n\mathsf q^n$ with $b_n\in\Z$ has a compositional inverse $\mathsf q(J)=J+\sum_{n\ge2}c_nJ^n$ with $c_n\in\Z$: the
coefficient of $J^n$ in $J(\mathsf q(J))$ is $c_n$ plus a polynomial with integer coefficients in $c_2,\dots,c_{n-1}$ and the $b_i$,
so the equation $J(\mathsf q(J))=J$ determines the $c_n$ recursively in $\Z$. The coefficients printed were computed with PARI's
\texttt{serreverse}. Since both series have integer coefficients, substituting an element of $\mathfrak m$ gives convergent
series, and substitution respects composition. Put $\mathsf q_0=\mathsf q(1/j_0)$; then $v(\mathsf q_0)=v(1/j_0)=-v(j_0)>0$ (the leading term
dominates), and $J(\mathsf q_0)=1/j_0$, so $j(E_{\mathsf q_0})=j_0$. If also $v(\mathsf q_1)>0$ and $j(E_{\mathsf q_1})=j_0$, then
$\mathsf q_1=\mathsf q(J(\mathsf q_1))=\mathsf q(1/j_0)=\mathsf q_0$.
\end{proof}

\begin{theorem}[curves with $v(j)<0$]\label{thm:reduction:tate-uniformisation}
Let $E/K$ have $v(j)<0$, let $q_E=\mathsf q(1/j_E)$ be as in \Cref{prop:reduction:j-inverse}, and let
$\gamma=-c_4/c_6\in K^\times/(K^\times)^2$, computed from any model of $E$. Then $E$ is the quadratic twist of $E_{q_E}$ by
$\gamma$: $E\simeq E_{q_E}$ over $K(\sqrt\gamma)$, and over $K$ if and only if $\gamma\in(K^\times)^2$. Moreover:
\begin{enumerate}[label=(\roman*)]
\item $E$ has split multiplicative reduction if and only if $\gamma\in(K^\times)^2$;
\item $E$ has nonsplit multiplicative reduction if and only if $K(\sqrt\gamma)/K$ is the unramified quadratic extension;
\item $E$ has additive reduction if and only if $K(\sqrt\gamma)/K$ is ramified.
\end{enumerate}
In the multiplicative case $n=v(\Delta_{\min})=-v(j)=v(q_E)$.
\end{theorem}

\begin{proof}[Sketch of proof]
As $v(j)<0$, $j\ne0,1728$, so $c_4c_6\ne0$. For $E_{\mathsf q}$, $\gamma(E_{\mathsf q})=(1+240s_3)/(1-504s_5)=1+24(10s_3+21s_5)/(1-504s_5)$ lies in
$1+24\mathfrak m$, which consists of squares by Hensel's lemma (for $p$ odd, $1+\mathfrak m$ consists of squares; for $p=2$, $x^2=1+8m$ with
$v(8m)>v(4)$ has a root near $1$). Two curves with the same $j\ne0,1728$ are twists of each other by
$d=c_6(E)c_4(E_{\mathsf q})/(c_6(E_{\mathsf q})c_4(E))=\gamma(E_{\mathsf q})/\gamma(E)$ (\Cref{ch:ellcurves}), and $d\equiv\gamma(E)$ modulo squares. This proves
the first statement. If $\gamma$ is a square, $E\simeq E_{q_E}$ has split multiplicative reduction by
\Cref{thm:reduction:tate-curve}(ii). If $K(\sqrt\gamma)/K$ is unramified, $E$ acquires multiplicative reduction over it, hence has
multiplicative reduction (\Cref{prop:reduction:etale-base-change}); for $p$ odd, computing $\gamma$ from a minimal model makes it a unit
($c_4$ and $c_6$ are units), which is not a square, so by Hensel's lemma its reduction is not a square, and
\Cref{prop:reduction:split}(ii) shows the reduction is nonsplit. If $E$ has multiplicative reduction and $p$ is
odd, $c_4$ and $c_6$ are units for a minimal model, so $K(\sqrt\gamma)$ is unramified; hence a ramified $\gamma$ gives additive reduction.
This proves (i) to (iii) for $p$ odd; for $p=2$ see \cite[V.5]{SilvermanATAEC}. The last sentence is
\Cref{thm:reduction:potential-criterion}(iii) and \Cref{prop:reduction:j-inverse}.
\end{proof}

\begin{definition}[Tate parameter and unit part]\label{def:reduction:tate-parameter}
Let $E/\Q_q$ have $v_q(j_E)<0$. The \emph{Tate parameter} of $E$ is $q_E=\mathsf q(1/j_E)\in\Q_q^\times$, and
\[
\nu=v_q(q_E)=-v_q(j_E)\ge1 ;
\]
at a multiplicative prime $\nu=n_q=v_q(\Delta_E)$ is the paper's $n$. The \emph{unit part} is $u=q_Eq^{-\nu}\in\Z_q^\times$, with
reduction $\bar u\in\F_q^\times$; for a prime $\ell\ne q$, $c(u)$ denotes the class of $\bar u$ in
$\F_q^\times/(\F_q^\times)^{\ell^{v_\ell(q-1)}}$ (the paper's $c$ in Lemma lem:tamfactor; the arithmetic of this class is in
\Cref{ch:rings}).
\end{definition}

\begin{proposition}[invariance of the Tate parameter]\label{prop:reduction:tate-parameter-invariance}
Let $E/\Q_q$ have $v_q(j_E)<0$.
\begin{enumerate}[label=(\roman*)]
\item $q_E$, $\nu$, $u$ and $c(u)$ depend only on $j_E$; in particular they are the same for $E$ and for every quadratic twist of
$E$, and $t_\ell(E,q)=v_\ell(\nu)$ is invariant under quadratic twist.
\item If $\varphi\colon E\to E'$ is an isogeny over $\Q_q$ of prime degree $d$, then $\nu(E')\in\{d\,\nu(E),\ \nu(E)/d\}$. Consequently $v_\ell(\nu)$
and $t_\ell(E,q)$ are unchanged by isogenies of degree prime to $\ell$.
\end{enumerate}
\end{proposition}

\begin{proof}
(i) is the definition: $q_E$ is a function of $j_E$, and twisting does not change $j$.
(ii) \emph{Sketch.} Over a finite extension $L$ of $\Q_q$ over which $E$ and $E'$ are both split, $E(\bar L)=\bar L^\times/q_E^{\Z}$, and
$\ker\varphi$ is a subgroup $C$ of order $d$. If $C=\mu_d$, then $w\mapsto w^d$ identifies $\bar L^\times/\mu_dq_E^{\Z}$ with
$\bar L^\times/q_E^{d\Z}$, so $E'\simeq E_{q_E^d}$ and $\nu(E')=d\nu(E)$. Otherwise $C$ is generated by the class of $\zeta q_E^{1/d}$ for some
$\zeta\in\mu_d$, and $\bar L^\times/\langle\zeta q_E^{1/d}\rangle$ is the Tate curve with parameter $\zeta q_E^{1/d}$, so $\nu(E')=\nu(E)/d$. That
an isogeny of Tate curves is induced by the corresponding map of multiplicative groups is part of Tate's theory
\cite[V.3, V.5]{SilvermanATAEC}. For the consequence: since $d\ne\ell$, $v_\ell(d\nu)=v_\ell(\nu/d)=v_\ell(\nu)$. (A
second proof: an isogeny of degree prime to $\ell$ induces an isomorphism $T_\ell E\simeq T_\ell E'$ of Galois modules, and
by Lemma lem:tamagawa-exp (proved in \Cref{ch:galrep}) $v_\ell(\nu)$ is the largest $t$ such that $E[\ell^t]$ is unramified after twisting
$E$ by the quadratic character of \Cref{thm:reduction:tate-uniformisation}, which commutes with isogenies.)
\end{proof}

\begin{proposition}[torsion of the Tate curve]\label{prop:reduction:tate-torsion}
Let $\mathsf q\in K^\times$ with $v(\mathsf q)>0$, $\ell$ a prime and $t\ge1$.
\begin{enumerate}[label=(\roman*)]
\item In $E_{\mathsf q}(\bar K)=\bar K^\times/\mathsf q^{\Z}$, fix an $\ell^t$-th root $\mathsf q^{1/\ell^t}$ and a primitive $\ell^t$-th root of unity
$\zeta$. Then $E_{\mathsf q}[\ell^t]=\{\zeta^a(\mathsf q^{1/\ell^t})^b:a,b\in\Z/\ell^t\}$, and there is an exact sequence of $G_K$-modules
\[
0\longrightarrow\mu_{\ell^t}\longrightarrow E_{\mathsf q}[\ell^t]\longrightarrow\Z/\ell^t\longrightarrow0,
\]
the quotient being generated by the class of $\mathsf q^{1/\ell^t}$.
\item Let $\ell\ne p$, choose compatible roots $\zeta_{\ell^k}$ and $\mathsf q^{1/\ell^k}$, and let $e_1=(\zeta_{\ell^k})_k$, $e_2=(\mathsf q^{1/\ell^k})_k$ be the
corresponding basis of $T_\ell E_{\mathsf q}$. Then $\sigma\in G_K$ acts by
\[
\begin{pmatrix}\chi_\ell(\sigma)&\kappa(\sigma)\\0&1\end{pmatrix},
\]
where $\kappa$ is the Kummer cocycle of $\mathsf q$: $\sigma(\mathsf q^{1/\ell^k})=\zeta_{\ell^k}^{\kappa(\sigma)}\mathsf q^{1/\ell^k}$ for all $k$ (\Cref{ch:galois}).
\item If $E$ is the twist of $E_{\mathsf q}$ by a quadratic character $\varepsilon$ of $G_K$ (\Cref{thm:reduction:tate-uniformisation}), then $G_K$
acts on $T_\ell E$ by
\[
\begin{pmatrix}\chi_\ell\varepsilon&\varepsilon\kappa\\0&\varepsilon\end{pmatrix}.
\]
\end{enumerate}
\end{proposition}

\begin{proof}
(i) A class $w\in\bar K^\times/\mathsf q^{\Z}$ is killed by $\ell^t$ if and only if $w^{\ell^t}=\mathsf q^b$ for some $b\in\Z$, that is,
$w=\zeta^a\mathsf q^{b/\ell^t}$ with $\mathsf q^{b/\ell^t}=(\mathsf q^{1/\ell^t})^b$. Changing $b$ by $\ell^t$ multiplies $w$ by $\mathsf q$, so $b$ matters only modulo
$\ell^t$, and $w$ is trivial if and only if $a\equiv b\equiv0$. The map $w\mapsto b$ is a homomorphism onto $\Z/\ell^t$ with kernel
$\mu_{\ell^t}$, and it commutes with $G_K$ because $\sigma(\mathsf q^{1/\ell^t})$ is again an $\ell^t$-th root of $\mathsf q$. By
\Cref{thm:reduction:tate-curve}(iii) these are the $\ell^t$-torsion points of $E_{\mathsf q}$, with their Galois action.
(ii) In additive notation $\sigma(e_1)=\chi_\ell(\sigma)e_1$ by definition of $\chi_\ell$, and $\sigma(e_2)=e_2+\kappa(\sigma)e_1$ by definition of
$\kappa$; the columns of the matrix are the images of $e_1$ and $e_2$.
(iii) $T_\ell E\simeq T_\ell E_{\mathsf q}\otimes\varepsilon$ as $G_K$-modules (\Cref{ch:ellcurves}): every matrix of (ii) is multiplied by
$\varepsilon(\sigma)$. In particular the upper right entry is $\varepsilon\kappa$, which satisfies the cocycle relation
$c(gh)=\chi_\ell\varepsilon(g)c(h)+c(g)\varepsilon(h)$ for these matrices; $\kappa$ alone does not when $\varepsilon\ne1$.
\end{proof}

\begin{warning}\label{warn:reduction:epskappa}
In the nonsplit and the additive potentially multiplicative cases the corner of the matrix is $\varepsilon\kappa$, not $\kappa$. The
proof of Lemma lem:tamfactor in the original version of the paper wrote $\kappa$; the revised version writes $\varepsilon\kappa$. In the nonsplit
multiplicative case (the setting of Lemma lem:tamfactor) $\varepsilon$ is unramified, so nothing in the conclusion changes: $\varepsilon$ is
trivial on inertia and only changes the Frobenius entry by a sign. In the additive potentially multiplicative case $\varepsilon$ is the
ramified character of \Cref{thm:reduction:tate-uniformisation}(iii), so it is not trivial on inertia; there the paper does not
use the matrix of $T$ but applies the lemma to $E_{q_E}$, since $\ad T$ is unchanged by a quadratic twist (\Cref{ch:blochkato}).
\end{warning}

\begin{paperlink}
Lemma lem:tamagawa-exp: ``The $\ell^t$-torsion is generated by $\mu_{\ell^t}$ and an $\ell^t$-th root of $q_E$'' is
\Cref{prop:reduction:tate-torsion}(i), applied over the unramified quadratic extension, where a nonsplit curve becomes a Tate
curve (\Cref{thm:reduction:tate-uniformisation}) and the inertia group is unchanged. The proof of the lemma (inertia acts trivially on
$E[\ell^t]$ if and only if $\ell^t\mid n_q$) is given in \Cref{ch:galrep}. Lemma lem:tamfactor uses the matrices of (iii), the
decomposition $q_E=q^\nu u$, and the restriction of $\kappa$ to inertia, which is $\nu$ times the tame character (\Cref{ch:local}); the
cohomology is done in \Cref{ch:cohomology}, \Cref{ch:galcoh} and \Cref{ch:blochkato}.
\end{paperlink}

\begin{proposition}[Tamagawa numbers at multiplicative primes]\label{prop:reduction:tate-tamagawa}
Let $E/K$ have multiplicative reduction of type $\kod{I}_n$, $k$ finite.
\begin{enumerate}[label=(\roman*)]
\item If the reduction is split, $E(K)\simeq K^\times/q_E^{\Z}$, $E_0(K)$ corresponds to $\mathcal O^\times q_E^{\Z}/q_E^{\Z}$, and the valuation
induces $E(K)/E_0(K)\simeq\Z/n$; so $c=n$.
\item In both cases $\Phig\simeq\Z/n$. If the reduction is nonsplit, Frobenius acts on $\Phig$ by $-1$, so $c=\#\Phi_E(k)$ is $1$ for $n$
odd and $2$ for $n$ even.
\end{enumerate}
\end{proposition}

\begin{proof}
(i) By \Cref{thm:reduction:tate-uniformisation}, $E\simeq E_{q_E}$, and $E(K)\simeq K^\times/q_E^{\Z}$ by \Cref{thm:reduction:tate-curve}(iii). Every
class has a representative $w$ with $0\le v(w)<n$, since multiplying by $q_E^{\pm1}$ changes $v(w)$ by $\pm n$. Write $X(w,q_E)=\sum_{i\in\Z}X_i-2s_1(q_E)$ and $Y(w,q_E)=\sum_iY_i+s_1(q_E)$
termwise. For $i=-m\le-1$ one has $X_i=q_E^mw/(q_E^m-w)^2$ and $Y_i=q_E^mw^2/(q_E^m-w)^3$, and $s_1(q_E)\in\mathfrak m$. If $v(w)=0$ and
$\bar w\ne1$, every term other than $X_0$, $Y_0$ lies in $\mathfrak m$, so $\phi(w)$ reduces to
$\bigl(\bar w/(1-\bar w)^2,\ \bar w^2/(1-\bar w)^3\bigr)$, a point of $y^2+xy=x^3$ other than the node $(0,0)$. (Directly:
$y^2+xy=\bar w^4/(1-\bar w)^6+\bar w^3/(1-\bar w)^5=\bar w^3\bigl(\bar w+1-\bar w\bigr)/(1-\bar w)^6=x^3$. Its slope is
$\lambda=y/x=\bar w/(1-\bar w)$; the tangents at the node are $y=0$ and $y=-x$, and the isomorphism of \Cref{prop:reduction:Ens}(ii)
with $\alpha=0$, $\beta=-1$ sends the point to $\lambda/(\lambda+1)=\bar w$. So on units the reduction map is $w\mapsto\bar w\in k^\times$.) If $\bar w=1$, $X$ has a pole and
$\phi(w)$ reduces to $\tilde O$. If $0<v(w)<n$, then $v(X_0)=v(w)$, $v(Y_0)=2v(w)$, $v(X_{-1})=v(Y_{-1})=n-v(w)$, and all the other terms have even larger
valuation; all are positive, so $\phi(w)$ reduces to the node. Hence $E_0(K)$ is the image of $\mathcal O^\times q_E^{\Z}$, and
$w\mapsto v(w)\bmod n$ gives $E(K)/E_0(K)\simeq\Z/n$.

(ii) Let $K'$ be the unramified quadratic extension, over which $E$ is split (\Cref{prop:reduction:split}(iii) for $p$ odd,
\Cref{thm:reduction:tate-uniformisation} in general). By (i) over every finite unramified extension of $K'$ the component group has
$\Z/n$ points, with trivial Galois action, so $\Phig\simeq\Z/n$ (\Cref{prop:reduction:etale-base-change}). If $E$ is nonsplit it is the
twist of $E_{q_E}$ by the unramified quadratic character, and the nontrivial element $\sigma$ of $\Gal(K'/K)$ acts on
$E(K')\simeq K'^\times/q_E^{\Z}$ by $w\mapsto\sigma(w)^{-1}$ (the twisted action, \Cref{ch:ellcurves}). On
$E(K')/E_0(K')\simeq K'^\times/\mathcal O'^\times q_E^{\Z}\simeq\Z/n$ this is $a\mapsto-a$. The invariants of $-1$ on $\Z/n$ are $\{a:2a=0\}$, of order
$\gcd(2,n)$.
\end{proof}

All $153{,}343$ pairs (curve, prime of multiplicative reduction) of Cremona's tables below conductor $10{,}000$ obey this rule.

\begin{theorem}[types at additive potentially multiplicative primes]\label{thm:reduction:potmult-additive}
Let $E/K$ have additive reduction with $v(j)<0$, and put $\nu=-v(j)\ge1$.
\begin{enumerate}[label=(\roman*)]
\item If $p$ is odd, $E$ has type $\kod{I}_\nu^*$, $f=2$ and $v(\Delta_{\min})=\nu+6$. Conversely, a curve of type $\kod{I}_n^*$ with $n\ge1$ at an
odd $p$ has $v(j)<0$, hence $n=-v(j)$.
\item If $p=2$ and $K=\Q_2$, $E$ has type $\kod{I}_m^*$ with $m=\nu+2f_2-4$, where $f_2=2a(\chi)\in\{4,6\}$ and $a(\chi)\in\{2,3\}$ is the conductor
exponent of the ramified quadratic character $\chi$ of \Cref{thm:reduction:tate-uniformisation}; so $m=\nu+4$ or $m=\nu+8$.
\end{enumerate}
\end{theorem}

\begin{proof}[Sketch of proof]
(i) By \Cref{thm:reduction:tate-uniformisation}, $\gamma$ is in a ramified class, so $\gamma=\varpi w$ with $w$ a unit up to squares, and
$E\simeq E''^{(\varpi)}$ where $E''=E_{q_E}^{(w)}$ has multiplicative reduction (its $\gamma$ is the unramified class of $w$) with
$v(\Delta''_{\min})=\nu$. Take a minimal model of $E''$ of the form $y^2=x^3+a_2x^2+a_4x+a_6$ (complete the square; $p$ is odd) and
translate the node to $(0,0)$, so that $\varpi\mid a_4,a_6$ and $\varpi\nmid a_2$ (the reduced cubic is $x^2(x+\tilde a_2)$ with a double, not
triple, root). The twist by $\varpi$ is $y^2=x^3+\varpi a_2x^2+\varpi^2a_4x+\varpi^3a_6$. As in the proof of
\Cref{prop:reduction:I0star-cubic}, Tate's algorithm passes Steps 2 to 5 and reaches Step 6 with
$P(T)=T^3+a_2T^2+a_4T+a_6\equiv T^2(T+\tilde a_2)$, which has a double and a simple root: the type is $\kod{I}_m^*$ for some $m\ge1$, and
the model is minimal, with $v(\Delta)=\nu+6$. By \Cref{prop:reduction:conductor-values}(iv), $f=2$, and Ogg's formula
(\Cref{thm:reduction:ogg}) gives $m+5=v(\Delta)+1-f=\nu+5$, so $m=\nu$.

For the converse, let $E$ have type $\kod{I}_n^*$, $n\ge1$, at odd $p$. As in the converse part of the proof of
\Cref{prop:reduction:I0star-cubic}, there is a minimal model $y^2=x^3+a_2x^2+a_4x+a_6$ with $\varpi\mid a_2$, $\varpi^2\mid a_4$, $\varpi^3\mid a_6$
whose cubic $\tilde P$ has a double and a simple root. The model $y^2=x^3+(a_2/\varpi)x^2+(a_4/\varpi^2)x+a_6/\varpi^3$ of the twist by $\varpi$
is integral and reduces to $y^2=\tilde P(x)$, which has a node. So the twist has multiplicative reduction, and
$v(j)=v(j(E^{(\varpi)}))<0$.

(ii) We do not prove this. It is a special case of Lorenzini's description of the minimal regular model of a quadratic twist of a
curve with multiplicative reduction \cite{Lorenzini10} (we follow the citation in \cite{Praveen26} and have not checked the
numbering there); $f_2=2a(\chi)$ is the conductor exponent of the twist of the Steinberg representation by $\chi$
(\Cref{ch:local}), and $a(\chi)=2$ for $\Q_2(\sqrt{-1})$, $\Q_2(\sqrt{3})$ and $3$ for $\Q_2(\sqrt{\pm2})$, $\Q_2(\sqrt{\pm10})$.
\end{proof}

On Cremona's tables below conductor $10{,}000$, all $17{,}761$ additive pairs with $v_p(j)<0$ have a symbol $\kod{I}_m^*$
($3{,}732$ at $p=2$, $7{,}805$ at $p=3$, the rest at $p\ge5$); at odd $p$ always $m=\nu$; at $p=2$, $m-\nu=4$ with $f_2=4$ in
$2{,}742$ cases and $m-\nu=8$ with $f_2=6$ in $990$ cases. Examples at $2$: \curve{208a1} ($\kod{I}_5^*$, $\nu=1$, $f_2=4$),
\curve{304b1} ($\kod{I}_7^*$, $\nu=3$, $f_2=4$), \curve{832c1} ($\kod{I}_9^*$, $\nu=1$, $f_2=6$), \curve{336d1} ($\kod{I}_{12}^*$, $\nu=8$,
$f_2=4$).

\begin{example}[Tate parameters of the paper's curves]\label{ex:reduction:tate-parameters}
The Tate parameters $q_E=q^\nu u$ at the potentially multiplicative primes of the running examples, computed from the series of
\Cref{prop:reduction:j-inverse} (and, at odd $q$, also with PARI's \texttt{ellinit(E, O(q\^{}k)).tate}, with the same result):
\begin{center}\small
\begin{tabular}{llrrrl}
\toprule
curve, prime & type & $a_q$ & $\nu$ & $\bar u\bmod q$ & remark\\
\midrule
\curve{11a1}, $11$ & $\kod{I}_5$ & $1$ & $5$ & $10\equiv-1$ & $(\F_{11}^\times)^5=\{\pm1\}$: $\bar u$ is a $5$th power\\
\curve{11a3}, $11$ & $\kod{I}_1$ & $1$ & $1$ & $8$ & \\
\curve{20a1}, $5$ & $\kod{I}_2$ & $-1$ & $2$ & $1$ & nonsplit\\
\curve{77c1}, $7$ & $\kod{I}_3$ & $-1$ & $3$ & $5$ & same $q_E$ as \curve{539c1} at $7$ ($\kod{I}_3^*$)\\
\curve{77c1} and \curve{539c1}, $11$ & $\kod{I}_2$ & $1$ & $2$ & $3$ & \\
\curve{140b1}, $7$ & $\kod{I}_5$ & $-1$ & $5$ & $1$ & same $q_E$ as \curve{980i1} at $7$ ($\kod{I}_5^*$)\\
\curve{140b1}, $5$ & $\kod{I}_1$ & $-1$ & $1$ & $2$ & nonsplit\\
\curve{980i1}, $5$ & $\kod{I}_1$ & $1$ & $1$ & $2$ & split: the twist by $-7$, $\chi_{-7}(5)=-1$\\
\curve{726a1}, $2$ & $\kod{I}_6$ & $-1$ & $6$ & $1$ & \\
\curve{726a1}, $3$ & $\kod{I}_3$ & $-1$ & $3$ & $2$ & \\
\curve{121d1}, $11$ & $\kod{I}_1^*$ & $0$ & $1$ & $8$ & same $q_E$ as \curve{11a3}\\
\curve{121d2}, $11$ & $\kod{I}_5^*$ & $0$ & $5$ & $10$ & same $q_E$ as \curve{11a1}\\
\bottomrule
\end{tabular}
\end{center}
For instance $q_{E}=10\cdot11^5+2\cdot11^6+6\cdot11^7+\cdots$ for \curve{11a1}. The values of $t_\ell$: $t_5(\curve{11a1},11)=1$;
$t_3(\curve{539c1},7)=t_3(\curve{77c1},7)=v_3(3)=1$; $t_5(\curve{980i1},7)=1$; $t_3(\curve{726a1},2)=v_3(6)=1$. At $q=3$, \curve{726a1} has
$v_3(\nu)=1$, which the appendix's \texttt{rule} counts as $t_3(E,3)$; the revised Definition def:terms defines $t_\ell(E,q)$ for every
bad $q$, including $q=\ell$.
\end{example}

\begin{warning}[PARI's Tate parameter at $2$]\label{warn:reduction:pari-tate}
In PARI/GP 2.17.2, the value \texttt{ellinit(E, O(2\^{}k)).tate[3]} does not return the Tate parameter when $v_2(j)=-1$: in all $1{,}262$ such pairs
among the curves of conductor below $3{,}000$ it returned $q_E^5$ (for instance $2^5+2^6+\cdots$ for \curve{208a1}, whose Tate parameter
is $2+2^2+2^4+\cdots$), while for every other pair with $v_p(j)<0$ in that range ($\nu\ge2$ at $2$, and all odd $p$) it agreed with the
series of \Cref{prop:reduction:j-inverse}. Before using PARI's value, check that $j(\mathsf q)=j_E$, or compute $q_E$ from the series.
The paper's computations use $v_q(j)$ directly and are not affected.
\end{warning}

\section{Quadratic twists and reduction}\label{sec:reduction:twists}

The paper twists constantly: by $\chi_{q^*}$ at a prime of type $\kod{I}_0^*$ or $\kod{I}_n^*$ (Theorems thm:qual (b), (c) and
thm:quant, Lemma lem:euler), by a product $\chi_S$ of such characters (Theorems thm:twist and thm:quant), and in the
definition of the twist-symmetrised weight $w_\ell(E,q)$ of its Section 1.4. This section proves the paper's Lemma
lem:twist (i) and (ii), which say what twisting does to Kodaira types, component groups and conductor exponents; part
(iii) of that lemma (Atkin and Li) is about newforms and is stated in \Cref{ch:hecke}, with a sketch via Galois conductors
(only the case $q^2\mid M$ is proved there). Recall (\Cref{ch:ellcurves}) that
for $d\in\Q^\times$ the twist $E^{(d)}$ of a model $y^2=x^3+a_2x^2+a_4x+a_6$ is $y^2=x^3+da_2x^2+d^2a_4x+d^3a_6$, with invariants
$d^2c_4$, $d^3c_6$, $d^6\Delta$ and the same $j$, and that $\rho_{E^{(d)},\ell}\simeq\rho_{E,\ell}\otimes\chi$, where $\chi$ is the quadratic character of $\Q(\sqrt d)$; it is
written $\chi_D$, indexed by the discriminant $D$ of $\Q(\sqrt d)$ (for $d=11$, $D=44$ and $\chi=\chi_{-4}\chi_{-11}$). For an odd prime $q$, the
paper's $E^{(q)}$ is the twist by $q^*=(-1)^{(q-1)/2}q$, that is, by $\chi_{q^*}$ and the field $\Q(\sqrt{q^*})$, which is unramified
outside $q$ (\Cref{ch:ant}).

\begin{lemma}[unramified twists; the paper's Lemma lem:twist (i)]\label{lem:reduction:twist-unramified}
Let $E/\Q$, $p$ a prime and $d\in\Q^\times$ such that $\chi=\chi_D$, the quadratic character of $\Q(\sqrt d)$ ($D$ its discriminant), is unramified at $p$ (for instance $\chi_{q^*}$ with $q\ne p$ an odd
prime; at $p=2$ this uses $q^*\equiv1\pmod4$). Then $E$ and $E^{(d)}$ have at $p$ the same reduction type, the same Kodaira
symbol, the same geometric component group $\Phig_p$, the same $v_p(\Delta_{\min})$, the same conductor exponent $f_p$, the same
$v_p(j)$ and the same Tate parameter. The Tamagawa number $c_p$ can change.
\end{lemma}

\begin{proof}
$E$ and $E^{(d)}$ become isomorphic over $\Q_p(\sqrt d)$, which is $\Q_p$ or the unramified quadratic extension of $\Q_p$. By
\Cref{prop:reduction:etale-base-change}, the reduction type, the Kodaira symbol, $\Phig_p$ and $v_p(\Delta_{\min})$ are unchanged by this
extension, so they agree for $E$ and $E^{(d)}$. The conductor exponents agree because $V_\ell E^{(d)}\simeq V_\ell E\otimes\chi$ and $\chi$
is trivial on $I_p$, so the two representations of $I_p$ coincide. The $j$-invariant, hence $v_p(j)$ and the Tate parameter, does
not change under twisting. An example where $c_p$ changes: \curve{11a1} has split multiplicative reduction of type $\kod{I}_5$ at
$11$, with $c_{11}=5$; its twist by $13$ (and $13$ is not a square modulo $11$, so $\chi_{13}$ is unramified at $11$ with
$\chi_{13}(\Frob_{11})=-1$) is \curve{1859a2}, of type $\kod{I}_5$ nonsplit at $11$, with $c_{11}=1$; so is its twist by $-3$,
\curve{99d2}.
\end{proof}

\begin{lemma}[ramified twists at odd primes; the paper's Lemma lem:twist (ii)]\label{lem:reduction:twist-ramified}
Let $q$ be an odd prime and $E^{(q)}=E^{(q^*)}$.
\begin{enumerate}[label=(\alph*)]
\item If $E$ has type $\kod{I}_n$ at $q$ with $n\ge1$, then $E^{(q)}$ has type $\kod{I}_n^*$ at $q$, with the same $n$, and conversely.
\item If $E$ has type $\kod{I}_0$ at $q$, then $E^{(q)}$ has type $\kod{I}_0^*$, and conversely.
\item If $q\ge5$, the twist exchanges $\kod{II}\leftrightarrow\kod{IV}^*$, $\kod{III}\leftrightarrow\kod{III}^*$, $\kod{IV}\leftrightarrow\kod{II}^*$,
and $v_q(\Delta_{\min})$ changes by $+6$ modulo $12$.
\item The conductor exponent $f_q(E^{(q)})$ is $0$ if $E^{(q)}$ has good reduction, $1$ if it has type $\kod{I}_n$ with $n\ge1$, and $2$
if it has additive reduction and is one of the curves in (a), (b), (c).
\end{enumerate}
At every prime $p\ne q$ the two curves have the same type, $\Phig_p$ and $f_p$ (\Cref{lem:reduction:twist-unramified}).
\end{lemma}

\begin{proof}
Throughout, $q$ is odd and we use the twisted model of a model $y^2=x^3+a_2x^2+a_4x+a_6$ (complete the square first; for $q$ odd
this does not change minimality at $q$). Its invariants are $(q^*)^2c_4$, $(q^*)^3c_6$, $(q^*)^6\Delta$: the twist by $q^*$, not
by $q$, though only $q$-adic valuations are used below, and those are the same for $q$ and $q^*$.

(a) Here $v_q(j)<0$. By \Cref{thm:reduction:tate-uniformisation} the reduction type of a curve with $v_q(j)<0$ is determined by
the square class of $\gamma=-c_4/c_6$: unramified for multiplicative, ramified for additive reduction. The twist multiplies $\gamma$
by $(q^*)^{2}/(q^*)^{3}=1/q^*$, which moves an unramified class to a ramified one and back: for odd $q$, $\Q_q(\sqrt\gamma)/\Q_q$
is trivial or unramified exactly when $v_q(\gamma)$ is even (then $\gamma$ is a unit times a square), and ramified when it is odd,
and $v_q(1/q^*)=-1$ changes the parity. So $E$ is multiplicative if and only if
$E^{(q)}$ is additive. By \Cref{thm:reduction:potmult-additive}(i) the additive one has type $\kod{I}_\nu^*$, and the multiplicative
one type $\kod{I}_\nu$, with the same $\nu=-v_q(j)$, as $j$ is unchanged. (The identification of the index of $\kod{I}_\nu^*$
with $\nu$ uses Ogg's formula, through \Cref{thm:reduction:potmult-additive}.)

(b) If $E$ has good reduction at $q$, \Cref{prop:reduction:I0star-cubic} shows that $E^{(q)}$ has type $\kod{I}_0^*$ (with the sign
$\epsilon=q^*/q$ in the cubic). Conversely, if $E$ has type $\kod{I}_0^*$, \Cref{prop:reduction:I0star-cubic} gives a model
$y^2=x^3+a_2x^2+a_4x+a_6$ with $q\mid a_2$, $q^2\mid a_4$, $q^3\mid a_6$ and $\tilde P$ separable. The twist by $q^*$ is
$y^2=x^3+q^*a_2x^2+(q^*)^2a_4x+(q^*)^3a_6$, and the change $(q,0,0,0)$ turns it into
$y^2=x^3+\epsilon a_{2,1}x^2+a_{4,2}x+\epsilon a_{6,3}$ with $\epsilon=q^*/q$ and $a_{i,j}=a_i/q^j$. This model is integral, and its cubic
has the roots $\epsilon\theta$ for the roots $\theta$ of $P$, which are distinct modulo $q$; so its discriminant is a unit and $E^{(q)}$ has good
reduction.

(c) Let $q\ge5$ and let $E$ be potentially good and additive, with minimal $v=v_q(\Delta_{\min})\in\{2,3,4,6,8,9,10\}$
(\Cref{cor:reduction:tame-table}). The twisted model has $v_q(\Delta)=v+6$ and $v_q(c_4)$ raised by $2$. If $v+6<12$ it is minimal
(\Cref{prop:reduction:min-criterion}(i)). If $v+6\ge12$, then, as $v_q(j)\ge0$ is unchanged, $3v_q(c_4')\ge v_q(\Delta')\ge12$ for the twisted
model, so $v_q(c_4')\ge4$ and one step of \Cref{prop:reduction:min-criterion}(ii) gives a minimal model with $v_q(\Delta)=v-6<12$. So the
new minimal valuation is $v+6$ or $v-6$, that is, $v+6\bmod12$: $2\leftrightarrow8$, $3\leftrightarrow9$, $4\leftrightarrow10$ (and $0\leftrightarrow6$, which is (b)).
By \Cref{cor:reduction:tame-table}(iii) this is $\kod{II}\leftrightarrow\kod{IV}^*$, $\kod{III}\leftrightarrow\kod{III}^*$,
$\kod{IV}\leftrightarrow\kod{II}^*$.

(d) The values $0$ and $1$ are \Cref{prop:reduction:conductor-values}(i), (ii). In the additive cases of (a), (b), (c) we have
$V_\ell E^{(q)}=V_\ell E\otimes\chi_{q^*}$, where $\chi_{q^*}$ is tamely ramified at $q$ (it is trivial on the pro-$q$ group $P_q$, as $q$ is
odd), and $V_\ell E$ is tame at $q$: it is unramified (case (b)), multiplicative (case (a), \Cref{prop:reduction:conductor-values}(ii)), or
additive with $q\ge5$ (case (c), \Cref{prop:reduction:conductor-values}(iv)). So $P_q$ acts trivially on $V_\ell E^{(q)}$,
$\operatorname{sw}=0$, and $f_q(E^{(q)})=2$ by \Cref{prop:reduction:conductor-values}(iii).
\end{proof}

All $90{,}970$ pairs (curve numbered $1$ in its isogeny class, odd bad prime $q$) of Cremona's tables below conductor $10{,}000$
behave as stated: the same $n$ in all $78{,}351$ cases of (a), the shift of $v_q(\Delta_{\min})$ by $6$ modulo $12$ in all $5{,}264$
potentially good additive cases with $q\ge5$ other than $\kod{I}_0^*$, and the same type, conductor exponent and $v_p(\Delta_{\min})$ at
every other bad prime.

\begin{warning}\label{warn:reduction:q-versus-qstar}
The paper's proof of Lemma lem:twist (ii) in its original version wrote the invariants of the twist as $q^2c_4$, $q^3c_6$, $q^6\Delta$,
which are those of the twist by $q$; the twist by $q^*$ has $(q^*)^2c_4$, $(q^*)^3c_6$, $(q^*)^6\Delta$. For $q\equiv3\pmod4$ the two are
different curves, differing by the twist by $-1$, which is ramified at $2$: the twist of \curve{121a1} by $-11$ is \curve{121c2}, of
conductor $121$, while its twist by $11$ is \curve{1936k2}, of conductor $2^4\cdot11^2$ (type $\kod{I}_4^*$ at $2$). Likewise the phrase
``the conductor exponents are the standard ones for $q$ odd (tame)'' is right for the exchanges of (a), (b), (c), and wrong for
general additive reduction at $q=3$, which can be wild (\Cref{rem:reduction:wild-at-3}).
\end{warning}

\begin{proposition}[the exchanges at $q=3$]\label{prop:reduction:twist-at-3}
Let $E/\Q$ have additive, potentially good reduction at $3$ of type $\kod{II}$, $\kod{III}$, $\kod{IV}$, $\kod{IV}^*$, $\kod{III}^*$ or $\kod{II}^*$, and
let $E^{(3)}=E^{(-3)}$. Then $f_3(E^{(3)})=f_3(E)$, and the type of $E^{(3)}$ at $3$ is $\kod{IV}^*$, $\kod{III}^*$, $\kod{II}^*$, $\kod{II}$, $\kod{III}$,
$\kod{IV}$ respectively.
\end{proposition}

\begin{proof}
The conductor: $V_\ell E^{(3)}=V_\ell E\otimes\chi_{-3}$ with $\chi_{-3}$ tamely ramified at $3$, so $P_3$ acts on both through the same
representation and the Swan conductors agree; $E^{(3)}$ is additive: it is not multiplicative, as $v_3(j)\ge0$, and it does not have good
reduction, since otherwise $E=(E^{(3)})^{(3)}$ would have type $\kod{I}_0^*$ by \Cref{lem:reduction:twist-ramified}(b). So
$V^{I_3}=0$ for both and $f_3(E^{(3)})=f_3(E)=:f$.
The discriminants: by the formula of \Cref{ch:ellcurves} for the twist of a general model, the twisted model has discriminant
$2^{12}(-3)^6\Delta$, so $v_3(\Delta_{\min}(E^{(3)}))\equiv v_3(\Delta_{\min}(E))+6\pmod{12}$.
The components: both curves are potentially good ($j$ is unchanged and $v_3(j)\ge0$), so neither has type $\kod{I}_n$ or, by the converse
in \Cref{thm:reduction:potmult-additive}(i), $\kod{I}_n^*$ with $n\ge1$; so both have $1\le m\le9$, and Ogg's formula
(\Cref{thm:reduction:ogg}) gives $v_3(\Delta_{\min})=m+f-1\in[f,f+8]$ for both. Two integers in an interval of length $8$ that are
congruent modulo $12$ after adding $6$ to one of them differ by exactly $6$; so $m(E^{(3)})=m(E)\pm6$. The types with $m\in\{1,2,3\}$ and
additive potentially good reduction are $\kod{II}$, $\kod{III}$, $\kod{IV}$, and those with $m\in\{7,8,9\}$ are $\kod{IV}^*$, $\kod{III}^*$,
$\kod{II}^*$. Since $m(E)\in\{1,2,3,7,8,9\}$ and $1\le m(E^{(3)})=m(E)\pm6\le9$, $m(E^{(3)})=m(E)+6$ if $m(E)\le3$ and $m(E)-6$ if $m(E)\ge7$. This
gives the exchanges.
\end{proof}

This answers, given Ogg's formula, the question left open in the revised paper (Open question 6) whether the exchanges of
Lemma lem:twist (ii) hold at $q=3$. All $5{,}122$ such pairs below conductor $10{,}000$ (curves numbered $1$) behave as stated. Examples:
\curve{54b1} (type $\kod{II}$, $f_3=3$) and its twist \curve{54a1} (type $\kod{IV}^*$, $f_3=3$); \curve{27a1} ($\kod{IV}^*$) and \curve{27a3}
($\kod{II}$); \curve{36a1} ($\kod{III}$, $f_3=2$) and \curve{36a3} ($\kod{III}^*$).

\begin{paperlink}
Section 1.4 of the paper lists the exchanges ``the first and the last for every odd $q$, the three others for $q\ge5$'';
\Cref{prop:reduction:twist-at-3} shows that the three others also hold at $q=3$. Lemma lem:euler (``$V=V'\otimes\chi_{q^*}$ with $V'$
Steinberg at $q$'' for $\kod{I}_n^*$, and unramified for $\kod{I}_0^*$) uses (a) and (b). Theorems thm:qual (b), (c) and thm:quant
twist a type $\kod{II}$, $\kod{II}^*$, $\kod{I}_0^*$ or $\kod{I}_n^*$ prime to a type $\kod{IV}^*$, $\kod{IV}$, good or multiplicative one; Theorem
thm:quant needs this at every odd $q$ in $S$, which (a) and (b) provide. The proof of Theorem thm:main uses the multiplicative twist at an
$\kod{I}_n^*$ prime; at $q=2$, where \Cref{lem:reduction:twist-ramified} does not apply, the local statement
\Cref{thm:reduction:tate-uniformisation} (a ramified quadratic twist of a Tate curve) is what is needed.
\end{paperlink}

\begin{proposition}[the twist-symmetrised weight]\label{prop:reduction:twist-weight}
For a bad prime $q$ and an odd prime $\ell$, let $w_\ell(E,q)$ be the paper's weight: $v_\ell(n)$ if the type at $q$ is $\kod{I}_n$, or
$\kod{I}_n^*$ with $q$ odd; $1$ if $\ell=3$ and the type is $\kod{IV}$ or $\kod{IV}^*$, or $\kod{II}$ or $\kod{II}^*$ with $q$ odd; $0$ otherwise.
Then
\[
w_\ell(E,q)=\max\bigl(v_\ell(\#\Phig_q(E)),\,v_\ell(\#\Phig_q(E^{(q)}))\bigr)\ \text{for odd }q,\qquad w_\ell(E,2)=v_\ell(\#\Phig_2(E)).
\]
\end{proposition}

\begin{proof}
For odd $q$ use the exchanges of \Cref{lem:reduction:twist-ramified} and \Cref{prop:reduction:twist-at-3} and the orders of
\Cref{tab:reduction:kodaira}. These exchanges pair up all types at an odd $q$ ($\kod{I}_n^*$ with $\kod{I}_n$ for $n\ge0$, and the six remaining
additive types among themselves), so the type of $E$ at a bad prime $q$ lies in exactly one of the pairs below, in one order or the
other. The maximum is symmetric in $E$ and $E^{(q)}$, and so is the definition of $w_\ell$ on each pair with both members bad (it
gives $v_\ell(n)$ to $\kod{I}_n$ and $\kod{I}_n^*$, and $[\ell=3]$ to $\kod{II}$, $\kod{IV}^*$, $\kod{IV}$, $\kod{II}^*$); for the pair
$\kod{I}_0\leftrightarrow\kod{I}_0^*$ only the member of type $\kod{I}_0^*$ is bad at $q$, with $w_\ell=0$. $\kod{I}_n\leftrightarrow\kod{I}_n^*$: $\max(v_\ell(n),v_\ell(4))=v_\ell(n)$, as $\ell$ is odd.
$\kod{I}_0\leftrightarrow\kod{I}_0^*$: $\max(0,v_\ell(4))=0$. $\kod{II}\leftrightarrow\kod{IV}^*$ and $\kod{IV}\leftrightarrow\kod{II}^*$:
$\max(v_\ell(1),v_\ell(3))=[\ell=3]$. $\kod{III}\leftrightarrow\kod{III}^*$: $v_\ell(2)=0$. At $q=2$: $\kod{I}_n$ gives $v_\ell(n)$,
$\kod{I}_n^*$ and $\kod{I}_0^*$ give $v_\ell(4)=0$, $\kod{IV}$ and $\kod{IV}^*$ give $v_\ell(3)=[\ell=3]$, the others $0$, which is the definition at
$q=2$.
\end{proof}

The paper states the max formula for $q\ge5$ only; by \Cref{prop:reduction:twist-at-3} it holds at $q=3$ as well.

\begin{example}[twist pairs]\label{ex:reduction:twist-pairs}
Curve-level twists used in the paper, each identified by computing the twisted model, its minimal model and its label (all
checked with PARI against Cremona's tables):
\begin{center}\small
\begin{tabular}{llrlll}
\toprule
$E$ & type at $q$ & $d$ & $E^{(d)}$ & type at $q$ & other bad primes\\
\midrule
\curve{11a1} & $\kod{I}_5$ at $11$ & $-11$ & \curve{121d2} & $\kod{I}_5^*$ ($v=11$) & none\\
\curve{11a3} & $\kod{I}_1$ at $11$ & $-11$ & \curve{121d1} & $\kod{I}_1^*$ ($v=7$) & none\\
\curve{121a1} & $\kod{II}$ ($v=2$) at $11$ & $-11$ & \curve{121c2} & $\kod{IV}^*$ ($v=8$) & none\\
\curve{121c1} & $\kod{IV}$ ($v=4$) at $11$ & $-11$ & \curve{121a2} & $\kod{II}^*$ ($v=10$) & none\\
\curve{539c1} & $\kod{I}_3^*$ at $7$ & $-7$ & \curve{77c1} & $\kod{I}_3$ & $\kod{I}_2$ at $11$ on both\\
\curve{539d1} & $\kod{I}_0^*$ at $7$ & $-7$ & \curve{11a3} & $\kod{I}_0$ & $\kod{I}_1$ at $11$ on both\\
\curve{980i1} & $\kod{I}_5^*$ at $7$ & $-7$ & \curve{140b1} & $\kod{I}_5$ & $\kod{IV}^*$ at $2$, $\kod{I}_1$ at $5$ on both\\
\curve{20a1} & $\kod{IV}^*$ at $2$ & $-4$ & \curve{80b1} & $\kod{I}_0^*$, $f_2=4$ & $\kod{I}_2$ at $5$ on both\\
\bottomrule
\end{tabular}
\end{center}
In the row of \curve{980i1}, the type $\kod{IV}^*$ at $2$ is preserved because $-7\equiv1\pmod8$, so $\chi_{-7}$ is trivial on $G_{\Q_2}$. The last
row is a wild twist at $2$, to which no rule of \Cref{lem:reduction:twist-ramified} applies. The first row also shows that
$e_\ell(E,q)$ is not twist invariant (the revised paper corrects the last sentence of Definition def:terms accordingly; see \Cref{ch:paper}): \curve{11a1} has $e_3(E,11)=0$, while \curve{121d2} has $e_3(E,11)=v_3(11^2-1)=1$
(\Cref{ch:lfunctions}); the Tate parameters at $11$ coincide (\Cref{ex:reduction:tate-parameters}).
\end{example}

\section{The wild primes \texorpdfstring{$2$ and $3$}{2 and 3}}\label{sec:reduction:wild}

At $p\ge5$ every additive reduction is tame and the Kodaira symbol determines $v(\Delta_{\min})$, $f$ and the image of inertia
(\Cref{cor:reduction:tame-table}, \Cref{thm:reduction:kraus}). At $p=2$ and $p=3$ none of this holds. The paper meets this in Remark
rem:wild (``At $q=2,3$ the Kodaira type does not determine $P_q$''), in Section sec:wild (the minima at $2$ and the weights $w_3$
at $3$), and in Definition def:terms (the index of a symbol $\kod{I}_n^*$ at $2$ is not the valuation of the Tate parameter). This
section collects what changes and gives the dictionary between the image of inertia at $2$ and the factorisation of the
$3$-division polynomial over $\Q_2$, which is how the image is computed in practice.

\begin{remark}[what changes at $2$ and $3$]\label{rem:reduction:wild-overview}
Let $p\in\{2,3\}$.
\begin{enumerate}[label=(\roman*)]
\item The type does not determine $v(\Delta_{\min})$ or $f$, although Ogg's formula $f=v(\Delta_{\min})+1-m$ still links them
(\Cref{thm:reduction:ogg}).
\item $f$ can exceed $2$: up to $8$ at $2$ and $5$ at $3$ (\Cref{thm:reduction:conductor-bounds}).
\item $G_{\mathrm{ss}}$ can be non-abelian, and of order divisible by $p$ (\Cref{thm:reduction:kraus}(ii), (iii)).
\item The type does not determine $G_{\mathrm{ss}}$, with two exceptions: $\kod{IV}$, $\kod{IV}^*$ at $2$ force $G_{\mathrm{ss}}\simeq\Z/3$
(\Cref{cor:reduction:IV-at-2}), and $\kod{III}$, $\kod{III}^*$, $\kod{I}_0^*$ at $3$ force $G_{\mathrm{ss}}\simeq\Z/4$, $\Z/4$, $\Z/2$
(\Cref{rem:reduction:wild-at-3}).
\item At $2$ a symbol $\kod{I}_m^*$ may be potentially good, and when it is potentially multiplicative $m\ne-v_2(j)$
(\Cref{thm:reduction:potmult-additive}(ii), \Cref{rem:reduction:Instar-at-2}).
\end{enumerate}
Hence the Euler factor of the adjoint is not a function of the type at $2$ and $3$, as the paper says in Remark rem:wild; what
determines it is the image of inertia (\Cref{ch:lfunctions}; the paper's Proposition prop:shapes).
\end{remark}

\begin{proposition}[inertia at $2$ and the $3$-division polynomial]\label{prop:reduction:inertia-at-2}
Let $E/\Q_2$ have additive, potentially good reduction. Let $x_1,\dots,x_4$ be the roots of the $3$-division polynomial $\psi_3$ (the
$x$-coordinates of the four lines of $E[3]$, \Cref{ch:ellcurves}) and $e_i$ the ramification index of $\Q_2(x_i)/\Q_2$. Then
$G_{\mathrm{ss}}$ is one of $\Z/2$, $\Z/3$, $\Z/4$, $\Z/6$, $Q_8$, $\SL_2(\F_3)$, and the multiset $(e_1,\dots,e_4)$ is
\begin{itemize}
\item $(1,1,1,1)$ if $G_{\mathrm{ss}}\simeq\Z/2$; then $E$ is the twist of a curve with good reduction at $2$ by a ramified quadratic
character $\chi$, and $f_2=2a(\chi)\in\{4,6\}$;
\item $(1,3,3,3)$ if $G_{\mathrm{ss}}$ is cyclic of order $3$ or $6$; order $3$ occurs exactly for the types $\kod{IV}$, $\kod{IV}^*$, with
$f_2=2$, and order $6$ is the twist of the order-$3$ case by a ramified quadratic character, with $f_2\ge3$;
\item $(2,2,2,2)$ if $G_{\mathrm{ss}}\simeq\Z/4$;
\item $(4,4,4,4)$ if $G_{\mathrm{ss}}\simeq Q_8$ or $\SL_2(\F_3)$.
\end{itemize}
\end{proposition}

\begin{proof}
The list of groups is \Cref{thm:reduction:kraus}(iii) ($G_{\mathrm{ss}}\ne1$ as the reduction is additive). By
\Cref{thm:reduction:serre-tate}, $G_{\mathrm{ss}}$ is the image of $I_2$ in $\Aut(E[3])$, a subgroup of $\SL_2(\F_3)$, and the
$x_i$ correspond to the four lines of $E[3]$. For a root $x_i$, $e_i=[\Q_2^{\ur}(x_i):\Q_2^{\ur}]$ (the residue degree is absorbed by
$\Q_2^{\ur}$), which is the size of the orbit of $x_i$, that is of its line, under $I_2$. So $(e_i)$ is the multiset of orbit sizes of
$G_{\mathrm{ss}}$ on the four lines, each orbit counted once for each of its lines. The element $-1$ acts trivially on lines, and
$\SL_2(\F_3)/\{\pm1\}\simeq A_4$ acts on the four lines as $A_4$ on four points. The unique element of order $2$ of $\SL_2(\F_3)$ is
$-1$; elements of order $3$ and $6$ map to $3$-cycles; elements of order $4$ map to double transpositions; $Q_8$ maps onto the
Klein group, which is transitive; and $\SL_2(\F_3)$ maps onto $A_4$. This gives the four patterns. If $G_{\mathrm{ss}}=\{\pm1\}$, the
character $\psi\colon I_2\to\{\pm1\}$ through which inertia acts is invariant under conjugation by $G_{\Q_2}$ (a conjugate of an element of $I_2$ acts
by the conjugate matrix, and $-1$ is central), and it extends to a quadratic character $\chi$ of $G_{\Q_2}$. Explicitly: choose a
Frobenius lift $\phi\in G_{\Q_2}$. Its image generates $G_{\Q_2}/I_2\simeq\hat\Z$ topologically, and since $\hat\Z$ is free as a
profinite group on one generator, $a\mapsto\phi^a$ is a section, so $G_{\Q_2}=I_2\rtimes\overline{\langle\phi\rangle}$ with
$\overline{\langle\phi\rangle}\simeq\hat\Z$ and every element written uniquely as $i\phi^a$ ($i\in I_2$, $a\in\hat\Z$). Put
$\chi(i\phi^a)=\psi(i)$. Then $\chi(i\phi^a\cdot j\phi^b)=\chi\bigl(i(\phi^aj\phi^{-a})\phi^{a+b}\bigr)=\psi(i)\psi(\phi^aj\phi^{-a})=\psi(i)\psi(j)$
by the invariance of $\psi$, so $\chi$ is a continuous character extending $\psi$ (this is the construction of the paper's
Proposition prop:shapes(ii)); then
$\rho\otimes\chi$ is unramified, so $E\otimes\chi$ has good reduction (criterion of N\'eron, Ogg and Shafarevich), and
$f_2(E)=a(V'\otimes\chi)=2a(\chi)$ for $V'$ unramified (\Cref{ch:local}), with $a(\chi)\in\{2,3\}$. The same argument applied to the
character $g\mapsto g^3\in\{\pm1\}$ of a cyclic $G_{\mathrm{ss}}$ of order $6$ shows that $E$ is the twist of a curve with
$G_{\mathrm{ss}}\simeq\Z/3$; order $6$ is wild since $2\mid6$, so $f_2\ge3$. Order $3$ is tame, so $f_2=2$, and by
\Cref{thm:reduction:tame-additive} and \Cref{cor:reduction:IV-at-2} this happens exactly for $\kod{IV}$, $\kod{IV}^*$.
\end{proof}

\Cref{tab:reduction:inertia2} records what occurs on Cremona's tables below conductor $10{,}000$: all $20{,}979$ pairs (curve, $p=2$)
with additive potentially good reduction at $2$, by type, $f_2$ and pattern. The patterns were computed by factoring $\psi_3$ over
$\Q$ and reading the ramification indices of the primes above $2$ in each factor with PARI.

\begin{table}[ht]
\centering\small
\begin{tabular}{lp{11cm}}
\toprule
type at $2$ & $f_2$: pattern (number of pairs)\\
\midrule
$\kod{IV}$ & $2$: $(1,3,3,3)$ $2{,}094$\\
$\kod{IV}^*$ & $2$: $(1,3,3,3)$ $2{,}656$\\
$\kod{II}$ & $4$: $(1,3,3,3)$ $469$, $(4,4,4,4)$ $478$; $6$: $(1,1,1,1)$ $136$, $(4,4,4,4)$ $395$; $7$: $(4,4,4,4)$ $236$\\
$\kod{II}^*$ & $3$: $(4,4,4,4)$ $1{,}171$; $4$: $(1,1,1,1)$ $291$; $6$: $(1,3,3,3)$ $108$, $(4,4,4,4)$ $96$\\
$\kod{III}$ & $3$: $(4,4,4,4)$ $1{,}062$; $5$: $(4,4,4,4)$ $644$; $7$: $(4,4,4,4)$ $210$; $8$: $(2,2,2,2)$ $98$, $(4,4,4,4)$ $96$\\
$\kod{III}^*$ & $3$: $(4,4,4,4)$ $1{,}582$; $5$: $(4,4,4,4)$ $314$; $7$: $(4,4,4,4)$ $210$; $8$: $(2,2,2,2)$ $98$, $(4,4,4,4)$ $96$\\
$\kod{I}_0^*$ & $4$: $(1,3,3,3)$ $596$, $(4,4,4,4)$ $720$; $5$: $(4,4,4,4)$ $792$; $6$: $(1,3,3,3)$ $196$, $(4,4,4,4)$ $200$\\
$\kod{I}_m^*$ (potentially good) & $3$: $(4,4,4,4)$ $1{,}575$; $4$: $(1,1,1,1)$ $466$, $(4,4,4,4)$ $1{,}305$; $5$: $(4,4,4,4)$ $528$; $6$: $(1,1,1,1)$ $250$,
$(1,3,3,3)$ $146$, $(4,4,4,4)$ $1{,}429$; $7$: $(4,4,4,4)$ $236$\\
\bottomrule
\end{tabular}
\caption{Additive potentially good reduction at $2$ below conductor $10{,}000$: conductor exponent and ramification pattern of the
roots of $\psi_3$ over $\Q_2$. The pattern $(1,3,3,3)$ with $f_2=2$ is $G_{\mathrm{ss}}\simeq\Z/3$, with $f_2\ge4$ it is $\Z/6$.}
\label{tab:reduction:inertia2}
\end{table}

In particular $f_2=2$ occurs only for $\kod{IV}$ and $\kod{IV}^*$, in accordance with \Cref{thm:reduction:tame-additive}.

\begin{paperlink}
The Euler factor $P_2(X)$ (\Cref{ch:lfunctions}; the paper's Proposition prop:shapes) is determined by $G_{\mathrm{ss}}$ and the action of
Frobenius on $(\ad V)^{I_2}$: cyclic of order $3$ or $6$ gives $P_2(X)=1+2X$ and $e_3(E,2)=1$ (the order-$6$ case reduces to the order-$3$
case, as $\ad V$ is unchanged by a quadratic twist); $\Z/2$ gives the cubic $(1-2X)(1-(a^2-4)X+4X^2)$ of the twist $E'$ with good reduction,
$a=a_2(E')$, and $e_3(E,2)=v_3(9-a^2)$; $\Z/4$ gives $1\pm2X$; $Q_8$ and $\SL_2(\F_3)$ act absolutely irreducibly, so $(\ad V)^{I_2}=0$ and
$P_2=1$. So a type $\kod{II}$ or $\kod{II}^*$ at $2$ can have $e_3(E,2)\in\{0,1,2\}$: \curve{336f1} has type $\kod{II}$ at $2$, $f_2=4$, pattern
$(1,3,3,3)$, $P_2=1+2X$, $e_3(E,2)=1$ and $\degphi=24$; \curve{1232l1}, the twist of \curve{77a1} by $\chi_{-4}$ with $a_2(\curve{77a1})=0$, has
type $\kod{II}^*$ at $2$, $f_2=4$, pattern $(1,1,1,1)$, $P_2=(1-2X)(1+2X)^2=1+2X-4X^2-8X^3$, $e_3(E,2)=2$ and $\degphi=288=2^5\cdot3^2$. The
original version of Remark rem:wild said that ``a type $\kod{II}$ at $2$ has $P_2=1$ and $e_3=0$''; the revised version says ``usually'',
with these two curves as counterexamples. A potentially good $\kod{I}_m^*$ at $2$ has the pattern $(4,4,4,4)$, hence $P_2=1$, in most cases
(\Cref{tab:reduction:inertia2}).
\end{paperlink}

\begin{remark}[the prime $3$]\label{rem:reduction:wild-at-3}
Let $E/\Q_3$ have additive reduction.
\begin{enumerate}[label=(\roman*)]
\item Types $\kod{III}$, $\kod{III}^*$, $\kod{I}_0^*$ are tame ($f_3=2$), with $G_{\mathrm{ss}}$ cyclic of order $4$, $4$, $2$.
\item Types $\kod{II}$, $\kod{II}^*$, $\kod{IV}$, $\kod{IV}^*$ are wild: $f_3\in\{3,4,5\}$ and $3\mid e$.
\item Type $\kod{I}_n^*$ ($n\ge1$) is potentially multiplicative, with $n=-v_3(j)$ and $f_3=2$.
\end{enumerate}
\end{remark}

\begin{proof}[Justification]
(iii) is \Cref{thm:reduction:potmult-additive}(i) and \Cref{prop:reduction:conductor-values}(iv). The tame/wild dichotomy in (i), (ii) is
\Cref{thm:reduction:tame-additive}. In (ii), wild means that the pro-$3$ group $P_3$ acts nontrivially, so its image, a nontrivial
$3$-group, lies in $G_{\mathrm{ss}}$. In (i), $G_{\mathrm{ss}}$ is a quotient of the procyclic tame inertia of order prime to $3$, so it is
cyclic of order $2$ or $4$ (\Cref{thm:reduction:serre-tate}(iii)). It has order $2$ if and only if $E$ is a ramified quadratic twist of a
curve with good reduction (the argument of \Cref{prop:reduction:inertia-at-2}), that is, up to an unramified twist, the twist by $\chi_{-3}$
of such a curve, which has type $\kod{I}_0^*$ by \Cref{lem:reduction:twist-ramified}(b) and \Cref{lem:reduction:twist-unramified};
conversely type $\kod{I}_0^*$ is such a twist. So $\kod{III}$, $\kod{III}^*$ have order $4$.
\end{proof}

Below conductor $10{,}000$, at $3$: $\kod{III}$, $\kod{III}^*$ ($1{,}553$ pairs each), $\kod{I}_0^*$ ($2{,}283$) and $\kod{I}_n^*$ ($7{,}805$) always have
$f_3=2$; $\kod{II}$ and $\kod{IV}^*$ both have $f_3=3,4,5$ in $879$, $280$, $107$ pairs, and $\kod{IV}$ and $\kod{II}^*$ both in $628$, $177$, $71$
pairs. The equal distributions are \Cref{prop:reduction:twist-at-3}: the twist by $\chi_{-3}$ is unramified outside $3$, preserves the
conductor, and exchanges these types.

\begin{paperlink}
This is why Lemma lem:IV excludes $q=3$ (at $3$ the types $\kod{IV}$, $\kod{IV}^*$ are wild, and the conclusion ``tamely ramified, Serre
conductor exponent $1$'' fails), and why the weights $w_3$ at $q=\ell=3$ in Conjecture conj:euler are not explained by Lemma lem:euler. The
Remark after the Example of Section 3 of the paper mentions $q=3$ in cases (a) and (b) of Theorem thm:qual; there the obstruction is
$q=\ell$, and the ramification at $3$ of the types $\kod{II}$, $\kod{IV}$ is also wild by (ii).
\end{paperlink}

\begin{remark}[symbols $\kod{I}_m^*$ at $2$]\label{rem:reduction:Instar-at-2}
At $2$ a symbol $\kod{I}_m^*$ is potentially good in $5{,}935$ of the $9{,}667$ pairs below conductor $10{,}000$ ($61\%$; among curves
numbered $1$ in their class, which below conductor $10{,}000$ are the optimal curves except in class \curve{990h} (whose optimal
curve is \curve{990h3}; it is multiplicative at $2$, as $2\parallel990$, and so does not enter this count), $2{,}340$ of $4{,}037$, that is $58\%$, the
figure quoted in the revised Remark rem:wild; the six curves with complex multiplication among them are all potentially good, so the
optimal non-CM curves give $2{,}334$ of $4{,}031$, again $58\%$). When it is
potentially multiplicative, $t_\ell(E,2)=v_\ell(\nu)$ with $\nu=-v_2(j)=m-4$ or $m-8$ (\Cref{thm:reduction:potmult-additive}(ii)), so
$v_\ell(m)$ is the wrong invariant: \curve{304b1} has type $\kod{I}_7^*$ and $\nu=3$, so $t_3(E,2)=1$ although $v_3(7)=0$; \curve{832c1} has
type $\kod{I}_9^*$ and $\nu=1$, so $t_3(E,2)=0$ although $v_3(9)=2$.
\end{remark}

\begin{paperlink}
Definition def:terms: ``a symbol $\kod{I}_n^*$ at $q=2$ can have potentially good reduction, and then $t_\ell=0$''; Section sec:wild: ``the
index $n$ of the symbol is not the valuation of the Tate parameter at $2$, which is why $n$ was not the right invariant''. In the proof of
Theorem thm:main a potentially multiplicative prime $2$ is handled by \Cref{thm:reduction:tate-uniformisation} (a ramified quadratic
twist of a Tate curve with $\nu=-v_2(j)$), not by Lemma lem:twist, which is for odd $q$ (the original proof of Theorem thm:main applied Lemma lem:twist there;
see \Cref{ch:paper}).
\end{paperlink}

\section{Cremona's tables and PARI}\label{sec:reduction:tables}

The paper's computations take the curves from Cremona's tables and recompute every local invariant with PARI (Section
sec:data and the appendix). This short section says what the tables contain and what PARI computes, and collects the local
data of the running examples in one table. File formats and the practical use of PARI are in \Cref{ch:computation}.

\begin{remark}[what the tables contain]\label{rem:reduction:cremona-tables}
Cremona's tables \cite{Cremona97,ecdata} list, for every elliptic curve over $\Q$ up to their conductor bound, the global minimal model
in the normalisation of \Cref{thm:reduction:global-minimal}, the conductor, the isogeny class, rank and torsion. The file
\texttt{allbsd} adds the product $\prod_pc_p$ of the Tamagawa numbers, the real period and the other quantities of the conjecture of Birch
and Swinnerton-Dyer; other files give the isogeny matrices, the modular degrees (\texttt{degphi}) and the mod $\ell$ images. Section
sec:data of the paper describes the \texttt{ecdata} repository as containing every curve ``with its Kodaira symbols, Tamagawa
numbers, modular degree, isogeny class, torsion'' and mod $\ell$ images, and adds that ``the Kodaira types and discriminant exponents
at the bad primes were recomputed with PARI/GP 2.15 from the $a$-invariants'' (its appendix calls \texttt{ellglobalred}). The files
used in this book (\texttt{allcurves}, \texttt{allbsd}, \texttt{allisog}, \texttt{alllabels}, \texttt{degphi}, \texttt{galrep},
conductor below $10{,}000$) are a subset of that repository and, unlike the paper's description of the full repository, do not list
the Kodaira symbols or the Tamagawa numbers prime by prime: here these are recomputed with Tate's algorithm (\texttt{ellglobalred}),
as in the paper's own check. File formats are in \Cref{ch:computation}. The minimal discriminant and $j$ follow from the model.
\end{remark}

\begin{example}[a PARI session]\label{ex:reduction:pari-session}
With PARI/GP 2.17.2, for \curve{726a1}, \curve{20a1} and \curve{11a1}:
\begin{lstlisting}[language=GP]
E = ellinit([1,1,0,-35,-51]);        \\ 726a1
R = ellglobalred(E)
  \\ [726, [1, 0, 0, 0], 4, [2, 1; 3, 1; 11, 2], [[1, 10, 0, 2], [1, 7, 0, 1], [2, 3, 0, 2]]]
R[4]                                  \\ [2, 1; 3, 1; 11, 2]
elllocalred(ellinit([0,1,0,4,4]), 2)  \\ 20a1: [2, -4, [1, 0, 0, 0], 3]
elllocalred(ellinit([1,0,1,170,-3237]), 7)   \\ 539c1: [2, -7, [1, 0, 0, 0], 4]
E11 = ellinit([0,-1,1,-10,-20]);     \\ 11a1
ellinit(E11, O(11^8)).tate[3]
  \\ 10*11^5 + 2*11^6 + 6*11^7 + 6*11^8 + 5*11^9 + ... + O(11^13)
valuation(E11.j, 11)                  \\ -5
\end{lstlisting}
The local data of \curve{726a1} are $[f,k,\ast,c]=[1,10,\ast,2]$, $[1,7,\ast,1]$, $[2,3,\ast,2]$ at $2,3,11$: types $\kod{I}_6$, $\kod{I}_3$, $\kod{III}$
(codes $10=6+4$, $7=3+4$, $3$), Tamagawa numbers $2,1,2$ with product $4$, and conductor $2\cdot3\cdot11^2$. For \curve{20a1} at $2$ the code
$-4$ is $\kod{IV}^*$, $f_2=2$, $c_2=3$. For \curve{539c1} at $7$ the code $-7=-(3+4)$ is $\kod{I}_3^*$. For \curve{11a1} the Tate parameter has
$\nu=5$ and $\bar u=10$ (\Cref{ex:reduction:tate-parameters}).
\end{example}

\begin{example}[local data of the running examples]\label{ex:reduction:running-table}
For each curve: conductor, minimal discriminant, and at each bad prime $q$ the type, $f_q$, $c_q$, $\#\Phig_q$ and $v_q(j)$ (all computed
with PARI; the hand computations of \Cref{ex:reduction:tate-running} agree).
\begin{center}\footnotesize
\begin{tabular}{lllp{7.4cm}}
\toprule
curve & $N$ & $\Delta_{\min}$ & $q$: type, $f_q$, $c_q$, $\#\Phig_q$, $v_q(j)$\\
\midrule
\curve{11a1} & $11$ & $-11^5$ & $11$: $\kod{I}_5$, $1$, $5$, $5$, $-5$\\
\curve{20a1} & $2^2\cdot5$ & $-2^8\cdot5^2$ & $2$: $\kod{IV}^*$, $2$, $3$, $3$, $4$; $5$: $\kod{I}_2$, $1$, $2$, $2$, $-2$\\
\curve{99a1} & $3^2\cdot11$ & $3^3\cdot11$ & $3$: $\kod{III}$, $2$, $2$, $2$, $9$; $11$: $\kod{I}_1$, $1$, $1$, $1$, $-1$\\
\curve{121a1} & $11^2$ & $-11^2$ & $11$: $\kod{II}$, $2$, $1$, $1$, $1$\\
\curve{121c1} & $11^2$ & $-11^4$ & $11$: $\kod{IV}$, $2$, $1$, $3$, $2$\\
\curve{121c2} & $11^2$ & $-11^8$ & $11$: $\kod{IV}^*$, $2$, $1$, $3$, $1$\\
\curve{121d1} & $11^2$ & $-11^7$ & $11$: $\kod{I}_1^*$, $2$, $2$, $4$, $-1$\\
\curve{539c1} & $7^2\cdot11$ & $-7^9\cdot11^2$ & $7$: $\kod{I}_3^*$, $2$, $4$, $4$, $-3$; $11$: $\kod{I}_2$, $1$, $2$, $2$, $-2$\\
\curve{539d1} & $7^2\cdot11$ & $-7^6\cdot11$ & $7$: $\kod{I}_0^*$, $2$, $2$, $4$, $0$; $11$: $\kod{I}_1$, $1$, $1$, $1$, $-1$\\
\curve{726a1} & $2\cdot3\cdot11^2$ & $2^6\cdot3^3\cdot11^3$ & $2$: $\kod{I}_6$, $1$, $2$, $6$, $-6$; $3$: $\kod{I}_3$, $1$, $1$, $3$, $-3$; $11$: $\kod{III}$, $2$, $2$,
$2$, $0$\\
\curve{980i1} & $2^2\cdot5\cdot7^2$ & $-2^8\cdot5\cdot7^{11}$ & $2$: $\kod{IV}^*$, $2$, $1$, $3$, $19$; $5$: $\kod{I}_1$, $1$, $1$, $1$, $-1$; $7$: $\kod{I}_5^*$, $2$,
$2$, $4$, $-5$\\
\bottomrule
\end{tabular}
\end{center}
The curves \curve{121c1}, \curve{121c2} at $11$ and \curve{980i1} at $2$ have $c_q=1<3=\#\Phig_q$; \curve{726a1} at $2$ and $3$ has $c_q<\#\Phig_q$
because the reduction is nonsplit.
\end{example}

\section{Exercises}\label{sec:reduction:exercises}

The exercises go from routine to harder. Several ask the reader to check a claim of the paper; PARI is useful for the
numerical ones (\Cref{ch:computation}).

\begin{exercise}\label{exer:reduction:minimal}
Show that an integral model with $v(\Delta)<12$ is minimal, and that for $p\ge5$ an integral model with $v(c_4)\ge4$ and
$v(\Delta)\ge12$ is not (\Cref{prop:reduction:min-criterion}). Decide whether the model $y^2=x^3-27c_4x-54c_6$ of \curve{11a1} is minimal at
$2$ and at $3$.
\end{exercise}

\begin{exercise}\label{exer:reduction:split}
For \curve{11a1} at $11$ and \curve{20a1} at $5$, decide whether the reduction is split, once from the polynomial $T^2+a_1T-a_2$ after moving
the node to $(0,0)$, and once from the Legendre symbol $\bigl(\frac{-c_6}{p}\bigr)$ (\Cref{prop:reduction:split}). Compare with $a_{11}=1$ and
$a_5=-1$.
\end{exercise}

\begin{exercise}\label{exer:reduction:Ens}
Write down explicitly the isomorphisms $\tilde E_{\mathrm{ns}}(k)\simeq k^+$ for $y^2=x^3$ and $\tilde E_{\mathrm{ns}}(k)\simeq k^\times$ for
$y^2=x^3+x^2$, and check directly that three collinear points are sent to elements with sum $0$, respectively product $1$
(\Cref{prop:reduction:Ens}).
\end{exercise}

\begin{exercise}\label{exer:reduction:torsion}
Let $E=\curve{11a1}$, which has good reduction at $3$ and $7$. Using \Cref{cor:reduction:torsion-injects}, show that the prime-to-$3$
part of $E(\Q)_{\mathrm{tors}}$ injects into $\tilde E(\F_3)$ and the prime-to-$7$ part into $\tilde E(\F_7)$. Count these groups (orders $5$
and $10$), deduce that $E(\Q)_{\mathrm{tors}}$ has order dividing $5$, and conclude that $E(\Q)_{\mathrm{tors}}\simeq\Z/5$ by exhibiting a point
of order $5$ (for instance $(5,5)$ on $y^2+y=x^3-x^2-10x-20$).
\end{exercise}

\begin{exercise}\label{exer:reduction:tate-hand}
Run Tate's algorithm by hand for \curve{121a1} and \curve{121c2} at $11$ and for \curve{539d1} and \curve{539c1} at $7$, and check the types
used in the Example of Section 3 of the paper (\Cref{ex:reduction:tate-running}).
\end{exercise}

\begin{exercise}\label{exer:reduction:twist121}
Check the paper's claim (Lemma lem:twist (ii)) that the twist of \curve{121a1} by $-11$ has type $\kod{IV}^*$ at $11$. Then show that its twist
by $+11$ has additive reduction at $2$ (it is \curve{1936k2}, of type $\kod{I}_4^*$ at $2$ with $v_2(j)=0$), so that the twist by $q$ and the
twist by $q^*$ really are different curves (\Cref{warn:reduction:q-versus-qstar}).
\end{exercise}

\begin{exercise}\label{exer:reduction:cubic}
Check the paper's sentence ``the cubic $T^3+a_2T^2+a_4T+a_6$ of Tate's algorithm has distinct roots modulo $q$'' for the twist of
\curve{11a3} by $-7$ (it is \curve{539d1}): compute a model $y^2=x^3+a_2x^2+a_4x+a_6$ of \curve{11a3}, the cubic of the twist with the sign
$\epsilon=-1$, and its roots modulo $7$ (\Cref{prop:reduction:I0star-cubic}).
\end{exercise}

\begin{exercise}\label{exer:reduction:weight}
Prove $w_\ell(E,q)=\max\bigl(v_\ell(\#\Phig_q(E)),v_\ell(\#\Phig_q(E^{(q)}))\bigr)$ for $q\ge5$ from \Cref{tab:reduction:kodaira} and
\Cref{lem:reduction:twist-ramified} (\Cref{prop:reduction:twist-weight}). Which curves of \Cref{ex:reduction:running-table} have
$c_q\ne\#\Phig_q$, and why?
\end{exercise}

\begin{exercise}\label{exer:reduction:ogg}
Verify Ogg's formula $f_q=v_q(\Delta_{\min})+1-m_q$ at every bad prime of the curves of \Cref{ex:reduction:running-table}
(\Cref{thm:reduction:ogg}).
\end{exercise}

\begin{exercise}\label{exer:reduction:defects}
Using \Cref{thm:reduction:kraus}, compute the semistability defect at each additive prime $q\ge5$ of the running examples. Show directly,
following the proof of \Cref{thm:reduction:kraus}, that \curve{726a1} acquires good reduction at $11$ over $\Q_{11}^{\ur}(11^{1/4})$ and not over
$\Q_{11}^{\ur}(\sqrt{11})$.
\end{exercise}

\begin{exercise}\label{exer:reduction:IV2}
Write out the proof of \Cref{cor:reduction:IV-at-2} for \curve{20a1}: identify $E[3]^{I_2}$ with $\Phig_2[3]$, and check that the $3$-division
polynomial of \curve{20a1} has one root in $\Q_2$ and three roots generating a totally ramified cubic extension. Then check the paper's claim
(Remark rem:wild, Proposition prop:forms) that a type $\kod{IV}$ or $\kod{IV}^*$ at $2$ has $P_2(X)=1+2X$: deduce it from
\Cref{cor:reduction:IV-at-2}(ii), (iv) as in the paper link after it, and compare with PARI's \texttt{lfunsympow} for \curve{20a1} and
\curve{52a1}.
\end{exercise}

\begin{exercise}\label{exer:reduction:tate-series}
Show that $a_6(\mathsf q)=-(5s_3+7s_5)/12$ has integer coefficients, and compute $j(E_{\mathsf q})$ up to the term $196884\,\mathsf q$
(\Cref{thm:reduction:tate-curve}). Invert the series to obtain $\mathsf q(J)=J+744J^2+750420J^3+\cdots$ (\Cref{prop:reduction:j-inverse}).
\end{exercise}

\begin{exercise}\label{exer:reduction:t3-539}
Check the paper's claim (Example of Section 3) that $t_3(\curve{539c1},7)=1$: compute $v_7(j)$, the Tate parameter of \curve{539c1} at $7$ to a
few $7$-adic digits and its unit part, and show that they coincide with those of \curve{77c1} (\Cref{ex:reduction:tate-parameters}).
\end{exercise}

\begin{exercise}\label{exer:reduction:11a1-unit}
For \curve{11a1} at $11$, show that the unit part of $q_E$ is $\equiv-1\pmod{11}$ and that it is a $5$th power in $\F_{11}^\times$. This is the
situation of the extra term of Lemma lem:tamfactor at $\ell=5$ (although $\ell=5$ is an Eisenstein prime for \curve{11a1}, so the lemma
applies and Theorem thm:main does not).
\end{exercise}

\begin{exercise}\label{exer:reduction:tamagawa-mult}
Let $E/K$ have multiplicative reduction of type $\kod{I}_n$, with $k$ finite, and let $K'/K$ be the unramified quadratic extension, with
$\Gal(K'/K)=\langle\sigma\rangle$. Using \Cref{thm:reduction:tate-uniformisation}, $E(K)\simeq K^\times/q_E^{\Z}$ in the split case, and
$E(K)\simeq\{w\in K'^\times/q_E^{\Z}:\sigma(w)=w^{-1}\}$ in the nonsplit case (the Frobenius action twisted by $w\mapsto w^{-1}$). Identify
$E_0(K)$ with the classes represented by units, and compute $E(K)/E_0(K)$ in both cases for general $n$: it is $\Z/n$ (split), and
$\Z/2$ or $0$ according as $n$ is even or odd (nonsplit). This reproves \Cref{prop:reduction:tate-tamagawa}(ii). Check $c_2=2$ and $c_3=1$ for
\curve{726a1}, which is nonsplit of type $\kod{I}_6$ at $2$ and $\kod{I}_3$ at $3$.
\end{exercise}

\begin{exercise}\label{exer:reduction:Instar2}
Check the paper's statement (Definition def:terms, Section sec:wild) that the index of a symbol $\kod{I}_m^*$ at $2$ is not the valuation of the
Tate parameter, on \curve{304b1} and \curve{832c1}: compute $m$, $v_2(j)$, $f_2$ and $t_3(E,2)$ (\Cref{thm:reduction:potmult-additive},
\Cref{rem:reduction:Instar-at-2}). Then compute PARI's \texttt{ellinit(E, O(2\^{}20)).tate[3]} for \curve{832c1} and explain the result
(\Cref{warn:reduction:pari-tate}).
\end{exercise}

\begin{exercise}\label{exer:reduction:typeII2}
Check the claim of the original version of the paper (Remark rem:wild) that ``a type $\kod{II}$ at $2$ has $P_2=1$'' against \curve{336f1} and
the type $\kod{II}^*$ curve \curve{1232l1}, using PARI's \texttt{lfunsympow} and the dictionary of \Cref{prop:reduction:inertia-at-2}. Explain
the answer through the inertia images ($\Z/6$, respectively $\Z/2$), and verify that $v_3(\degphi)=e_3(E,2)=2$ for \curve{1232l1}
($\degphi=288$).
\end{exercise}

\section{Notes and references}\label{sec:reduction:notes}

\paragraph{Sources followed.} Sections \ref{sec:reduction:minimal} to \ref{sec:reduction:potential} follow Silverman's first book
\cite[Chapter VII]{SilvermanAEC} (minimal models in VII.1, reduction in VII.2 and VII.5, the formal group in Chapter IV and VII.2, torsion
in VII.3 and VII.7, potential good reduction in VII.5); the proofs of \Cref{prop:reduction:Ens} and \Cref{prop:reduction:filtration} are
the standard ones written out. N\'eron models, Kodaira symbols and Tate's algorithm follow \cite[Chapter IV]{SilvermanATAEC} (Sections
IV.5 to IV.9, with the table of Kodaira types and Tamagawa numbers in IV.9 and Ogg's formula in IV.11), with \cite{BLR} for the general
theory and \cite[Section 3.2]{Cremona97} for the explicit form of the algorithm. The Tate curve follows \cite[Chapter V]{SilvermanATAEC}
(V.3 for the uniformisation, V.5 for curves with $v(j)<0$). The original sources are N\'eron \cite{NeronModels}, Kodaira \cite{Kodaira}, Tate
\cite{Tate75}, Ogg \cite{Ogg}, Saito \cite{Saito88} and Serre and Tate \cite{SerreTate68}; Serre's paper \cite[Section 5.6]{Serre72}
discusses the image of inertia at $2$ and $3$.

\paragraph{Theorem numbers.} We cite theorem numbers only where they have been checked: Dokchitser and
Dokchitser \cite[Theorem 3.1]{DokDok15} and its source Kramer and Tunnell \cite[Proposition 8.20]{KramerTunnell82} (the number was checked in the arXiv
version of \cite{DokDok15} when this book was planned). Elsewhere we cite sections. In particular Kraus's paper \cite{Kraus90}, which gives the semistability defect at
all primes and the full $2$- and $3$-adic classification, is cited as a whole: the numbering of its theorems is quoted inconsistently in
later papers (see Billerey \cite{Billerey09} for the $2$-adic statements and \cite{PinedaMartin26} for a recent restatement). Kraus's
conditions on $c_4$, $c_6$ (\Cref{prop:reduction:min-criterion}(iii)) are in a different paper \cite{Kraus89}. Papadopoulos
\cite{Papadopoulos93} tabulates the Kodaira types at $2$ and $3$ in terms of valuations and congruences. The relation between Kodaira
symbols and wild ramification for twists of curves with multiplicative reduction at $2$ is in Lorenzini \cite{Lorenzini10}; see also
\cite{Praveen26}.

\paragraph{The types $\kod{IV}$, $\kod{IV}^*$ at $2$.} \Cref{cor:reduction:IV-at-2} has two independent proofs: the one given here,
through the component group (\Cref{prop:reduction:component-invariants}), which is the proof of Lemma lem:IV in the revised paper; and
one through the conductor, from \Cref{thm:reduction:tame-additive}: $f_2=2$ forces $\operatorname{sw}=0$, so the image of inertia is a
cyclic group of odd order dividing $24$, hence of order $3$. The original version of the paper left this case as an exception in Proposition
prop:forms.

\paragraph{The exchanges at $3$.} \Cref{prop:reduction:twist-at-3} is proved here from Ogg's formula; we have not found it stated in the
literature in this form, though it can presumably be read off the tables of \cite{Papadopoulos93}.

\paragraph{Computations.} All numbers in the chapter were computed with PARI/GP 2.17.2 \cite{PARI} through \texttt{cypari2}, on Cremona's
tables \cite{ecdata} for conductor below $10{,}000$ ($64{,}687$ curves, $211{,}455$ pairs of a curve and a bad prime). The population checks
quoted in the text are: Kraus's conditions (\Cref{prop:reduction:min-criterion}); $a_p=\bigl(\frac{-c_6}{p}\bigr)$
(\Cref{prop:reduction:split}); the formulas for Step 2 of Tate's algorithm; the table for $p\ge5$; the values of $c$ by type; the twist
exchanges ($90{,}970$ pairs); the types at additive potentially multiplicative primes; the conductor exponents at $2$ and $3$; and the inertia
patterns at $2$ ($20{,}979$ pairs). The Tate parameters were computed from the series of \Cref{prop:reduction:j-inverse}; the discrepancy of
PARI's \texttt{.tate} at $2$ is recorded in \Cref{warn:reduction:pari-tate}.
